\documentclass[twoside]{article}
\usepackage[
    a5paper,
    margin=0.5in,
    includehead,
    headsep=4pt,
    headheight=13pt
]{geometry}
\usepackage{gregosheet}
\usepackage{gregosheet-psalm}
% Font sizes
\newcommand{\musicfontsize}{20}
\newcommand{\lyricfontsize}{10}
\newcommand{\psalmfontsize}{10}

% Fonts
% Change this to any font command you want:
% \rmfamily (roman/serif), \sffamily (sans-serif), \ttfamily (monospace)
% or use \setmainfont{FontName} or \newfontfamily for specific fonts
\newfontfamily\lyricfont{Cambria}
\newfontfamily\psalmfont[BoldFont={Cambria Bold}]{Cambria}
\newfontfamily\headingfont[BoldFont={Cambria Bold}]{Cambria}
\setmainfont{Cambria}

% Spacing
\newcommand{\systemvskip}{10pt}
\newcommand{\lyricvskip}{2pt}
\newcommand{\blockvskip}{10pt}

% Page layout
\usepackage{fancyhdr}
\pagestyle{fancy}
\fancyhf{}
\fancyhead[CO]{\textcolor{red}{\MakeUppercase{\leftmark}}}
\fancyhead[RO]{\thepage}
\fancyhead[LE]{\thepage}
\fancyhead[CE]{\textcolor{red}{\MakeUppercase{\rightmark}}}
\renewcommand{\headrulewidth}{0pt}
\setlength{\parindent}{0pt}

% Title style
\makeatletter
\renewcommand{\maketitle}{
  \begin{titlepage}
    \centering
    \vspace*{\fill}
    {\headingfont\fontsize{32}{38}\selectfont\bfseries\MakeUppercase{\@title}\par}
    \vskip 2em
    {\headingfont\fontsize{16}{20}\selectfont\MakeUppercase{\@subtitle}\par}
    \vspace*{\fill}
  \end{titlepage}
}
\newcommand{\subtitle}[1]{\gdef\@subtitle{#1}}
\newcommand{\@subtitle}{}
\makeatother

% Section style
\usepackage{titlesec}
\titleformat{\section}
  {\centering\headingfont\fontsize{20}{24}\selectfont\bfseries}
  {}{0pt}{\MakeUppercase}
\titlespacing*{\section}{0pt}{\blockvskip}{\blockvskip}

% Subsection style
\titleformat{\subsection}
  {\centering\headingfont\fontsize{16}{20}\selectfont\bfseries\addfontfeature{LetterSpace=10.0}}
  {}{0pt}{\MakeUppercase}
\titlespacing*{\subsection}{0pt}{\blockvskip}{\blockvskip}

% Subsubsection style
\titleformat{\subsubsection}
  {\centering\headingfont\fontsize{12}{14}\selectfont\bfseries\scshape}
  {}{0pt}{}
\titlespacing*{\subsubsection}{0pt}{\blockvskip}{\blockvskip}

% Part title command
\newcommand{\parttitle}[2][]{
  \vskip\blockvskip
  \par\noindent
  \headingfont\fontsize{10}{12}\selectfont
  \textcolor{red}{\MakeUppercase{#2}}%
  \ifx&#1&\else\hfill\textcolor{red}{#1}\fi
  \nopagebreak
  \par
}

% Rubric command
\newcommand{\rubric}[1]{
  \par\noindent
  {\fontsize{8}{10}\selectfont\textcolor{red}{#1}}
  \par
}

% Red text command
\newcommand{\redtext}[1]{
  \par\noindent
  \textcolor{red}{#1}
  \par
}

\newcommand{\p}{\textcolor{red}{P.: }}

\newcommand{\h}{\textcolor{red}{H.: }}

\newcommand{\e}{\textcolor{red}{E.: }}

\newcommand{\lec}{\textcolor{red}{L.: }}

\newcommand{\diac}{\textcolor{red}{D.: }}

% Red dash command
\newcommand{\reddash}{\textcolor{red}{\textemdash{}}\xspace}
\usepackage{xspace}

% Lectio environment
\newenvironment{lectio}[1]{
  \par\noindent
  {\centering\textcolor{red}{\textit{#1}}\par}
  \setlength{\parindent}{0.2in}
  \everypar{\setlength{\parindent}{0.2in}}
}{
  \par
}

% Versicle-Response command
\newcommand{\versicle}[2]{
  \par\noindent
  \hangindent=1em\hangafter=1
  \emergencystretch=1em
  \makebox[1em][l]{\textcolor{red}{V.}}#1
  \par\noindent
  \hangindent=1em\hangafter=1
  \emergencystretch=1em
  \makebox[1em][l]{\textcolor{red}{F.}}#2
  \par
}

% Hymn stanzas environment
\usepackage{multicol}
\newenvironment{hymnstanzas}{
  \setlength{\leftmargin}{0pt}
  \setlength{\rightmargin}{0pt}
  \setlength{\columnsep}{4em}
  \begin{multicols}{2}
  \raggedright
}{
  \end{multicols}
}

\newcommand{\stanzabreak}{%
  \par\addvspace{0.8\baselineskip}%
}


% Single centered stanza (for odd last stanza)
\newcommand{\centerstanza}[2][0.4]{
  \par
  \begin{center}
  \begin{minipage}{#1\textwidth}
  \raggedright
  #2
  \end{minipage}
  \end{center}
  \par
}

% Structured text environment
\newenvironment{preces}[2]{
  \par\noindent
  #1
  \par\noindent
  \hskip0.2in\textit{#2}
  \par
  \setlength{\parindent}{0pt}
  \everypar{\hangindent=0.2in\hangafter=1}
}{
  \par
}

\newcommand{\printsheet}[1]{%
  \ifcsname sheet@#1\endcsname
    \csname sheet@#1\endcsname
  \else
    \textbf{[Missing sheet: #1]}%
  \fi
}

\usepackage{ragged2e}


% ============================================================================
% I. PROPRIUM DE TEMPORE
% ============================================================================

% ----------------------------------------------------------------------------
% I. Tempus Adventus
% ----------------------------------------------------------------------------

% ----------------------------------------------------------------------------
% a. Usque ad diem 16 decembris
% ----------------------------------------------------------------------------

% --- Pro cunctis diebus ---

% ··· Invitatorium ···

\defsheet{RegemVenturum}
  {<-3-3r-3----:-3--3--3--4---3--4t-5---,-5---3--4t---4-4---3-.}
  {  A Ki-rályt, az el-jö-ven-dő U--rat * jöj-je-tek, i-mád-juk!}

% --- Hebdomada 1 Adventus ---

% --- Dominica ---

% ··· I Vesperas ···

% Antiphona 2 (Adv43)
\defsheet[tone={5-c}]{EcceDominusVeniet}
  {<-7-5--7--4u-5--4--3----:-4----4--3---5---7i----7--7-,-5--7--7----8-6--7--tT4-:-3----4--4----4t---5--4t-3--3-?,}
  {  Í-me az Úr el-jö-vend * s_ve-le min-den szent-je-i,  és lé-szen a-ma na-pon   nagy fé-nyes-ség, al-le-lu-ja.}

% Magnificat (Adv40)
\defsheet[tone={1la-a}]{EcceNomenDomini}
  {<X-1-1---3--1--0---34t-5--:-4-6--5---4---4--3t-4-,-4--4--4-2--4---tT4-3-:-2--3---4--eE2-1---1-?,}
  {   Í-me, az Úr-nak ne--ve * a tá-vol-gól kö-ze-leg és az Ő ra-gyo-gá--sa  be-töl-ti a   föl-det.}

% ··· Laudes ···

% Antiphona 1 (Adv42)
\defsheet[tone={8szo-c}]{InIllaDie}
  {<-4-4--4t-4---:-£------------------5u---4-5---rR3-3--,-3r-5-7---4-:-4---4---5---4---eE2--1--:-2--3r-4--4-?,}
  {  A-ma na-pon * megcsordulnak_a_he-gyek é-des-ség-gel, és a dom-bok tej-jel-méz-zel foly-nak, al-le-lu-ja.}

% Antiphona 2 (Adv60)
\defsheet[tone={1la-a}]{MontesEtColles}
  {<X-1--3----4--3t--5---:-5-4--2---3---4----eE2-1--1--,-1--2--3---4---5---4---3r-4-:-5---7---4---5--5u--5--,-5----tT4-2---3--rR3-1--3-1---1--0-:-1-0---3--3---2--3---4--eE2-1--1-?,}
  {   He-gyek és hal-mok * é-ne-kel-nek majd az  Úr-nak, az er-dők-nek min-den fá-i   tap-sol-nak ke-zük-kel, mert el--jön az Úr, az U-ral-ko-dó  ö-rök-or-szá-gá-ba, al-le--lu-ja.}

% Antiphona 3 (Adv44)
\defsheet[tone={4-a}]{EcceVeniet}
  {<-1-2---4t-5--:-4t-6----tT4-5--5-,-2--5----4---Ï3-wW1-wW1-:-0--1--2---2--–3---4--3--2--2-?,}
  {  Í-me, el-jő * a  nagy pró-fé-ta, ki majd meg-új-ít--ja    Je-ru-zsá-le-met, al-le-lu-ja.}

% Benedictus (A) (Adv85)
\defsheet[tone={1re-a}]{DiesDominiSicutFur}
  {<-3--0--1---1tu-5-:-5----5-5u--5----:-[------------tT4¨5uj5T4-:-4--2---3---4-eE2--1--1---,-3----2-1----1-1--1---2---3--4t-4--3--:-2--3---4--eE2-1---0q-1-?,}
  {  Az Úr-nak nap-ja  mint a tol-vaj, * éjszaka_érke-zik,         le-gye-tek a-zért ké-szen, mert a-mely ó-rá-ban nem is vé-li-tek, el-jön az Em--ber-fi-a.}

% ··· II Vesperas ···

% Antiphona 1 (Adv43)
\defsheet[tone={8szo-c}]{Iucundare}
  {<-3--4--5----4uiI7U6¨7iI7i-iI7-:-uU6-5--uU6-tT4-4--,-4--4----4--eE2-1-:-3--3--3---4t--4--3r-rR3--2--3r-4--4-?,}
  {  Vi-ga-dozz im------------már * Si--on le--á---nya, uj-jong re-pes-ve, Je-ru-zsá-lem le-á--nya, al-le-lu-ja.}

% Magnificat (A) (Adv26)
\defsheet[tone={8szo-c}]{SpiritusSanctus}
  {<-1-3-----3--3---4t-4---:-4--4-----5--4---3--4t-5-,-5--556uU6-tT4-4--:-4----4---5--4---3--eE2-1----:-1--3--4---5---4--4-?,}
  {  A Szent-lé-lek Is-ten * le-száll re-ád  Má-ri-a,  ne félj   te--hát, mert kit mé-hed-be fo--gadsz, az Is-ten-nek Fi-a.}

% Magnificat (B) (Adv49)
\defsheet[tone={7do-d}]{NeTimeasMaria}
  {<-8--6-----89p-oO8-8--:-©-----------------6---7---8--tT4-4--,-¦-------------5u-6----:-5----4--5-------7--6--4--4-?,}
  {  Ne félj, Má--ri--a, * mert_kegyelmet_ta-lál-tál az Úr--nál. Íme,_méhedben fo-gansz, s_fi-at szülsz, al-le-lu-ja.}

% Magnificat (C) (Adv64)
\defsheet[tone={1la-a}]{BeataEsMaria}
  {<-1t--5u--5---:-5--5u-5--:-5--5---rR3r-1--,-1----3---2--1--1---2---3---4tT4-3-:-2---3--4--3---2----1--qQ0-1--eE2-1--1-?,}
  {  Bol-dog vagy, Má-ri-a, * mi-vel hit--tél, mert tel-je-se-dik ben-ned mind-az, mit az Úr mon-dott né-ked al-le--lu-ja.}

% --- Feria II ---

% Benedictus (Adv57)
\defsheet[tone={1la-a}]{Leva}
  {<-1tu--5----:-7--5--4---4tT4-:-2-3----4---3---2--1-1--,-1--34t--4-:-5--4--3----4--4---5--,-1-2--3---4--4--4--eE2-:-1----3---eE2-1---1-3--eE2-1----1--?,}
  {  Nézz fel, * Je-ru-zsá-lem,   e-meld föl sze-me-i-det, és lásd meg ki-rá-lyod ha-tal-mát, í-me jön Üd-vö-zí-tőd,  hogy fel-old-jon a bi-lin-csek-ből.}

% Magnificat (Adv51)
\defsheet[tone={1la-a}]{AngelusDomini}
  {<-1--3--3---rR3-1---qQ0-:-3--4----5---4--3--3r-3-,-3---4--3--2---1e-eD0-:-2e-4-----eE2-1---qQ0--1--eE2-1--1-?,}
  {  Az Úr-nak an--gya-la  * kö-szön-töt-te Má-ri-át, s_ő mé-hé-ben fo-gant  a  Szent-lé--lek-től, al-le--lu-ja.}

% --- Feria III ---

% Benedictus (Adv90)
\defsheet[tone={1la-a}]{EgredieturVirga}
  {<-0----1---1tu-5---:-[--------------tT4¨5uj5T4-:-5--4---3--3-3r--3-:-3--2--3---4--eE2-1--1-,-3-----2---1--2--3---4t--4---3--:-3--2--3---4--eE2-1--1-?,}
  {  Vesz-sző sar-jad * Jessze_gyökeré-ből,         be-töl-ti a föl-det az Úr-nak di-cső-sé-ge, s_meg-lát-ja ak-kor min-den test az Úr-nak üd-vös-sé-gét.}

% --- Feria IV ---

% Benedictus (Adv52)
\defsheet[tone={7la-d}]{IerusalemRespice}
  {<-4t-6--7i--8----:-8--9----ü---8--9---8-,-8--8---45z--6--5z-4--4-?,}
  {  Je-ru-zsá-lem, * te-kint Nap-ke-let-re, és lás-sál, al-le-lu-ja.}

% Magnificat (Adv53)
\defsheet[tone={2-f}]{DeSion}
  {<-1--1r-3---:-3---0-1---1---,-3--3--1-0q-1-:-3--0--1eE2-3---1-?,}
  {  Si-on-ból * jön a Tör-vény, az Úr i-gé-je  Je-ru-zsá--lem-ből.}

% --- Feria V ---

% Magnificat (Adv53)
\defsheet[tone={4-a}]{BenedictaTu}
  {<-1--2----4t---5--:-4t-6--5---4----5--5---,-2--5----4-Ï3-2--1---wW1-:-0---1---2r---2--2-?,}
  {  Ál-dott vagy te * az as-szo-nyok kö-zött, ál-dott a te mé-hed-nek   gyü-möl-cse, Jé-zus.}

% --- Feria VI ---

% Benedictus (Adv54)
\defsheet[tone={1la-a}]{EcceVenietDeus}
  {<-1tu-5---:-5u-5---5--4--3---4--1-,-3--2---1---3--3--4tT4-:-3--3---eE2-1--4---eE2-qQ0-:-1--eE2-1--1-?,}
  {  Í---me, * el-jön az Is-ten-Em-ber Dá-vid-nak há-zá-ból,   el-fog-lal-ni tró-nu--sát,  al-le--lu-ja.}

% Magnificat (Adv54)
\defsheet[tone={4-a}]{ExAegyptoVocavit}
  {<-1-2----4t--5---:-4---5---6--tT4-5--5-5--,-5---2----5----4--Ï3-wW1-wW1-:-0--1--–2e--4--3--2--2-?,}
  {  E-gyip-tom-ból * hív-tam az én  fi-a-mat, jön már, hogy üd-vö-zít-se    né-pe-met, al-le-lu-ja.}

% --- Sabbato ---

% Benedictus (Adv55)
\defsheet[tone={4-a}]{SionNoliTimere}
  {<-4t-5---:-4--5z---tT4-5--5---,-2-5--4-Ï3-2--1--wW1-:-0--1---2---–3---4--3--2--2?,}
  {  Si-on, * ne félj már to-vább, í-me a te Is-te-ned   el-jön hoz-zád, al-le-lu-ja.}

% --- Hebdomada 2 Adventus ---

% --- Dominica ---

% ··· Laudes ···

% Antiphona 1 (Adv59)
\defsheet[tone={1la-a}]{UrbsFortitudinis}
  {<X-1t-5-:-[---------7--5---4--5-----:-3----5---4---5--rR3-2e-4--eE2-1--1-,-5uj5-5-5-5i--8iK5-5-,-5---7----5--5---4---5--rR3-4-:-2----4--2----4--wW1-0--:-1e-2--1--1-?,}
  {   Si-on  erősséges vá-ros ne-künk, * mert biz-tos fa-la  az Üd-vö--zí-tő, és   ő a bás-tyá--ja, hát nyis-sá-tok meg ka-pu--it, mert ve-lünk az Is--ten, al-le-lu-ja.}

% Antiphona 2 (Adv44)
\defsheet[tone={7do-a}]{OmnesSitientes}
  {<ô-tg3¨5z-5---:-5--5zZ5-4--4----:-4---4---4--rR3-wW1-2---1--:-5--6---7iI7-zZ5-4--5uU6-5--:-4---2---4--5---4---eE2-1--:-2e-2e-1--1-?,}
  {   Ó,     mind, ti szom-ja-zók, * jer-tek az é---lő  víz-hez, ke-res-sé---tek az U----rat, míg meg-ta-lál-hat-já--tok, al-le-lu-ja.}

% Antiphona 3 (Adv60)
\defsheet[tone={3-c}]{EcceDominusNosterCumVirtute}
  {<-4-5---7--7--:-7--7--6---5---7--tT4-4--,-5----7----6--5-4---5---4-:-4---3--4--5--eE2-1-:-2--3r-2--2-?,}
  {  Í-me, az Úr * ha-ta-lom-mal ér-ke--zik, hogy szol-gá-i-nak sze-mét meg-vi-lá-go-sít-sa, al-le-lu-ja.}

% Benedictus (B) (Adv68)
\defsheet[tone={4-a}]{EcceMitto}
  {<-1-2--4--4t--5---:-4---5--6--tT4-5---5--5--,-2-5--4-Ï3--wW1-wW1-:-0--1--2---4---2-2-?,}
  {  Í-me el-kül-döm * hír-nö-kö-met hoz-zá-tok, a-ki u-tat ké--szít  az én szí-nem e-lőtt.}

% ··· II Vesperas ···

% Antiphona 1 (Adv58)
\defsheet[tone={1la-a}]{EcceInNubibusCaeli}
  {<-1--0--1--3---2--1-qQ0-:-3--4--5---4----tT4-3-:-3r---2--eE2-1---qQ0-:-1--eE2-1--1-?,}
  {  Ím az ég fel-hő-i-ben * ér-ke-zik majd az  Úr  nagy ha-ta--lom-mal,  al-le--lu-ja.}

% Antiphona 2 (Adv59)
\defsheet[tone={1la-a}]{EcceApparebit}
  {<-1t-5-:-[----------7--tT4t-:-3---4--5---eE2-1--,-5--5--4-----5--5u--5--:-5--5----55i-8iK5-5-,-5----7--5----4----5--rR3r-:-3-----3---4--eE2-qQ0-:-1--eE2-1--1-?,}
  {  Í--me, megjelenik az Úr,  * nem ál-tat té--ged, ha ha-laszt-ja jöt-tét, te csak vár-jad  őt, mert bi-zony-nyal el-jön,   s_nem kés-le-ke--dik,  al-le--lu-ja.}

% Magnificat (AB) (Adv53)
\defsheet[tone={8do-c}]{Veniet}
  {<-7---7---7-7iI7-7----:-¦-------6----5u-4--,-3--5---7---7--7----7i--7-:-7---7--6--5---55--4---5zZ5-4--3r-4-?,}
  {  Jön már u-tá---nam, * aki_erő-sebb ná-lam, ki-nek nem va-gyok mél-tó  még sa-ru-szí-ját sem meg--ol-da-ni.}

% --- Feria II ---

% Benedictus (Adv82)
\defsheet[tone={8szo-c}]{DicitDominus}
  {<-4---3----1--3ed134tT4-:-0--1--2---3r-4--5----4tT43r-rR3-,-5----7--4--5---4----5--4-:-4-3---4----5--4---4rR3E1-:-2--3r-4--4-?,}
  {  Így szól az Úr:       * Pe-ni-ten-ci-át tart-sa-----tok,  mert el-kö-zel-gett im-már a men-nyek or-szá-ga,      al-le-lu-ja.}

% --- Feria III ---

% Magnificat (Adv15)
\defsheet[tone={5-c}]{VoxClamantisInDeserto}
  {<-¢-----------4---3-:-5-7----7ui-7----:--8--8----8--7i--5-:-5--5--5---4t-3--,-3-4---4---5--3--3r--4-:-7--5--4----5--4--3---3-?,}
  {  A_Kiáltónak sza-va  a pusz-tá--ban: * „Ké-szít-sé-tek el  az Úr-nak út-ját, e-gye-nes-sé te-gyé-tek Is-te-nünk ös-vé-nye-it!”}

% --- Feria IV ---

% Benedictus (Adv61)
\defsheet[tone={8szo-c}]{SuperSoliumDavid}
  {<-4--4---1eE2-3r--:-4tT4-3--4tT4¨5uj5-,-4--4--5-7--4---5--4-:-4-4---eE2-1---3----4t-4---4rR3d1-:-2--3r-4--4-?,}
  {  Dá-vid trón-ján * ül   az Úr,         és az ő or-szá-gá-ban u-ral-ko--dik mind-ö--rök-ké,      al-le-lu-ja.}

% Magnificat (Adv68)
\defsheet[tone={4-a}]{SionRenovaberis}
  {<-4t-5---:-4--5----6----tT4-5-5----,-2---5--4---Ï3-2-1--wW1-:-0-1--2--4---2---2-?,}
  {  Si-on, * bi-zony majd meg-ú-julsz, meg-lá-tod az i-ga-zat,  a-ki el-jön hoz-zád.}

% --- Feria V ---

% Benedictus (Adv84)
\defsheet[tone={5-c}]{PonentDominoGloriam}
  {<-¦-----------uU6-5-:-5-5--4---tT4-3---,-4--3--5-77i-7-:-8-ö--8--¦-----------7i--5--,-5--8---7-6u-tT4-5-:-4--4t--4t--3--3-?,}
  {  Megadják_az Úr--nak a di-csé-re--tet * és az ő hí--rét a tá-vo-li_szigetek zen-gik, mi-vel í-me el--jő  és már nem ké-sik.}

% Magnificat (Adv69)
\defsheet[tone={4-a}]{QuiPostMeVenit}
  {<-1-2--4-4--5z--5----:-4-5----zh4--5--5--,-2---5--4----Ï3-wW1-wW1-:-0---1--2--3--4--–3-2---2-?,}
  {  A-ki u-tá-nam jön, * e-lőbb volt ná-lam, nem va-gyok én mél-tó    meg-ol-da-ni sa-ru-szí-ját.}

% --- Feria VI ---

% Benedictus (Adv69)
\defsheet[tone={1re-a}]{DicitePusillanimes}
  {<-3----4t--5---:-7--5-----4--2e-4-:-4-3--2---1ed1-0q--1--,-4-4---2--4tT4-3--:-3-4--2-eE2--1--1-?,}
  {  Mond-já-tok: * ki-csiny-hi-tű-ek, e-rő-söd-je---tek meg, í-me, az Is---ten, a mi U-runk el-jő.}

% --- Sabbato ---

% Benedictus (Adv70)
\defsheet[tone={1la-a}]{Levabit}
  {<-1----0--1---2e-4-----eE2--1--qQ0-:-3-5---4--3---4-3---,-3--3---4---tT4--3-:-3r---2-3----2-----1--0q-1-?,}
  {  Nagy je-let tá-maszt majd az Úr  * a nem-ze-tek e-lőtt, és ösz-sze-gyűj-ti  mind a szét-szórt Iz-rá-elt.}

% --- Hebdomada 3 Adventus ---

% ··· I Vesperas ···

% Antiphona 1 (Adv74)
\defsheet[tone={8szo-c}]{IerusalemGaude}
  {<-1--2--3---4----eE2-1-----:-3----3-4t--4--,-5u-7-:-5u-5---4----5---4-:-3---4--eE2-1--:-3--4t-4--4-?,}
  {  Je-ru-zsá-lem, ör--vendj * nagy ö-röm-mel, í--me, el-jön majd hoz-zád sza-ba-dí--tód, al-le-lu-ja.}

% Magnificat (Adv72)
\defsheet[tone={1la-a}]{AnteMe}
  {<X-1---3----1--0---3-4t--5----:-6-----5----4---4-3t-4--,-4--4---4--2---4---5---4---3-:---2----3r--eE2--1---0q--1-?,}
  {   Nem volt Is-ten e-lőt-tem, * s_nem lesz más u-tá-nam, ne-kem ha-jol meg min-den térd, s_en-gem vall min-den nyelv.}

% ··· Laudes ···

% Antiphona 1 (Adv73)
\defsheet[tone={1re-a}]{VenietDominus}
  {<-5--4---4--5--4----:-3---4--5-----eE2-1---1--:-1---1--2--3--4t--4-:-4-5--4---3---4t--3--3--,-3--4---5---4---5---4--2--3---4---2---3---qQ0-:-1--eE2-1--1-?,}
  {  El-jön az Is-ten, * nem ha-laszt-ja  jöt-tét, meg-vi-lá-go-sít-ja  a sö-tét-ség rej-te-két, ki-nyi-lat-koz-tat-ja ma-gát min-den nép-nek,  al-le--lu-ja.}

% Antiphona 2 (Adv75)
\defsheet[tone={5-c}]{MontesEtOmnesColles}
  {<-¢-----------4---3---:-3---3-5---7i--7--:-7-7---7--7i-8-:-7--8-ö---9--iI7-7--:-8---7i-5-:-5--5--5--4t-3---,-3----3r-4t---3r-4---:-7--7---4--tT4-3-?,}
  {  A_hegyek_és hal-mok * meg-a-láz-kod-nak, a gör-be-sé-gek ki-e-gye-ne-sed-nek, s_a rö-gök út-tá si-mul-nak, jöjj el hát, U--runk, ne kés-le-ked-jél!}

% Antiphona 3 (Adv74)
\defsheet[tone={1la-a}]{DaboInSionSalutem}
  {<X-1---3-1---0--3--4---3--4t-5----:-6----5--4---4---4---3--4---5--eE2-qQ0-:-1--eE2-1--1-?,}
  {   Meg-a-dom az Üd-vöt Si-on-nak, * s_Je-ru-zsá-lem-nek di-cső-sé-ge--met,  al-le--lu-ja.}

% ··· II Vesperas ···

% Antiphona 3 (Adv74)
\defsheet[tone={1la-a}]{IusteEtPieVivamus}
  {<-1-0--1---2e----4--eE2-1--1---,-0---1--3---1-:-1-0---1---4--3----:-0q-3--3---0e--1-?,}
  {  I-ga-zul s_jám-bo-ran él-jünk, vár-va vár-juk a bol-dog re-ményt, az Úr-nak jöt-tét.}

% Magnificat (C) (Adv68)
\defsheet[tone={8do-c}]{VenietFortior}
  {<-7---7---¨-7-7iI7-7--:-¦-------6----5u-4--,-3--5---7---7--7----7i--7-:-7---7--6--5---55--4---5zZ5-4--3r-4-?,}
  {  Jön már * u-tá---nam, aki_erő-sebb ná-lam, ki-nek nem va-gyok mél-tó  még sa-ru-szí-ját sem meg--ol-da-ni.}

% --- Feria II ---

% Benedictus (Adv66)
\defsheet[tone={8szo-c}]{DeCaeloVeniet}
  {<-4-7-----7---8--9--8----:-9-7---8--7--4--4--4---,-7--7--8-7--5u-7-:-7--7iI7-tT4-3--4--4-?,}
  {  A menny-ből el-jö-vend * U-ral-ko-dó Is-te-nünk, és az ő ke-zé-ben or-szág és  ha-ta-lom.}

% Magnificat (Adv82)
\defsheet[tone={8szo-c}]{BeatamMeDicent}
  {<-5u--5---4---5---4--4---:-5--4---5---rR3-4---4t-3--3--,-4----4-4--3--5---7----7--8--7iI7-:-5---5---7--7--4--4-?,}
  {  Bol-dog-nak mon-da-nak * en-gem min-den nem-ze-dé-kek, mert a-lá-za-tos szol-gá-ló-ját    meg-lát-ta az Is-ten.}

% --- Feria III ---

% Benedictus (Adv83)
\defsheet[tone={4-a}]{ElevareElevare}
  {<-1----2----4t---5----:-4-5---6-----5--4--5---5--,-2----5--4---Ï3-wW1-wW1-:-0--1--2---4--2-2-?,}
  {  Kelj föl, kelj föl, * e-mel-kedj, Je-ru-zsá-lem, vesd le már i--gá--dat,  Si-on rab le-á-nya.}

% Magnificat (Adv52)
\defsheet[tone={1la-a}]{AntequamConvenirent}
  {<X-1--1-3----1---0--34t-5----:-6---5--£------------3t-4-,-£--------------2---4---5--4--2-:-2e-4-----eE2-1---qQ0--1--eE2-1--1-?,}
  {   Mi-e-lőtt egy-be-kel-tek, * úgy ta-láltatott_Má-ri-a,  hogy_gyermeket hor-doz mé-hé-ben a  Szent-lé--lek-től, al-le--lu-ja.}

% --- Feria IV ---

% Benedictus (Adv85)
\defsheet[tone={2-f}]{Consolamini}
  {<-3--3----3--eE2-1----:-1--0----1--3rR3-3----1--0----1----3rR3-3--:-eE2-qQ0--1--eE2-1--1-?,}
  {  Vi-gasz-ta-lód-jál, * vi-gasz-ta-lód--jál, vá-lasz-tott né---pem, így szól az Úr--is-ten.}

% Magnificat (Adv68)
\defsheet[tone={7do-d}]{TuEsQuiVenturusEs}
  {<-8--8----6--8o-8--7iI7U6-4---:-5--5-----7-7--:-5---6---5--4--4---,-£-------------3-----4t-eE2-1--:-3--4t-4--4-?,}
  {  Te vagy az el-jö-ven----dő, * te vagy, U-ram, kit vár-va vá-runk, hogy_megszaba-dítsd né-pe--det, al-le-lu-ja.}

% --- Feria IV ---

% Benedictus (Adv82)
\defsheet[tone={2-f}]{ConsurgeConsurge}
  {<-1r-3---:-eD0q-1-----:-3----1---3-3---1--0q-1--,-3-0----1--3---2---1---1-?,}
  {  In-dulj, in---dulj, * vedd fel e-rős-sé-ge-det, U-runk ha-tal-mas kar-ja!}

% --- Feria IV ---

% Magnificat (Adv84)
\defsheet[tone={1la-a}]{HocEstTestimonium}
  {<X-1--qQ0-3--4--3---4t-5----:-6--5----4--4---3t--4---,-4--4-2--4---tT4-3---:-2-3r---eE2--1--1-?,}
  {   Ez a   ta-nú-ság-té-tel, * me-lyet Já-nos mon-dott: ki u-tá-nam jön majd, e-lőbb volt ná-lam.}

% ----------------------------------------------------------------------------
% b. Post diem 16 decembris
% ----------------------------------------------------------------------------

% --- Hebdomada 4 Adventus ---

% ··· I Vesperas ···

% Antiphona 1 (Adv119)
\defsheet[tone={1la-a}]{EcceVenietDesideratus}
  {<X-1-1--3-1----0--34t-5------:-6---5---4--4---4--3t-4--,-£----------2----4---tT4-3-:-3--3--4---eE2-qQ0-:-1--eE2-1--1-?,}
  {   Í-me a vágy-va vá--gyott, * min-den-né-pek re-mé-nye, betelik_fé-nyes-ség-gé--vel az Úr-nak há--za,   al-le--lu-ja.}

% Antiphona 2 (NZS26)
\defsheet[tone={4-a}]{VeniDomine}
  {<-1----2---4-4t-5-----:-4--5---6--5---4--5x--5x-,-5--2----5---4--Ï3---2--1--w\\W1\\-:-0--1---–2e-4--3--2y-2y-?,}
  {  Jöjj el, ó U--runk, * és már ne kés-le-ked-jél, ol-dozd föl go-nosz te-te-it        né-ped-nek Iz-ra-el-nek.}

% Antiphona 3 (NZS232)
\defsheet[tone={4-a}]{EcceIamVenit}
  {<-1-2--4---4t--5---:-5--4-5---6---tT4-5--5-,-2----5---4--Ï3--2--1---w\\W1\\-:-0-----1--2--4-2----2-?,}
  {  Í-me már itt van * az i-dők tel-jes-sé-ge, mely-ben Is-ten el-kül-dé        szent Fi-át a föld-re.}

% ··· Laudes ···

% Antiphona 1 (Adv118)
\defsheet[tone={1la-a}]{CaniteTuba}
  {<-3----0--1----1tu-5----:-5--5--tT4¨5uj5T4-:-4----5--4---33--3--23r-eE2-1---1-,-3-2---1---1---1---2--3--4tT4-3--:-3--qQ0-1--4---2--eE2-1--1-?,}
  {  Szól-ja-tok, kür-tök, * Si-on-ban,         mert kö-zel van az Úr--nak nap-ja. Í-me, jön meg-men-te-ni min--ket, al-le--lu-ja, al-le--lu-ja.}

% Antiphona 2 (Adv120)
\defsheet[tone={1la-a}]{DominusVeniet}
  {<-qQ0-3---4t-5---:-4----7---5--4---7-55-1---3----1--,-5uj55-4----7--7--7uU6Z5-5--:-5--4---3--4t--4-:-5---4--3---4--1--,-4-rF1-3--1----3-3---2e-qQ0-:-0--1--eE2-3r-4--:-0--1--3r-2e--4--eE2-1--1-?,}
  {  Jön már az Úr, * most fus-sa-tok e-lé ezt mond-ván: ó,    nagy Fe-je-de-----lem, ki-nek or-szá-ga  nem ér-het vé-get! E-rős Is-ten, U-ral-ko-dó,   bé-ke Ki--rá-lya, al-le-lu-ja, al-le--lu-ja.}

% Antiphona 3 (Adv121)
\defsheet[tone={2-f}]{Omnipotens}
  {<-3---3---3--1--¨-1-0----1e-3--3--0---,-2-3--4--3--1---3---3---2e--4-eE2-qQ0--:-1--eE2-1--1-?,}
  {  Min-den-ha-tó * I-géd, ó  Is-te-nünk, a te ki-rá-lyi szé-ked-ből a-lá--száll, al-le--lu-ja.}

% Benedictus (B) (Adv68)
\defsheet[tone={7la-d}]{QuomodoFietIstud}
  {<-4--7---7---8--7---5u--:-¦-----------4---5uj5-:-7--7---8o--8oO8-:-7----7--5--4--5---7--:-7---4--3r-4--,--8o---8---9õ-õÔö-oO8-9--8---9---8oO8-:-8-9-----8--7---7i-7---8--7-----5u-7-,-6--7-8--7---7---uU6-5--uU6-tT4-4-:-5---5--5---zZ5-4--4-?,}
  {  Mi-kép-pen le-het ez, * Istennek_an-gya-la,    mi-vel fér-fit,   hogy tő-le fo-gan-jak, nem is-me-rek? „Hall-jad Má-ri--a,  Is-ten szü-ze,    a Szent-lé-lek Is-ten le-száll re-ád, és a Ma-gas-ság-be--li e---re--je  meg-ár-nyé-koz té-ged.”}

% Benedictus (C) (Adv85)
\defsheet[tone={4-a}]{ExQuoFactaEst}
  {<-1--2-----4---4t--5---:-4--5----6--tT4-5---5--,-2t--ÏrR3-wW1-wW1-:-0--1---2---4-2----2-?,}
  {  Mi-helyt hal-lot-tam * kö-szön-té-sed sza-vát, fel-uj---jon-gott  mé-hem-ben a gyer-mek.}

% ··· II Vesperas ···

% Antiphona 2 (Adv120)
\defsheet[tone={1re-a}]{EruntPravaInDirecta}
  {<-3-0---1--1t-5u-:-[---------tT4--5----:-3r-5--4---2---3---4--3--2--1---1--,-3----4---3--3-0--:-2e-4--3---3--eE2-qQ0-:-1--eE2-1--1-?,}
  {  A gör-be-sé-gek  egyenessé lesz-nek, * a  rö-gös föl-dek út-tá si-mul-nak, jöjj el, Ó, U-runk és ne kés-le-ked-jél,  al-le--lu-ja.}

% --- Antiphonae ad Laudes et Vesperas a die 17 ad diem 23 decembris ---

% ··· Feria II ···

% Antiphona 1 (NZS231)
\defsheet[tone={4-a}]{EcceVenietDominus}
  {<-1-2--4--4t-5---:-4---5---6--5--4----5-5-,-2---5----4-Ï3-2--1----w\\W1\\-:-0-1-2--4--2---2-?,}
  {  Í-me jő az Úr, * min-den ki-rá-lyok U-ra, bol-dog, a-ki ké-szen áll      ő-e-lé-be fut-ni.}

% Antiphona 2 (Adv69)
\defsheet[tone={7la-d}]{CantateDominoCanticum}
  {<-4--7--7---8----9--8----:-8--8-9--7---5--7--7--,-7--7---5--7--iI7--7-:-7-7iI7-5---5---4--4-?,}
  {  Az Úr-nak zeng-je-tek, * új é-ne-ket az Úr-nak, di-csé-re-te szól-jon a föld min-den vé-gén.}

% ··· Feria III ···

% Antiphona 1 (NZS233)
\defsheet[tone={4-a}]{EgredieturDominus}
  {<-1--2---4t-5--:-4--5-6-----tT4-5---5--,-2---5----4---Ï3-wW1-w\\W1\\-:-0---1---2--4--2--2-?,}
  {  Ki-lép az Úr * az ő szent he--lyé-ről, jön mint sza-ba-dí--tó,       meg-men-te-ni né-pét.}

% --- Antiphonae ipsis diebus decembris a 17 ad 24 assignatae ---

% ··· Die 17 decembris ···

% Benedictus (Kar16)
\defsheet[tone={8szo-c}]{Scitote}
  {<-5---3--3r--4--:-4----5--4---3---4--4---5--3y--3y-,-3--5----7i--6---7---5--:-4---5---4---3--eE2-1--:-3--4t-4x-4x-?,}
  {  Tud-já-tok meg, hogy kö-zel van Is-ten or-szá-ga;  bi-zony mon-dom nék-tek: már nem kés-le-ke--dik, al-le-lu-ja.}

% Magnificat
\defsheet[tone={2-f}]{OSapientia}
  {<-1e3+-ed1-0q---1----:-Ÿ----------------------3---1e---1--0í-,-Ÿ-----------0q-1-:-Ÿ-------------3--1--1---3-3--1---3--3--4tT4¨5uj5-,-5-----5--4---3--3--2---1e-0í-,-0e---3-:-3----3----1---0q-1qQ0Ôð-:-0q-3-3---0---0e-1í-?,}
  {  Ó    Böl-cses-ség, * ki_a_Fölségesnek_szájá-ból szár-ma-zol, ki_a_minden-sé-get egyik_végétől má-si-kig e-rő-sen át-fo-god,        s_sze-lí-den el is ren-de-zed: Jöjj el, s_mu-tasd meg né-künk     az o-kos-ság út-ját!}

% ··· Die 18 decembris ···

% Benedictus (Adv98)
\defsheet[tone={1la-a}]{VigilateAnimo}
  {<-1---1----0---1---3---2--1---qQ0--,-3--4---tT4-3-:-3r-2---eE2-1--1-?,}
  {  Vir-rasz-sza-tok lel-ke-tek-ben, * kö-zel van már az Úr, az  Is-ten.}

% Magnificat
\defsheet[tone={2-f}]{OAdonai}
  {<-1e3+-ed1-0q-1----:-Ÿ-----------3---1--1--0í-,-0--Ÿ------------------3---1----1---3---1--3---4tT4¨5uj55\%-,-5----4--¢-----------3---2-----1e-0í-,-0e---3y--1--3-----1---0q--1qQ0Ôð-:-0--1----3----0e--1í-?,}
  {  Ó    Á---do-náj, * Izrael_házá-nak ve-zé-re,  ki Mózesnek_a_vöröslő tűz-láng-ban meg-je-len-tél,           s_ne-ki Sion_hegyén tör-vényt ad-tál: Jöjj el, és válts meg min-ket,     ki-nyúj-tott kar-ral!}

% ··· Die 19 decembris ···

% Magnificat
\defsheet[tone={2-f}]{ORadixIesse}
  {<-1e3+-3----1---ed1-0q-1---:-Ÿ------------------3--1e--1--0í--,-0---1---1--:-Ÿ----------3---1--3---4tT4¨5uj55\%-,-5--5--5-5--4---3--eE2-1e--0í-,-0e---33+-1--3---3--3----1---0q--1qQ0Ôð`-:-0q--3--3---0--0e--1í-?,}
  {  Ó    Jesz-sze gyö-ke-re, * ki_jelként_állsz_a né-pek kö-zött, kit lát-ván, a_királyok meg-né-mul-nak,           ki-re a né-pek vá-gya-koz-nak, Jöjj el, és sza-ba-díts meg min-ket,      már ne kés-le-ked-jél!}

% ··· Die 20 decembris ···

% Magnificat
\defsheet[tone={2-f}]{OClavisDavid}
  {<-1e3+-ed1-ed1-0q--1í---:-Ÿ-----------3---1--1--0í-,-1--0q--1-----"-1--3------1-1--3--3---4tT4¨5uj55\%-,-5--5------4--3------3--2---1e---0í-,-0e---33+-1--3----3--1-ed1-0q--1-"-Ÿ---------3--1---1--0í-,-1---3--3---1---0q-1qQ0Ôð`-:-0-0--1---3--0---0e-1í-?,}
  {  Ó    Dá--vid kul-csa, * Izrael_házá-nak jo-ga-ra   ki meg-nyitsz, és nincs, a-ki be-zár-ja,            be-zársz, és nincs, ki meg nyit-ja:  Jöjj el, és vidd ki a bör-tön-ből a_bilincs-be ver-te-ket, kik sö-tét-ben ül-nek,      a ha-lál ár-nyé-ká-ban.}

% ··· Die 21 decembris ···

% Magnificat
\defsheet[tone={2-f}]{OOriens}
  {<-1e3+-ed1ed1-0q-1í---:-1-1---3----1--1--0í-,-1-0---1---3---4tT4¨5uj55\%-,-5----5---5----4--3--3----2---1e--0í-,-1---3--3---1---0q-1qQ0Ôð`-:-0-0--1---3--0---0e-1í-?,}
  {  Ó    Nap----ke-let, * ö-rök fény su-ga-ra,  i-gaz-ság nap-ja,            Jöjj el, s_vi-lá-go-síts meg min-ket, kik sö-tét-ben ü-lünk,      a ha-lál ár-nyé-ká-ban.}

% ··· Die 22 decembris ···

% Magnificat
\defsheet[tone={2-f}]{ORexGentium}
  {<-1e3+-3--1---ed1-0q-1----:-Ÿ---------------1---3--1ed1Q0-0í-,-1--3----3---4tT4¨5uj55\%-:-5---5--4---3---3--2---1e-0í-,-0e---33+-1--3------3---1--0q-1qQ0Ôð-:-0---1-3---3---0--0e--1í-?,}
  {  Ó    né-pek Ki--rá-lya, * s_mindnyájuk_vá-gya-ko-zá-----sa,  te Szeg-let-kő,            két né-pet egy-be-fog-la-ló,  Jöjj el, és mentsd meg az em-bert,    kit a por-ból al-kot-tál!}

% ··· Die 23 decembris ···

% Magnificat
\defsheet[tone={2-f}]{OEmmanuel}
  {<-1e3+-ed1-ed1-0q-1í--:-1---1e-1-"---1--1---1----3--1ed1Q0í-,-3--3---1---3--3----4tT4¨5u-5x-:-5-----5--4---3--3--2--1e-0í-,-0e---3-!-1--3--3--1----0q--1qQ0Ôð`-:-0q-3-3y---0--0e-1í-?,}
  {  Ó    Em--má--nu-el! * Tör-vé-nyünk és Tör-vény-ho-zónk!     Né-pek-nek re-mény-sé------ge,  s_min-de-nek Üd-vö-zí-tő-je!  Jöjj el, és üd-vö-zíts min-ket,      mi U-runk Is-te-nünk!}

% --- Die 24 decembris ---

% ··· Laudes ···

% Antiphona 1 (Adv93)
\defsheet[tone={3-c}]{TuBetlehem}
  {<-4--5u--7--7----:-uU6-5u-7----4-:-5u--5----rR3-2---3r-4---,-5----¦-----------5-7--4--:-4--5---5--5----4--5--4--3----4--rR3-2-?,}
  {  Te Bet-le-hem, * Jú--da föld-je, nem vagy a   leg-ki-sebb, mert belőled_jön a Ve-zér, ki kor-má-nyoz-za né-pe-met, Iz-ra--elt.}

% Antiphona 3 (Adv131)
\defsheet[tone={8do-c}]{Crastina}
  {<-¦---------------¨-¦-----------5---7--6--5---4t-:-3---5----¦--------------5---6--5--4-?,}
  {  A_holnapi_napon * eljön_hozzá-tok az Üd-vös-ség, így szól a_seregeknek_U-ra, Is-te-ne.}

% Benedictus (Kar24)
\defsheet[tone={4-a}]{CompletiSunt}
  {<-4--4t--5---:-4--5--6-tT4-5---5--5-,-2t---ÏrR3-wW1--wW1-:-0--1--2---4----2--2y-?,}
  {  Be-tel-tek * Má-ri-á-nak nap-ja-i,  hogy meg--szül-je    el-ső szü-lött Fi-át.}

% ··· Tertia ···

% Antiphona (Adv129)
\defsheet[tone={8szo-c}]{IudaeaEtIerusalem}
  {<-4--4--rR3d1-3--4t-5--5---4----:-4--5u---4---4tT4R3-3-,-uU6-5--5i--6-7--tT4-4--:-4--rR3-1-:-1--2--3---4t---5--4rR3d1-:-2--3r-4--4-?,}
  {  Jú-de-a     és Je-ru-zsá-lem, * ne félj már töb----bé, jöj-je-tek e-lő hol-nap, és í---me, ve-le-tek lesz az Úr,      al-le-lu-ja.}

% ··· Sexta ···

% Antiphona (Adv129)
\defsheet[tone={8do-c}]{HodieScietis}
  {<-uU6-5--7---7---zZ5-4----:-4----5--4--34t-5-,-7--uU6-5---6---7---5--4-:-4--4-3--4t--4--4-?,}
  {  Még ma meg-tud-já--tok, * hogy el-jő az  Úr, és re--gel meg-lát-já-tok az ő di-cső-sé-gét.}

% ··· Nona ···

% Antiphona (Adv129)
\defsheet[tone={1la-a}]{CrastinaDie}
  {<X-1--1--3---1--0---34t-5---:-4-6----5---4--4----3t-4-,-4--4-4---2--4---tT4-3--:-3-2--3---4--3--2--1--1-?,}
  {   El-tö-röl-te-tik hol-nap * a föld-nek go-nosz-sá-ga, és u-ral-ko-dik raj-tunk a vi-lág Üd-vö-zí-tő-je.}

% ----------------------------------------------------------------------------
% II. Tempus Nativitatis
% ----------------------------------------------------------------------------

% ----------------------------------------------------------------------------
% a. Usque ad USQUE AD SOLLEMNITATEM EPIPHANIAE
% ----------------------------------------------------------------------------

% --- Pro Cunctis Diebus ---

% ··· Invitatorium ··· (NZS137)

\defsheet{ChristusNatusEstNobis}
  {<-33r---3-:-3--3---3---4--3----4t-5-----:-5---3--4t---4-4---3-.}
  {  Krisz-tus ma meg-szü-le-tett né-künk: * jöj-je-tek, i-mád-juk.}

% --- Die 25 decembris - In Nativitate Domini ---

% ··· Officium lectionis ···

% Antiphona 1 (NZS138)
\defsheet[tone={2-g}]{DominusDixit}
  {<ô-2--2y-¨-2---2--2r-2--:-2--4--44-2r---1í-,-3r-5--5----5--rR3-2r-2y-.}
  {   Az Úr * mon-dá né-kem: én Fi-am vagy te,  én ma szül-te-lek té-ged.}

% Antiphona 2 (Gondolkodó 4.)
\defsheet[tone={8szo-c}]{TamquamSponsus}
  {<-4----5--4--rR3-:-3--5--uU6--5u-zZ5-:-4--4-5----5--4--4-?,}
  {  Mint Vő-le-gény, ki-lé-pett az Úr    az ő nász-há-zá-ból.}

% Antiphona 3 (Gondolkodó 4.)
\defsheet[tone={1re-a}]{DiffusaEst}
  {<-3---4t--55--¨-2--3---4---eE2-1--1--,-3---4t-44---3--eE2-qQ0-1----eE2-1---1-?,}
  {  Ked-ves-ség * öm-lött el aj--ka-don: meg-ál-dott az Is--ten mind-ö---rök-ké.}

% ··· Laudes ···

% Antiphona 1 (Kar37)
\defsheet[tone={2-g}]{QuemVidistis}
  {<ô-4---4---4--2----2----4--4----:-1----2--4---3--2r---4--,-2--1--2---4--2---4--2---:-2--2--1--2-4---4---1---4--2---,-1w--4---4--4--2---4---2--:-2----1---2---4--3--2---4---4--,-4---5--3-4---2--1----2--:-1--2--4--4---1--1r-2y-2y-?,}
  {   Kit lát-ta-tok, pász-to-rok? * Mond-já-tok el gyor-san! Ad-já-tok hí-rül ne-künk, ki az ki a föl-dön meg-je-lent? Lát-tuk az Új-szü-löt-tet, s_an-gya-lok kó-ru-sát lát-tuk, kik az U-rat di-csér-ték, al-le-lu-ja, al-le-lu-ja.}

% Antiphona 2 (Kar40)
\defsheet[tone={7la-a}]{Angelus}
  {<ô-1--1t-5----:-[----------------4--56uU6Z5-5\\u\\-:-4-5---7----7-6--5---7---6--5---6---5--,-5--3---5---6---5--4----5-4---rR3-2e-:-1--3--4--wW1-0---1w-2--1í-1í-?,}
  {   Az an-gyal * így_szólt_a_pász-to-rok-----hoz:     í-me, nagy ö-rö-met hir-de-tek nek-tek, mi-vel meg-szü-le-tett e nap né--künk az Üd-vö-zí--tő, al-le-lu-ja.}

% Antiphona 3 (Kar45)
\defsheet[tone={8do-c}]{ParvulusFilius}
  {<-7--6-----5--7--4x-:-3---5--7----7-6--tT4-5--4x--,-7----7--7--8o---8v--:--9-8---7--8---uU6-5x--:-7--8--uU6-5-:-6--5z-4x-4x-?,}
  {  Ki-csiny-ke Fi-ú  * szü-le-tett e na-pon ne-künk, s_ne-ve ez lesz majd: „E-rős-sé-ges Is--ten”, al-le-lu--ja, al-le-lu-ja.}

% Benedictus
\defsheet[tone={8szo-c}]{Gloria}
  {<-3--4t--4---:-4-4--4---5---4---3--4t--5-,-5--6-7---6---5--4---4-:-4-4--4-4--5--4--tT4-3--eE2-1--:-3--4t-4x-4x-?,}
  {  Di-cső-ség * a ma-gas-ság-ban Is-ten-nek és a föl-dön bé-kes-ség a jó-a-ka-ra-tű em--be-rek-nek, al-le-lu-ja.}

% ··· II. Vesperas ···

% Antiphona 1 (Kar80)
\defsheet[tone={1re-a}]{TecumPrincipium}
  {<-3--0--1--1t-5u-5---:-[---------------tT45uj5T4-:-4-3-----4---55-rR3r-1--,-3--2---1---1----2--3---4tT4-3-:-3-3---2--3--4----eE2-1-1-?,}
  {  Ti-éd az el-ső-ség * hatalmadnak_nap-ján,        a szent-ség fé-nyé--ben, ma-gam-ból szül-te-lek Té---ged a haj-na-li csil-lag e-lőtt.}

% Antiphona 2 (Kar84)
\defsheet[tone={8szo-c}]{ApudDominum}
  {<-4--7\\-uJ4-:-6u-5--4---tT4-3y-,-3t-7--7--uJ4-6-7--5--5-6---5--4x-4x-?,}
  {  Az Úr--nál * az ir-gal-mas-ság, és bő-sé-ges ő-ná-la a sza-ba-dí-tás.}

% Antiphona 3 (Kar65)
\defsheet[tone={1la-a}]{InPrincipio}
  {<-1---1t--5----:-[-----------5u-5--:-5--5---4----3--4-1-,-3-2---1--1---1---2--3----4tT4-3---:-3-2--3---4--3--2--1í-1í-?,}
  {  Kez-det-ben, * minden_idők e--lőtt Is-ten volt az I-ge, Ő az, ki meg-szü-le-tett né---künk, a vi-lág Üd-vö-zí-tő-je.}

% Magnificat
\defsheet[tone={6-a}]{Hodie}
  {<X-3-4t-5---:-55T45z-5---4---3t-5x--,-3-3r-4-:-4--4--4--3--4t-zZ5-rR3-3y--,-3-4z-6-:-¥------------------4--zZ5T4-3t--5-:--3--4---6---6c--4--6---tT4-3t--5x-,-5-7--™7iI7uU6Z5-:-5--5---4--5---6--5-4--3t--5x--3r---4x-,-3--eE2-1í--1-2--3---4t--4x--3--eE2-q\\Q0\\-:-1--3--3y-3y-?,}
  {   E na-pon * Krisz--tus szü-le-tett, e na-pon az Üd-vö-zí-tő meg-je--lent, e na-pon a_földön_énekelnek az an----gya-lok, ör-ven-dez-nek az ark-an--gya-lok, e na-pon          új-jon-ga-nak az i-ga-zak is, mond-ván: Di-cső-ség a ma-gas-ság-ban Is-ten-nek,      al-le-lu-ja!}

% --- Dominica infra octavam Nativitatis - S. Familiae Iesu, Maria et Ioseph ---

% ··· I. Vesperas ···

% Antiphona 2 (Adv132)
\defsheet[tone={7do-d}]{IosephFiliDavid}
  {<-8--6-----8--9---8--8--:-©-------------6---7---8--5--4--4-,-¦-----------------5u-6---:-6-5-----4--5---uU6-4--4-?,}
  {  Jó-zsef, Dá-vid fi-a, * ne_félj_magad-hoz ven-ni Má-ri-át, mert_aki_tőle_szü-le-tett, a Szent-lé-lek-től va-ló.}

% ··· II. Vesperas ···

% Antiphona 3 (Kar122)
\defsheet[tone={7la-a}]{PuerIesus}
  {<ô-1-4----4---5z-5---:-rR3-2--4t-4-,--4---5---3-----4----wW1-2-:-2--qQ0-1---2--3--2--1---0q-1-?,}
  {   A gyer-mek Jé-zus * nö--ve-ke-dett kor-ban s_böl-cses-ség-ben az Is--ten és em-be-rek e--lőtt.}

% --- Die 26 decembris ---

% ··· Vesperas ···

% Magnificat (Kar121)
\defsheet[tone={8szo-c}]{DumMediumSilentium}
  {<-4--4---34t--4-----:-4--4--3----4t-eE2W1-12e-:-1----3r-5---4--5--3--5---7--6--7---tT4R3-3---,-3--4---5----4---3-4t-5---:-4---3---4--5--eE2-1-:-4--5--7---5---4---3----4t-5--4x-4x-?,}
  {  Mi-kor mély csend * bo-rí-tott be min---dent, s_az éj-sza-ka út-já-nak fe-lé-nél tar---tott, ak-kor jött el, ó U--runk, min-den-ha-tó i---géd ki-rá-lyi szé-ked-ből, al-le-lu-ja.}

% --- Die 29 decembris ---

% ··· Officium lectionis ···

% Antiphona 2 (Gondolkodó 4.)
\defsheet[tone={3-c}]{Orietur}
  {<-4u-7---:-7--7--4---3--5u-7-:-7-7--7---5---4---5--4--4-,-4---4-3---rR3-2-?,}
  {  Ki-kél * az Úr nap-ja-i--ban a bé-kes-ség-nek bő-sé-ge, s_Ő u-ral-ko--dik.}

% --- Die 30 decembris ---

% ··· Officium lectionis ···

% Antiphona 1 (Gondolkodó 4.)
\defsheet[tone={8szo-c}]{VeritasDeTerra}
  {<-4-7--uJ4-:-6---7--5---5----4-tT4--3--,-3--5--7-7---7-:-4-6--7--5----5-4-----4-?,}
  {  A hű-ség * fel-nö-vek-szik a föld-ből, és az i-gaz-ság a-lá-te-kint a menny-ből.}

% --- Die 31 decembris ---

% ··· Officium lectionis ···

% Antiphona 1 (Gondolkodó 4.)
\defsheet[tone={4-a}]{LaetenturCaeli}
  {<-1--2---4--4---4--5z-5---:-4----5---6---5---4-5---,-5--2t-4---Ï3-wW1-2wW1-:-0q---2--4---2--2-?,}
  {  Ví-gad-ja-nak az e--gek * s_ör-ven-dez-zen a föld, az Úr szí-ne e---lőtt   mert im-már el-jött.}

% Antiphona 3 (Gondolkodó 4.)
\defsheet[tone={6-a}]{Notum}
  {<-3---3--3---2--4--55-¨-3--3--1--0-,-3--3-3--2---4--55---5--tT4-rR3-3-?,}
  {  Meg-mu-tat-ta az Úr * al-le-lu-ja, az ő üd-vös-sé-gét, al-le--lu--ja.}

% ··· Laudes ···

% Benedictus
\defsheet[tone={7la-d}]{FactaEstCumAngelo}
  {<-4--4---4i--8---:-8--8--7-----6-8----8---8-9--8--7---9--9--iI7-7-:-©----------oO8--7---8o-iI7--7--,-7--tT4-4-:-4-4--5---tT4-3---4--4t--4--,-3--5-7---7---7--8o--8-:-£-----------5--7--tT4-3----4t-5--4--4-?,}
  {  Kö-rül-vet-te  * az an-gyalt a meny-nye-i se-re-gek so-ka-sá--ga, dicséretet zeng-vén és mond-ván: Di-cső-ség a ma-gas-ság-ban Is-ten-nek, és a föl-dön bé-kes-ség a_jóakaratú em-be-rek-nek, al-le-lu-ja.}

% --- Die 1 ianuarii -  In sollemnitate Sanctae Dei Genetricis Mariae ---

% ··· I. Vesperas ···

% Antiphona 1 (Kar128)
\defsheet[tone={6-a}]{OAdmirabileCommercium}
  {<X-3ed1Q0-:-¢----------3r--4--4z----rR3-3---:-¢---------4---3---4--4t-3y-3y-4zZ5T4-4---6-----5---4--3-:-4--3----3----4---3---3--3-3zZ5T4tT4-4--,-4----4t--3--3----eE2-q\\Q0\\-:-0--1----3-3---3--4t--4-:-tT4R3-4tT4-3y-3y-?,}
  {   Ó        csodálatra mél-tó szent cse-re: * az_emberi nem-nek Al-ko-tó-ja tes----tet s_lel-ket vé-ve  ke-gyes volt szü-let-ni a Szűz------től, s_ki-lép-ve mint em--ber,      ne-künk a-ján-dé-koz-ta  is----ten--sé-gét.}

% ··· Laudes ···

% Antiphona 3 (NZS33)
\defsheet[tone={2-g}]{GenuitPuerperaRegem}
  {<ô-4----3--4-----2------3-2--1--3---2y-----:-2--2--wW1--2---4-2eE2W1-1í-,-2----ñw-2-2----rR3-2--3r-5zZ5¨6uU6c-:-6--5----4--3--4----5-3--rR3-2y-,-2---4----3---4---3--2---eE2-1w-2wW1×ñ-:-1--2-rR3-1--3r-2y--2y-:-ËqQ0-1--2y-2y-?,}
  {   Nagy Ki-rályt szült, í-me az Asz-szony, * ne-ve fény-lik ö-rök----ké,  mint jó a-nya, ví--ga-do-zott,        de szűz ma-ra-dott i-ga-zá--ban. Nem volt ily cso-da még so--ha-sem,     az-u-tán se ha-son-ló,  al---le-lu-ja.}

% Benedictus
\defsheet[tone={1re-a}]{MirabileMysterium}
  {<-1---3--1--0---3----4--3t-5--:-5--4--2---3---4-eE2-1í-1í--,-1---3-1---0-3--4t---5---:-5----X6-5---4--4----4--3t-4--:-4-4--2-----4--5---4--3---:-2-----3---4--3-2--1---q\\Q0\\-,-0---1----3---3---3--4t--4---:-4-----3----4t--eE2-1---0q--1í-?,}
  {  Cso-dá-la-tos misz-té-ri-um * je-le-nik meg e nap né-künk: meg-ú-jul a te-remt-mény, mert em-ber-ré lett az Is-ten; a-mi volt, az meg-ma-radt, s_föl-vet-te a-mi nem volt,     nem szen-ved-vén ve-gyü-lést, s_nem szen-ved-vén meg-osz-lást.}

% ----------------------------------------------------------------------------
% III. Tempus Quadragaeismae
% ----------------------------------------------------------------------------

% ----------------------------------------------------------------------------
% a. Usque ad sabbatum hebdomadae quintae
% ----------------------------------------------------------------------------

% --- Pro cunctis diebus ---

% ··· Invitatorium ···

\defsheet{ChristumDominumProNobis}
  {<-0-----1ed1-2-23rR32¨2e1wW1-1--:-1--1--1eE2-3r-4t--4---4---3--4tT4R3-2eE2W1-1----,-0---012e-wW1--1eE2¨3r-4tT4R3r-eE2-.}
  {  Krisz-tus  U-run-----------kat, az ér-tünk kí-sér-tet-tet és gyö----tör----tet, * jöj-je---tek, i-------mád-----juk!}

% vel:
% TODO

% ··· Officium lectionis ···

% Hymnus (In officio dominicali)
\defsheet{ExMoreDoctiMystico}
  {<-3--4t----5----5--rR3-4----4t--5-:-4-2---3--4---eE2-1---1w-2---;-2----4---4t--tT4-wW1-3---wW1-1w-:-0--2----3r--2---qQ0-1w---wW1-1-?,-3-4t-5-5-rR3-4-4t-5-:-4-2-3-4-eE2-1-1w-2-;-2-4-4t-tT4-wW1-3-wW1-1w-:-0-2-3r-2-qQ0-1w-wW1-1-?,-3-4t-5-5-rR3-4-4t-5-:-4-2-3-4-eE2-1-1w-2-;-2-4-4t-tT4-wW1-3-wW1-1w-:-0-2-3r-2-qQ0-1w-wW1-1-?,-3-4t-5-5-rR3-4-4t-5-:-4-2-3-4-eE2-1-1w-2-;-2-4-4t-tT4-wW1-3-wW1-1w-:-0-2-3r-2-qQ0-1w-wW1-1-?,-3-4t-5-5-rR3-4-4t-5-:-4-2-3-4-eE2-1-1w-2-;-2-4-4t-tT4-wW1-3-wW1-1w-:-0-2-3r-2-qQ0-1w-wW1-1-?,-1wW1-0q-.}
  {  Meg-szen-telt ré-gi  szép szo-kás a böj-ti meg-tar-tóz-ta-tás;  tart-suk meg új--ra  gon-do--san, mi negy-ven nap-ra  szab-va  van. Mó-zes és a szent lát-no-kok pél-dá-ja már e-lénk ra-gyog, s_Krisz-tus, ki Úr min-den fö-lött, ma-gá-ra vett ily böjt-i-dőt. Kí-mé-lőn szól-jon hát sza-vunk, le-gyen sze-gé-nyebb asz-ta-lunk, ne győz-zön paj-zán tré-fa-ság, se ál-mos, tét-len lus-ta-ság. Ke-rül-jünk min-den vét-ke-set, szí-vün-ket rossz ne ront-sa meg, el-len-ség hogy-ha meg-kí-sért, hi-tünk a csap-dát zúz-za szét. Bol-dog Há-rom-ság, add meg ezt, osz-tat-lan Egy-ség, tedd meg ezt, ne-künk az üd-vös böjt-i-dő le-gyen sok jót gyü-möl-csö-ző. Á-men.}

% Hymnus (In officio feriali)
\defsheet{NuncTempusAcceptabile}
  {<-3--4t---5-5---rR3-4--4t-5-:-4--2----3---4---eE2-1----1w-2---;-2-4---4t---tT4--wW1-3---wW1-1w-:-0-2----3r-2---qQ0-1w--wW1-1-?,-3-4t-5-5-rR3-4-4t-5-:-4-2-3-4-eE2-1-1w-2-;-2-4-4t-tT4-wW1-3-wW1-1w-:-0-2-3r-2-qQ0-1w-wW1-1-?,-3-4t-5-5-rR3-4-4t-5-:-4-2-3-4-eE2-1-1w-2-;-2-4-4t-tT4-wW1-3-wW1-1w-:-0-2-3r-2-qQ0-1w-wW1-1-?,-3-4t-5-5-rR3-4-4t-5-:-4-2-3-4-eE2-1-1w-2-;-2-4-4t-tT4-wW1-3-wW1-1w-:-0-2-3r-2-qQ0-1w-wW1-1-?,-1wW1-0q-.}
  {  Ez most a ked-ve--ző i--dő, ke-gyel-mek nap-ja  ránk ra-gyog: a lan-kadt föld-re  hoz-za  már  a böjt a  bol-dog gyó-gyu-lást. Krisz-tus nagy fé-nye tün-dö-köl, üd-vünk-nek nyit-ja nap-ja-it, míg bűn-nel seb-zett szí-ve-ket fe-gyel-me gyó-gyít, jó-ra int. U-runk, tart-has-sunk üd-vö-sen test-ben-lé-lek-ben böj-tö-ket, s_u-tunk-on biz-tos fény le-gyen ö-rök Hús-vé-tod ö-rö-me. Ke-gyes Há-rom-ság, té-ge-det i-mád-ja-nak már min-de-nek, s_ki-ket me-gú-jít nagy ke-gyed, zeng-jük di-cső, új é-ne-ked. Á-men.}

% ··· Laudes ···

% In officio feriali

% Hymnus
\defsheet{IamChristeSolIustitiae}
  {<X-3--3---4---5---4---4-----3--2---:-4--2-----3----1---0--1----3--3--,-3----3---4--5----4---4--3-2---:-4--2---3--1---0---3----1--1-?,-02eE2W1Q0-1-.}
  {   Üd-vös-ség nap-ja, Krisz-tu-sunk, ha-sítsd szét vak-sá-gunk kö-dét, hogy fel-ra-gyog-jon az e-rény, mi-dőn re-ánk nap-palt ho-zol. Á---------men.}

% Responsorium breve
\defsheet{IpseLiberabitMe}
  {<X-3-3---3--4---3---4t-tT4-,-3-4--5--4----ed1--3r-4--3-?,-5--5-4---3--3--4t---5-?,-5--5---5-.}
  {   Ő sza-ba-dít meg en-gem * a va-dá-szok csap-dá-já-tól. És a dur-va be-széd-től. Di-cső-ség...}

% In officio dominicali
\defsheet{PrecemurOmnesCernui}
  {<X-3---3---4--5----4----4--3-2--:-4----2----3---1--0--1-3--3--,-3----3---4---5---4--4----3---2-:-4-2---3---1---0--3-----1--1-?,-02eE2W1Q0-1-.}
  {   Tér-del-ve mond-junk hő i-mát, zeng-jünk bűn-bá-nó é-ne-ket, tart-sák fel for-ró köny-nye-ink a vét-ket tor-ló szent ke-zet! Á---------men.}

% Responsorium breve
\defsheet{ChristeFiliDeiVivi}
  {<X-¢----------------4--3---4t-5--,-5--4---3----rR3-3--?,-[------------------4--5--4--:-5---rR3-4t--5-?,-5--5---5-.}
  {   Krisztus,_az_élő Is-ten Fi-a, * Ir-gal-mazz ne--künk. Akit_a_mi_gonoszsá-ga-in-kért tör-tek ösz-sze. Di-cső-ség...}

% ··· Tertia ···

% Hymnus
% TODO

% Antiphona
% TODO

% ··· Sexta ···

% Hymnus
\defsheet{QuaChristus}
  {<-0-1----2e-2-qQ0-1----2--1---:-1--2-----3---rR3-2--12e-eE2-2---,-2-4---4t---tT4--eE2--12e-eE2-2--:-0-1---2e--2-qQ0-1----2---1-.}
  {  U-runk ez ó-rán szom-ja-zott, ke-reszt-fán tes-te lan-ka--dott: a-kik most zsol-tárt mon-da--nak, a-zok-nak é-gi  szom-jat ad.}

% Antiphona
\defsheet[tone={tirr-b}]{VivoEgoDicitDominus}
  {<-uj5-4z--6---:-6----7--uU6-4t-,-7---7-7--6---5--4--5---5-rR3-2e-wW1-,-1--2---45z-5----5----rR3-2e--4---rR3-2-2---?,}
  {  É---lek én, * mond-ja az  Úr:  nem a-ka-rom ha-lá-lát a bű-nös-nek,  ha-nem in--kább hogy meg-tér-jen és  él-jen.}

% ··· Nona ···

% Hymnus
% TODO

% Antiphona
% TODO

% ··· Vesperas ···

% In officio dominicali

% Hymnus
\defsheet{AudiBenigneConditor}
  {<-qQ0-1---eE2-3-----2e---1--0-1---:-eE2-3r--rR3-1--3--1w--1--0--,-0-1----2----eE2-1--2---3-4--:-0q-3r---2---eE2---1---wW1-0--1-.}
  {  Sza-vun-kat halld meg, jó U-runk: a   hoz-zád es-dő szó-za-tot, a-mely-lyel sír-va kér i-mánk e  negy-ven szent nap al--ko-nyán.}

% In officio feriali

% Hymnus
\defsheet{IesuQuadragenariae}
  {<-3--2--4----tT4---3----4t--tT4R3-2-:-1--1t--5---5---tT4-4z-tT4-4t--;-5--5----4--tT4-3---4----eE2-1--:-3----2--4---tT4-3---4t-tT4R3-2-.}
  {  Úr Jé-zus, szent negy-ven-na-pos    böj-töd vál-lal-juk új-ra most: üd-vünk-re pél-dát adsz ne-künk, hogy fé-kez-zük ter-mé-sze-tünk.}

% Responsorium breve
\defsheet{EgoDixi}
  {<X-3---4----3----4t-5---,-4--5---rR3--4---3-?,-5---5-----5---4--4--,-4----4---4---5---rR3-4t-5-,-5--5---5-.}
  {   Így szól-tam: U-ram, * Kö-nyö-rülj raj-tam! Gyó-gyíts meg en-gem, mert vét-kez-tem el-le-ned. Di-cső-ség...}

% ············································································
% Hebdomada II Quadragesimae
% ············································································

% --- Feria VI ---

% ··· Vesperas ···

% Magnificat
\defsheet[tone={3-c}]{QuaerentesEumTenere}
  {<-4--5---6--7-5---7---6----:-4--4---5---5-zZ5-4--,-6----6---5--uU6-4-:-5--rR3-rR3-2-?,}
  {  El-fog-ni i-gye-kez-tek, * de fél-tek a nép-től, mert pró-fé-tá-nak tar-tot-ták őt.}

% --- Sabbato ---

% ··· Laudes ···

% Benedictus
\defsheet[tone={1la-a}]{VadamAdPatremMeum}
  {<X-1--1t-5----:-7-5----4-:-rf2-23r-eE2-1--1-,-3-123E2s0-:-34tT45z-5----4t-4---3---3---2e-4--eE2-1--1-?,}
  {   El-me-gyek * a-tyám-hoz és  mon-dom ne-ki: A-tyám,     fo------gadj en-gem nap-szá-mo-sa-id  kö-zé.}

% ············································································
% Hebdomada III Quadragesimae
% ············································································

% --- Dominica ---

% ··· I. Vesperas ···

% Antiphona 1
\defsheet[tone={8szo-c}]{ConvertiminiAdDominum}
  {<-4---4u--5---rR3-4t-4----:-8--7--5--7---7--,-©----------------8Oo8-uU67iI7-:-5---5--4----7--7-4----4--?,}
  {  Tér-je--tek az  Úr-hoz, * Is-te-ne-tek-hez: jóságos_és_irgal-mas  ő,        meg-bo-csát-ja a rosz-szat.}

% Antiphona 2
\defsheet[tone={8szo-c}]{Sacrificabo}
  {<-3--5---7-zh4-:-6--7--5--4---4--4--tT4-3-,-3t-7--zh4-6u-5---4--4--?,}
  {  Di-cső-í-tő  * ál-do-za-tot mu-ta-tok be, és az Úr  ne-vét hí-vom.}

% Antiphona 3
\defsheet[tone={3-c}]{NemoTollitAMe}
  {<-4---5z-¥------------------:-7--uU6-5--7--uU6-,-7--8---8--6---7-5---5-4-:-4--4--4t--eE2-1----2e---4---2-?,}
  {  Sen-ki sem_veszi_el_tőlem * az é---le-te-met,  ha-nem ma-gam a-dom o-da  és ma-gam ve--szem visz-sza azt.}

% ··· Laudes ···

% Antiphona 1
\defsheet[tone={1re-a}]{RegnavitDominus}
  {<-3--1---1---1t-5---:-5--4--5---rR3-34t-eE2-1---:-1--1eE2-3---4-5---rR3-2e-4---eE2-1--1-?,}
  {  Or-szá-gol az Úr, * ha-ta-lom-ba  öl--tö--zött, és most fel-ö-vez-te  ma-gát e---rő-vel.}

% Antiphona 2
\defsheet[tone={1la-a}]{VimVirtutisSuae}
  {<-0--1---1-1--1tu-5---:-[------------tT45uj5T4-,-5----4--3-3--3---3---3--3---3r-3-:-2e--4---eE2-1-?,}
  {  Sa-ját e-re-jé--ből * megkezdett_a tűz,        hogy fi-a-id meg-sza-ba-dul-ja-nak sér-tet-le--nül.}

% Antiphona 3
\defsheet[tone={2-g}]{RegesTerrae}
  {<-3----3--2--1---1-:-2--3---4---5--4--:-2--3----4--3---2--1--1-?,}
  {  Föld ki-rá-lya-i * és min-den né-pek, di-csér-jé-tek az Is-tent!}

% Benedictus (Anno A)
\defsheet[tone={8do-c}]{Aqua}
  {<-7-uU6--:-tT4-5z-tT4--5-4--,-4-4--2--4---eE2-3---:-1---2e---4t-5---4---4-?,}
  {  A víz, * a---me-lyet a-dok, a-ki ab-ból i---szik, nem szom-ja-zik töb-bé.}

% ············································································
% Sacrum triduum paschale
% ············································································

% --- Feria V in Cena Domini ---

% ··· Vesperas ···

% Responsorium breve
\defsheet{ChristusFactusEstBrevis}
  {<X-33r---3-:-¢---------------4tg3E2¨11rf2¨33r-3-:--3t¨34t77¨iI7i™uU6u¨55-44-:-4-5--7uj5g3¨zh4tT4t¨3rR3-3-.}
  {   Krisz-tus engedelmes_volt ér---------------tünk mind          @       @    a ha-lá------------------lig.}

% --- Feria VI in Passione Domini ---

% ··· Vesperas ···

% Responsorium breve
\defsheet{ChristusFactusEstLonga}
  {<X-33r---3-:-¢---------------4tg3E2¨11rf2¨33r-3-:--3t¨34t77¨iI7i™uU6u¨55-44-:-4-5--7uj5g3¨zh4tT4t¨3rR3-3-:--3-3--5z----445z¨uU6h44t¨zZ5T4-3--3--3--eD0¨3r¨tT4R3E24t3tT4t¨rR3-.}
  {   Krisz-tus engedelmes_volt ér---------------tünk mind          @       @    a ha-lá------------------lig, a ke-reszt-nek                ha-lá-lá-ig.}

% --- Dominica Paschae in Resurrectione Domini ---

% ··· Completorium ···

% Responsorium breve
\defsheet{HaecDies}
  {<-3--3--4-3--:-3--3----3--3--3----2----4--5---,-5--tT4-rR3-3---:-1----3---4----rR3-3-.}
  {  Ez az a nap, me-lyet az Úr szer-zett ne-künk, ör-ven-dez-zünk, s_ví-gad-junk raj-ta.}

% ----------------------------------------------------------------------------
% IV. Tempore Paschale
% ----------------------------------------------------------------------------

% ············································································
% b. Usque ad Ascensionem Domini
% ············································································

% --- Pro cunctis diebus ---

% ··· Hora Media ···

% Hymnus (Sexta)
\defsheet{VeniteServiSuplices}
  {<-5-5---5--5----5----5---4--4--,-5---5--5---4---3--4---5--5--,-5----5--5---4---3--4--3--2--,-3----4---4---3--4----5---4--4-?,-4tT4-3\\r\\-.}
  {  E-sen-gő szol-gák, jöj-je-tek, tár-já-tok fel ma szí-ve-tek, nyis-sá-tok dal-ra aj-ka-tok, hogy bol-dog Is-tent áld-ja-tok! Á----men.}

% --- In Ascencione Domini ---

% ··· Vesperas ···

% Hymnus
\defsheet{IesuNostraRedemptio}
  {<ô-2--2----1---2---3---4---2--2-:-2-3---4---3---1--2----4--334tT4R3y-,-2--3---4--5---3--4--3--2-:-2-3---4---3---2--3----2--1-----?,-1wW1-Ë0\\q\\-.}
  {   Jé-zus, meg-vál-tás kút-fe-je, a-kit sze-ret-ve szom-ja-zunk,       vi-lág te-rem-tő Is-te-ne, i-dők tel-jén ki föld-re szállsz, Á----men.}

% Antiphona
\defsheet[tone={8szo-c}]{AlleluiaTPSexta}
  {<-7--7--7--7---¨-5--4--5u-7-:-7--tT4-3r-4-?,}
  {  Al-le-lu-ja, * al-le-lu-ja, al-le--lu-ja!}

% ----------------------------------------------------------------------------
% V. Tempus per Annum
% ----------------------------------------------------------------------------

% ············································································
% In Dominicis
% ············································································

% --- Ad I. Vesperas (Cyclus unius anni) ---

% Magnificat
\defsheet[tone={1la-a}]{SustinuimusPacem}
  {<-0q--0---1-eE2-qQ0-:-2e-rR3-2eE2W1-0q-1--,--tT4-5---3--4---rR3-:-2--4t-wW1-3--wW1-0q--1--,-5--4--5---uU6--uj5-tT4R3-4--5--5--qQ0-,-2e-4--3--wW1-:-2-3---4t-qQ0-0---:-3--3--eE2-1wW1-0q-1-?,}
  {  Vár-tuk a bé-két, * és nem jött,  U--runk; a   jót ke-res-tük,  és í--me, za-var van itt: el-is-mer-jük, U---runk, bű-ne-in-ket;  ne ha-ra-gudj  ö-rök-ké mi--ránk, Iz-ra-el  Is---te-ne.}

% --- Anno A ---

% Dominica XIV

% Benedictus
\defsheet[tone={1re-a}]{ConfiteorTibiPater}
  {<X-0--1---1t-5u-5----:-tT45uj5-4---:-3--4--5----rR3r-1-,-3----2--Ÿ------------------------------2---3---4t-4--3--,-¢-----------------2--3---4---eE2-1---1-?,}
  {   Di-cső-í--te-lek, * A-------tyám, ég és föld U----ra, mert az okosak_és_a_bölcsek_elől_elrej-tet-ted e--ze-ket, és_a_kicsinyeknek ki-nyi-lat-koz-tat-tad.}

% Dominica XV

% Benedictus
\defsheet[tone={8szo-c}]{CumTurbaPlurima}
  {<-4-7--5---rR3--4t-4---:-7---7--8--uj5-7--7---,-©-------------------7--8--8oO8-uU67iI7-,-5---5--5--7---7---4---4---:-4--4----4-4---eE2-1-:-3----3--4---5--4---4-?,}
  {  A-mi-kor nagy tö-meg * vet-te kö-rül Jé-zust, és_a_városokból_hoz-zá si-et---tek,      pél-da-be-szé-det mon-dott: Ki-ment a mag-ve--tő, hogy el-ves-se mag-ját.}

% Dominica XXI

% Benedictus
\defsheet[tone={8szo-c}]{QuemDicuntHomines}
  {<-rF1-eE2-3r---4---:-4--4t-tT4-3---5--5zZ5-4---3r-4-;--4uU6-7i-iI7-5---7-uU6h4-5---rR3-5---zZ5-4--,-4--4t--tT4tT4R3-5uu-5zZ5-4---4-;--rR3d13r-r3ttT4tg3r35u-7i----8--:-8--uU6-4tT4R3-5uj5-zZ5-4--4-,--4--4t-5---4t--tT4R3-3--:-3t---7i78o-uU6--5--5--;-5i-8--uU6-5-4--4t---rR3-3-:-3--3---4t-5èu-zZ5-5-57iiI7-tT4R3-5zZ5-4-?,}
  { „Ki--nek tart-ják * az em-be--rek az Em---ber-fi-át?” mond-ta Jé--zus a ta----nít-vá--nya-i---nak. Fe-lel-vén      Pé--ter, mon-dá: „Te      vagy          Krisz-tus, az é---lő     Is---ten Fi-a.” „Én is mon-dom ne----ked, hogy te    vagy Pé-ter, és er-re  a kő-szik-lá--ra  fo-gom é--pí--ten-i egy----há----za---mat.”}

% Magnificat
\defsheet[tone={8szo-c}]{Quodcumque}
  {<-3r-4--¨-5---4-----3r-2e--1--;-1---0----34t-5--5uj5-5z----4---4-,-7--uj5-7i--8---8o-uU6--4-34t-5--;-7u8µ4-4----5--rf2-3--4-2-----3---1-;--5----tT4-5--5--5--5uj5-zZ5-4---4-?,}
  { „A-mit * meg-kötsz a  föl-dön, meg lesz köt-ve a    menny-ben is; és a---mit fel-ol-dasz a föl-dön, fel   lesz ol-doz-va a menny-ben is,” mond-ta  az Úr Si-mon  Pé--ter-nek.}

% --- Anno C ---

% Dominica XV

% Benedictus
\defsheet[tone={8szo-c}]{HomoQuidam}
  {<-4---4u-5---:-[------------------rR3-4t-4---:-7---7---8--7--5--7--7---:-¦------------------8--7---8oO8-uU67i-7--:-5--4---5---7---7---4----4-?,}
  {  Egy em-ber * Jeruzsálemből_Jeri-kó--ba ment. Rab-lók ke-zé-be ke-rült. Ezek_kifosztották, vé-res-re   ver---ték, és fél-hol-tan ott-hagy-ták.}

% Dominica XIX

% Benedictus
\defsheet[tone={8szo-c}]{NoliteTimere}
  {<-4--4u----5--:-3--4t----5--4---,-8----8---8o--8-----7-7----uU67iI7-:-5----5---5---5--7--7--4---4-?,}
  {  Ne félj, te * ki-csiny-ke nyáj, mert úgy tet-szett A-tyád-nak       hogy nek-tek ad-ja or-szá-gát.}

% ············································································
% In Solemnitatibus Domini
% ············································································

% --- Sanctissimi Corporis et Sanguinis Christi ---

% I. Vesperas

% Antiphona 1
\defsheet[tone={2-g}]{MiseratorDominus}
  {<ô-4--4--4---2---rR3-2---:-4-4--4----2-4----5z-5x--,-5--5--5--6--5---rR3--2--:-2---2--1---ñ--1w-2-?,}
  {   Az ir-gal-mas Is--ten * e-le-delt a-dott né-künk, em-lé-ke-ze-tet szer-zett cso-da-tet-te-i--nek.}

% Antiphona 2
\defsheet[tone={5-g}]{QuiPacem}
  {<-2--0--2---4t-4----:-4--4---5---7--6--tT4-4-,-4-5--4--3--4t-rR3E2-2-:-2-1w-eE2-1--0í--0í-?,}
  {  Ki bé-két te-remt * az Egy-ház ha-tá-ra--in, a ga-bo-na ja-vá----val e-lé-git ki min-ket.}

% Magnificat
\defsheet[tone={1la-a}]{OQuam}
  {<-3--0--1----1tu-5---:-[----------------tT45uj5T4-:-4----5---4---3--3r-3-:-3-2--3-4---3---2--1e--3-:-3--4---5---3----2--1--qQ0-:-1--eE2-1í-1í-?,}
  {  Ó, mi-lyen é---des * hozzánk_a_te_lel-ked:        hogy ked-ves-sé-ge-det a fi-a-kon meg-mu-tas-sad ke-nye-ret adsz az ég-ből,  al-le--lu-ja.}

% Laudes

% Antiphona 1
\defsheet[tone={7do-a}]{AngelorumEsca}
  {<ô-5--5---3---5-6--7--zZ5-5---:-5--5--3---4-5--wW1-1--,-£------------2r----3-:-2--1---2---4----3----1í-1í-?,}
  {   An-gya-lok e-le-de-lé--vel * él-te-ted a Te né--ped, s_kenyeret_a menny-ből bő-ség-gel nyúj-tasz né-kik.}

% Hora Medias

% Antiphona (tertia)
\defsheet[tone={4-a}]{Desiderio}
  {<-1----2--4t---5----:-4----5-6---5--4--5--5--5x-,-2--5----4---Ï3--2--1--wW1-:-0--1-2r---2----2--2y-.}
  {  Vágy-va vágy-tam, * hogy e hús-vé-ti la-ko-mát  el-költ-hes-sem ve-le-tek,  mi-e-lőtt szen-ve-dek.}

% --- Sanctissimi Cordis Iesu ---

% Laudes

% Antiphona 2
\defsheet[tone={1re-a}]{VeniteAdMe}
  {<-3---0--1---2tu-5----:-5--5----5---tT45uj5T4-:-4-4---5--4--3r--3--:-2-----3r--eE2-1í---1í-,-1----3---wW1-2--3----4t-4--3--:-2----3--4--eE2-1í----1í-.}
  {  Jöj-je-tek ho--zám, * ti mind-nyá-jan,        a-kik fá-ra-doz-tok, s_ter-hel-ve  vagy-tok, s_én meg-eny-hít-lek ti-te-ket, mond-ja az Úr  Krisz-tus.}

% Hora Media

% Antiphona (tertia)
\defsheet[tone={7do-a}]{PopuleMeus}
  {<ô-5tg3-56uj5-5----:-5---5--3---4--5--wW1-1í-,-£-------------------2r-3y-:-2--1----2rR3-1í-1í-?,}
  {   Én   né----pem, * mit te-tem én el-le--ned? Vagy_miben_ártottam né-ked? Fe-lelj meg  né-kem!}
\makeatletter

\makeatletter

% ············································································
% Hymnus
% ············································································

\newcommand{\sheet@TeLucisAnteTerminum}{%
  \begin{gregosheet}
    \melody{<-7--7---7-7---7--7--6--5---:-7--7---7-----7--4----5--4--3--,-5----5--6----7--6--5----4--5zZ5-:-5-5---4----3--4----5--4--4-.}
    \lyrics{  Im-már a nap le-ál-do-zott, Te-rem-tőnk, ké-rünk té-ge-det, légy ke-gyes és ma-radj ve-lünk,  ő-riz-zed, óv-jad  né-pe-det!}
  \end{gregosheet}
}

% vel:
\newcommand{\sheet@ChristeQuiSplendor}{%
  \begin{gregosheet}
    \melody{<-7-----7----7---7---7--7---6--5---:-7--7--7--7----4-5----4--3--,-5---5----6---7----6---5----4-5zZ5-:-5--5----4----3--4---5-----4--4-.}
    \lyrics{  Krisz-tus, tün-dök-lő nap-pa-lunk, az éj ár-nyát e-losz-la-tod, kit fény-ből fény-nek vall a hit    és fény-nyel ál-dod szent-je-id.}
  \end{gregosheet}
}

% ············································································
% Responzorium breve
% ············································································

% Responsorium breve in Tempore Paschali
% TODO

% Responzorium breve
\newcommand{\sheet@InManusTuasDomine}{%
  \begin{gregosheet}
    \melody{<X-3r-3---,-4--4---4--3-4t--4---ed1-3r-3-,,-5---5---5z--5---5---4---4t-4-:--5--4--4---rR3-4-4t-,,-5--5---5---5--5-5z--5-:--5-5--4-4-:--5--5-4-----3--4t--tT4-.}
    \lyrics{   U-ram, * ke-zed-be a-ján-lom lel-ke-met. Meg-vál-tot-tál min-ket U-runk, hű-sé-ges Is-te-nünk. Di-cső-ség az A-tyá-nak, a Fi-ú-nak, és a Szent-lé-lek-nek.}
  \end{gregosheet}
}

% Nunc dimittis
\newcommand{\sheet@SalvaNos}{%
  \begin{gregosheet}
    \melody{<X-1-----3----1-0----3---4----3--4t---5x---:-6-5----4---4--4--3t---4y-,-4----4-4-----2---4---5---4----3y---:-3--2--3---4---eE2-1---0q--1í-,,-3----4--[--------------------4-3----,-[-----------------4--3---4t--4-,,-5----5---5---4---3--3--,-5-5--5--4---3--4t--4-,,-5-5---4---3---3---,-5---5---5--5---4---3--4t-4--,,-3r---[------------------------5---5--4--3---,-[------------------------4--3--4t-4-,,-3--4---[-----------------4--3-3---,-5--5-----4--3---4--4t--4-,,-3--4---[---------------------------4---3---3---,-[----------------4---3---4t-4rR3E2W1í-,,-1-----3-.}
    \lyrics{   Tarts meg, U-ram, vir-rasz-tá-sunk-ban, * ő-rizz meg al-vá-sunk-ban, hogy a Krisz-tus-sal vir-rasz-szunk, és bé-kes-ség-ben nyu-god-junk! Most bo-csátod_el_szolgádat, U-ram, * a_te_igéd_szerint bé-kes-ség-ben. Mert lát-ták sze-me-im * a te üd-vös-sé-ged-et.  A-kit ren-del-tél * min-den né-pek szí-ne e--lőtt, Hogy fény_legyen_a_népek_vilá-gos-sá-gá-ra, * és_dicsőségére_népednek, Iz-rá-el-nek. Di-cső-ség_az_Atyának_és Fi-ú-nak * és Szent-lé-lek Is-ten-nek, Mi-kép-pen_kezdetben_vala,_most_és min-den-kor * és_mindörökkön_ö-rök-ké! Á--men.         Tarts meg...}
  \end{gregosheet}
}

% ············································································
% Dominica Post I Vesp.
% ············································································
% Antiphona 1
\newcommand{\sheet@MiserereMihiDomine}{%
  \begin{gregosheet}[tone=8. szó tónus]
    \melody{<-4--4---4----5---4----tT4-3y---,-5--7----6u---5x-:-5-4---3--4t-4y-,,-7-7-6-7-5-4-.}
    \lyrics{  Kö-nyö-rülj raj-tam, U---ram, * és hall-gasd meg  i-mád-sá-go-mat!}
  \end{gregosheet}
}

% Antiphona 2
\newcommand{\sheet@InNoctibus}{%
  \begin{gregosheet}[tone=4. tónus]
    \melody{<X-0--1--3r--4-:-5---6--5---4--4t-4-,,-7-7-6-7-6-4-.}
    \lyrics{   Éj-je-len-te  áld-já-tok az U--rat!}
  \end{gregosheet}
}

% ············································································
% Dominica Post II Vesp.
% ············································································
% Antiphona
\newcommand{\sheet@QuiHabitat}{%
  \begin{gregosheet}[tone=8. szó tónus]
    \melody{<-4-4--4---4---4t-tT4-3---:-3--5---5u-7---4--3r-4--,-5-7---7---6u-5-:-5--4---3--4t--4--4-,,-7-7-6-7-6-4-.}
    \lyrics{  A ma-gas-ság-be-li--nek * ol-tal-má-ban la-ko-zik, a min-den-ha-tó  ár-nyé-ká-ban ma-rad.}
  \end{gregosheet}
}

% ············································································
% Feria II
% ············································································
% Antiphona
\newcommand{\sheet@IntendeVoci}{%
  \begin{gregosheet}[tone=8. szó tónus]
    \melody{<-4--6----uU6-5z--:-5---5--uj5-zZ5-4-,,-7-7-6-7-5-4-.}
    \lyrics{  Kö-nyör-gé--sem * sza-vá-ra  fi--gyelj.}
  \end{gregosheet}
}

% ············································································
% Feria III
% ············································································
% Antiphona
\newcommand{\sheet@NeIntres}{%
  \begin{gregosheet}[tone=8. szó tónus]
    \melody{<-4t-5------:-5u--7i-7----6---5----4-4-,,-7-7-6-7-5-4-.}
    \lyrics{  Ne szállj * per-be szol-gád-dal, U-ram.}
  \end{gregosheet}
}

% ············································································
% Feria IV
% ············································································
% Antiphona 1
\newcommand{\sheet@EstoMihi}{%
  \begin{gregosheet}[tone=6. tónus]
    \melody{<-3----4--5----rR3-3----:-3--1---3--4--4--3--3-,,-6-6-5-4-5-4-.}
    \lyrics{  Légy né-kem, U---ram, * ol-tal-ma-zó Is-te-nem!}
  \end{gregosheet}
}

% Antiphona 2
\newcommand{\sheet@DeProfundis}{%
  \begin{gregosheet}[tone=8. szó tónus]
    \melody{<-7-7----7--7---5---:-5--4--5---6---5----4--3r-4-,,-7-7-6-7-5-4-.}
    \lyrics{  A mély-sé-gek-ből * ki-ál-tok hoz-zád, ó, U--ram}
  \end{gregosheet}
}

% ············································································
% Feria V
% ············································································
% Antiphona
\newcommand{\sheet@CaroMea}{%
  \begin{gregosheet}[tone=4. tónus]
    \melody{<-tT4-3---:-4--3----wW1--3--rR3--2---2-,,-5-5-4-5-4-2-.}
    \lyrics{  Tes-tem * el-nyug-szik re-mény-ség-ben.}
  \end{gregosheet}
}

% ············································································
% Feria VI
% ············································································
% Antiphona
\newcommand{\sheet@DomineDeusSalutisMeae}{%
  \begin{gregosheet}[tone=4. tónus]
    \melody{<-3-1----3---3--4--3--4--4t-4--,-3--2---2e--4-:-3--2---3---4--3--2-,,-5-5-4-5-4-2-.}
    \lyrics{  U-ram, sza-ba-dí-tó Is-te-nem, éj-jel-nap-pal te-hoz-zád ki-ál-tok!}
  \end{gregosheet}
}

% ············································································
% Antiphonae finales ad Beatam Mariam Virginem
% ············································································
% Tempus Quadragaesimae
% TODO

% Simpl.
\newcommand{\sheet@AveReginaCaelorum}{%
  \begin{gregosheet}
    \melody{<X-3-2--1--0--1--3---4--3----:-5-7---6--4--5--4--3--5--4--;-3---2--1--0----1---3---4---3-:-5--4---6---5--4---1---4--3-,-3---4---5---5--4---5--6-5-:-7--6---5--4---3---1--4-3-,-6--5---4-6---5--4--34t-5-;-7-6---67i-5--5-----4---3--44-3-.}
    \lyrics{   A-ve Re-gi-na cae-lo-rum, * a-ve, Do-mi-na an-ge-lo-rum, sal-ve ra-dix, sal-ve, por-ta, ex-qua mun-do lux est or-ta. Gau-de, Vir-go glo-ri-o-sa, su-per om-nes spe-ci-o-sa; va-le, o val-de de-co-ra, et pro no-bis Chris-tum ex-o-ra.}
  \end{gregosheet}
}

% Tempus Pachale
% TODO

% Simpl.
\newcommand{\sheet@ReginaCaeli}{%
  \begin{gregosheet}
    \melody{<X-3--4--3--4---5--:-6---5--4-:-6--5--4--3-,-3---7-7----8--7--6--5--4---5--6-:-6--5--4--3-,-7--7---8---7-:-7--3---4--3--:-4--5--6--7-,-7-3--4---6--5---4--3-:-2--4--4--3-.}
    \lyrics{   Re-gi-na cae-li * lae-ta-re, al-le-lu-ja: Qui-a quem me-ru-is-ti por-ta-re, al-le-lu-ja: Re-sur-rex-it, si-cut di-xit, al-le-lu-ja: O-ra pro no-bis De-um, al-le-lu-ja.}
  \end{gregosheet}
}

% Tempus Per Annum
% TODO

% Simpl.
\newcommand{\sheet@SalveRegina}{%
  \begin{gregosheet}
    \melody{<-0---2---4--5--4---,-5--7---6--5--4--5---4--4-,-7--4---5---33-1-:-2--3----4---2----wW1-0-,,-4--5--6---7--4-:-5--6--7---6--5--4-5-4-,,-7--4--5---3--1--2--:-2--4---5---7--5----4-:-5--4---3---2--1--2---1---0-,,-4-5---6--7-:-4--5--7--6--5---4--,-7--4---5--3--1--2--3--4---3---5-4--4-:-3--2---1---wW1-0-,,-4--5z-7--:-6--4--5---4---4----5---7----6---5--4-,-0--4---5----7---6--5--4--2--3--wW1-0-,,-2å4-22--0--,,-4ç6uU6-tT4-4-,,-77¨4tg3¨ed1-2e--4-:-0---3--2--qQ0-0-.}
    \lyrics{  Sal-ve, Re-gi-na, * ma-ter mi-se-ri-cor-di-ae, vi-ta, dul-ce-do, et spes nos-tra, sal-ve.  Ad te cla-ma-mus ex-su-les fi-li-i E-vae. Ad te sus-pi-ra-mus, ge-men-tes et flen-tes in hac lac-ri-ma-rum val-le.  E-ia, er-go, ad-vo-ca-ta nos-tra, il-los tu-os mi-se-ri-cor-des o-cu-los ad nos con-ver-te.  Et Ie-sum, be-ne-dic-tum fruc-tum vent-ris tu-i,  no-bis post hoc ex-si-li-um os-ten-de.  O   cle-mens, o      pi--a,   o           dul-cis Vir-go Ma-ri--a.}
  \end{gregosheet}
}

\makeatother

% ============================================================================
% DOMINICA
% ============================================================================

% ············································································
% I. Vesperas
% ············································································

% Hymnus
\defsheet{DeusCreatorOmnium}
  {<ô-2--2----2--1--2---4--3--2í-:-2--3--4---5---4---5---5--6x-,-6--6---5---6-4---5----rR3-2í-:-2--1---2--4----5---2--1---2í-?,-1e4R3E2W1-2í-.}
  {   Is-ten, vi-lá-got al-ko-tó,  vi-lá-got sar-kán for-ga-tó,  ki nap-nak é-kes fény-ru--hát, éj-nek ke-gyel-mes ál-mot adsz. Á---------men.}

% Antiphona 1
\defsheet[tone={8szo-c}]{DomineClamavi}
  {<-¦-----------5--4----45u-7c---:-5--4---5---4x-?,}
  {  Tehozzád_ki-ál-tok, U---ram, * si-ess hoz-zám!}

% Antiphona 2
\defsheet[tone={8szo-c}]{PortioMeaDomine}
  {<-¦-------------5--4----45u--7c--:-7--7-5---4---5----4x-?,}
  {  Osztályrészem te légy U----ram * az é-lők-nek föld-jén!}

% Antiphona 3
\defsheet[tone={1la-a}]{HumuliavitSemetipsum}
  {<-1---3-1---0--3--4t-5---:-XzZ5-4--3t-4--,-2-4----tT4-3-:-¢------------2--3--4----eE2-1---1-?,}
  {  Meg-a-láz-ta ön-ma-gát * az   Úr Jé-zus, e-zért Is--ten felmagasztal-ta őt mind-ö---rök-re.}

% Responsorium breve
\defsheet{QuamMagnificataSuntRBr}
  {<X-3----4----3---4t-5--,-5-5--5--4--3---rR3-3--?,-5---5----5z--4-:-4--3---4t-?,-[-------------5z--5-:-5-5--4-4-:-5--5-4-----3--4t--tT4-.}
  {   Mily nagy-sze-rű-ek * A te mű-ve-id, U---ram.  Böl-cses-ség-gel al-kot-tad.  Dicsőség_az_A-tyá-nak a Fi-ú-nak és a Szent-lé-lek-nek.}

% ············································································
% Invitatorium
% ············································································

\defsheet{VeniteExsultemusDomino}
  {<X-3---4t-5----:-5--5---5---3----4t-4\\-3\\-.}
  {   Jöj-je-tek, * ör-ven-dez-zünk az Úr--nak.}

% ············································································
% Officium lectionis
% ············································································

% Hymnus
\defsheet{PrimoDierumOmnium}
  {<ô-2--2--2---1---2---4---3--2í--:-2--3----4--5------4--5---5--6x--,-6----6----5---6--4----5--rR3-2í-:-2--1----2--4--5----2-1--2í-?,-1e4R3E2W1-2í-.}
  {   El-ső nap min-den nap kö-zött, me-lyen ég s_föld al-kot-ta-tott, s_me-lyen fel-tá-madt Al-ko--tónk le-győz-te ér-tünk a ha-lált, Á---------men.}

% Antiphona 1
\defsheet[tone={1re-a}]{InLege}
  {<-3--1äeE2s0-¨-3---5--4---tT45zZ5-:-5---4---3---3-4--eE2-1-:-3r-2---3--1---1-?,}
  {  Az Úr      * tör-vé-nyé-ben       van tel-jes a-ka-rat-tal éj-jel és nap-pal.}

% Antiphona 2
\defsheet[tone={1la-a}]{Praedicans}
  {<-1t--5z--5--¨-5--6--4---5--5---4---2-:-4--2-3-----4--eE2-1--0w--rf2-eE2-1-?,}
  {  Hir-det-ni * az Úr-nak pa-ran-csa-it  az ő szent he-gyé-re ren-del-te--tett.}

% Antiphona 3
\defsheet[tone={8szo-c}]{TuEsGloria}
  {<-3r-4----¨-5--4---eE2-3--:-1--2----3--4---5--rR3--4t-5--:-7--7-6--7----5---7--4--3-,-3--4---5----7---7i--7--6-:-7-5--rR3---4t-5----4-?,}
  {  Te vagy * di-cső-sé--gem, te vagy ol-tal-ma-zóm, U--ram, te e-mel-ted fel fe-je-met és meg-hall-gat-tál en-gem a te szent he-gyed-ről.}

% ············································································
% Laudes
% ············································································

% Hymnus
\defsheet{AeterneRerumConditor}
  {<-5--5---5--5---5--5--4--4-,-5--5--5----4-3---4--5--5-,--5-5---5-4--3---4---3--2----,-3-------4--4--3--4---5-4--4-?,-4tT4-3\\r\\-.}
  {  Vi-lág te-rem-tő Is-te-ne, ki Úr vagy é-jen és na-pon, i-dőt i-dő-vel vál-to-gatsz, s_nincs mű-ve-id-ben u-na-lom, Á----men.}

% Antiphona 1
\defsheet[tone={7la-d}]{AdTeDeLuce}
  {<-4---6---7--8--8---¨-9---8--7---7----7--7uU6-,-7----8oO8-uU67iI7-5--6uU6-4--4-?,}
  {  Hoz-zád éb-re-dek * vir-ra-dat-kor, Is-ten,   hogy lás--sam     ha-tal--ma-dat.}

% Antiphona 2
\defsheet[tone={8szo-c}]{HymnumDicamus}
  {<-4----4----tT4-3-----¨-5--5èuU6-5tT4-5z-5--4----rR3d1-3--3r-4--4-?,}
  {  Mond-junk him-nuszt * az Úr----nak, Is-te-nünk-nek,  al-le-lu-ja.}

% Antiphona 3
\defsheet[tone={8szo-c}]{FiliiSion}
  {<-7--7--7--uU6-5x-:-5--5----4--5---6--5--4----4----4--5z-4--4-?,}
  {  Si-on fi-a---i  * uj-jong-ja-nak Ki-rá-lyuk-ban, al-le-lu-ja.}

% Responsorium breve
\defsheet{ChristiFiliDeiVivi}
  {<X-¢----------------4--3---4t-tT4-,-5--4---3----rR3-3--?,-5z-4-:-4--4-5---rR3--4t--5-?,--[-------------5z--5-:-5-5--4-4-:-5--5-4-----3--4t--tT4-.}
  {   Krisztus,_az_élő Is-ten Fi-a,  * Ir-gal-mazz ne--künk. Ak-i   az A-tya jobb-ján ülsz. Dicsőség_az_A-tyá-nak a Fi-ú-nak és a Szent-lé-lek-nek.}

% ············································································
% Hora Media
% ············································································

% Hymnus
\defsheet{RectorPotens}
  {<-5-5-----5---5--5---5--4--4\\-,-5-5---5----5---4--5---5---5x-,-5---5---5--4--3---4--3---2\\-,-3--4---4---4--3--5--4---4x-?,-4tT4-3\\r\\-.}
  {  U-runk, fel-sé-ges Is-te-nünk, a lét rend-jét ki meg-sza-bod: reg-gel vi-lá-gos-sá-got adsz, és dél-ben tű-ző ra-gyo-gást, Á----men.}

% Antiphona 1
\defsheet[tone={8szo-b}]{BonumEst}
  {<XB-1----2e-3r-3---:-3--4---3---23r-wW1-0í-,-0----2--3--4---3---34zZ56uU6-:-4--3--2e--3y-?,}
  {    Jobb az Úr-ban * bi-zal-mat tar-ta--ni,  mint fe-je-del-mek-ben         bi-za-kod-ni.}

% Antiphona 2
\defsheet[tone={4-a}]{FortitudoMea}
  {<-0--1e-3-2---1w-2---:-2--4--5---tT4R3-rR3-1w-2-,-2e-4-eE2--wW1-0---1---1--3rR3-2--2-?,}
  {  Az én e-rős-sé-gem * és di-csé-re----tem az Úr  és Ő lett né--kem sza-ba-dí---tá-som.}

% Antiphona 3
\defsheet[tone={8szo-c}]{Confitebor}
  {<-3--3---3r-4---¨-5--7----7i-7--:-7----6---5----6u--5---4--4-?,}
  {  Há-lát a--dok * ne-ked, U--ram, mert meg-hall-gat-tál en-gem.}

% ············································································
% II. Vesperas
% ············································································

% Hymnus
\defsheet{LucisCreator}
  {<ô-2----2--2--1--2---4--3-2í--:-2--3---4---5---4--5-----5--6x-,-6----6--5----6--4--5---rR3-2í-:-2--1---2-4---5---2---1--2í-?,-1e4R3E2W1-2í-.}
  {   Fény al-ko-tó-ja, jó U-runk, ki min-den nap-ra fényt ho-zol, s_az új fény el-ső haj-na--lán  ad-tál a lét-nek kez-de-tet, Á---------men.}

% Antiphona 1
\defsheet[tone={1la-a}]{Virgam}
  {<-3-0--1-1tu-5---:-5---5--tT45uj5T4-,-5--4--3-----3--2--3r-eE2-1-1---,-1-2e--4tT4-3-:-¢----------2--3--4--3--2--1---1-?,}
  {  A te e-rőd-nek * pál-cá-ját         ki-bo-csájt-ja az Úr Si--on-ból: u-ral-kod--jál a_te_ellen-sé-ge-id kö-ze-pet-te.}

% Antiphona 2
\defsheet[tone={4-a}]{ExAegypto}
  {<-6-tT4--5z--5---¨-6--7--6---5---4---5z---6-5-?,}
  {  E-gyip-tom-ból * ki-ve-zet-tél min-ket, U-runk.}

% Antiphona 3
\defsheet{SalusEtGloria}
  {<----7--7--8oO8-7i-:-[---------------------------4--5----rR3-:-----4--4t-3--3-,------[--------------------------4--3--4t-5-:-----7--7--iI7-5-:-5--4tT4-3--3-?,-----7--7--8oO8-7i-:-[----------------------------------4----5--rR3-:---4--4t-3--3-,------[-----------------------------------4----3--4t-5-:--------7--7--iI7-5-:-5--4tT4-3--3-?,-----7--7--8oO8-7i-:-[----------------------5---5---5--5--4--5--rR3-:------4--4t-3--3-,------[------------------------------4---3--4t--5-:------7--7--iI7-5-:-5--4tT4-3--3-?,-----7--7--8oO8-7i-:-[----------------------------4--5---4---rR3-:----4--4t-3--3-,------[------------------4--3-----4t--5-:-------7--7--iI7-5-:-5--4tT4-3--3-?,-----7--7--8oO8-7i-:-[----------------------4--5-rR3-:-----4--4t-3--3-,------5--5-----5--4---3--4t--5-:-------7--7--iI7-5-:-5--4tT4-3--3-?,-----7--7--8oO8-7i-:-[--------------------------------4---5---rR3-:-----4--4t-3--3-,------[----------------4---3---4t-5-:-------7--7--iI7-5-:-5--4tT4-3--3-.}
  {<V.> Al-le-lu---ja.  Üdv,_dicsőség_és_hatalom_Is-te-nünk-nek, <F.> al-le-lu-ja, <V.> mert_igazak_és_jóságosak_í-té-le-te-i. <F.> Al-le-lu--ja, al-le---lu-ja. <V.> Al-le-lu---ja.  Dicsérjétek_Istenünket,_ti,_összes szol-gá-i, <F.> al-le-lu-ja, <V.> és_akik_félve_tisztelitek_őt,_kicsi-nyek és na-gyok. <F.> Al-le-lu--ja, al-le---lu-ja. <V.> Al-le-lu---ja.  Átvette_uralmát_Urunk, min-den-ha-tó Is-te-nünk, <F.> al-le-lu-ja, <V.> örvendezzünk,_ujjongjunk_és_di-cső-ít-sük őt! <F.> Al-le-lu--ja, al-le---lu-ja. <V.> Al-le-lu---ja.  Elérkezett_a_Bárány_menyegző-jé-nek nap-ja, <F.> al-le-lu-ja, <V.> már_fel_is_készült me-nyasz-szo-nya. <F.> Al-le-lu--ja, al-le---lu-ja. <V.> Al-le-lu---ja.  Dicsőség_az_Atyának_és Fi-ú-nak, <F.> al-le-lu-ja, <V.> és Szent-lé-lek Is-ten-nek. <F.> Al-le-lu--ja, al-le---lu-ja. <V.> Al-le-lu---ja.  Miképpen_kezdetben_vala,_most_és min-den-kor, <F.> al-le-lu-ja, <V.> és_mindörökkön_ö-rök-ké, á--men. <F.> Al-le-lu--ja, al-le---lu-ja.}

% Responsorium breve
\defsheet{BenedictusEs}
  {<X-3--4----3-----4t-5-----,-5--4--3----rR3-3--?,-5--5---4---5--4---4-:-5--5--5----4---3-4t--5-?,-[-------------5z--5-:-5-5--4-4-:-5--5-4-----3--4t--tT4-.}
  {   Ál-dott vagy, U--runk, * Az ég-bolt fö--lött. Di-csé-ret-re mél-tó  és ma-gasz-tos ö-rök-ké.  Dicsőség_az_A-tyá-nak a Fi-ú-nak és a Szent-lé-lek-nek.}

% ============================================================================
% FERIA II
% ============================================================================

% ············································································
% Invitatorium
% ············································································

\defsheet{IubilemusDeo}
  {<X-¢---------4--3--4t---5---,-5--5---3-----4t-4--3-.}
  {   Szabadító Is-te-nünk-nek * vi-gad-junk, az Úr-nak!}

% ············································································
% Officium lectionis
% ············································································

% Hymnus
\defsheet{SomnoRefectisArtubus}
  {<-3--3---4---5----4----4---3-2--:-4---2----3--1-0------1---3--3---,-3-3------4---5---4----4-3--2---:-4-2----3--1---0----3----1--1--?,-0w32W1Q0-1í-.}
  {  Al-vás-ban fris-sült már a test ott-hagy-va á-gyunk, fel-ke-lünk. A-tyánk, hoz-zád szól é-ne-künk; e-seng-ve kér-jük: légy ve-lünk. Á--------men.}

% ············································································
% Laudes
% ············································································

% Hymnus
\defsheet{SplendorPaternaeGloriae}
  {<-3--3---4---5----4----4---3-2--:-4---2----3--1-0------1---3--3---,-3-3------4---5---4----4-3--2---:-4-2----3--1---0----3----1--1--?,-0w32W1Q0-1í-.}
  {  A-tyá-nak fé-nye, fény-su-gár, nap-ból su-gár-zó nap-vi-lág, fény fé-nye, nap-pal nap-ja vagy: for-rás, hol min-den fény fa-kad. Á--------men.}

% Responsorium breve
\defsheet{BenedictusDominusASaeculo}
  {<X-4--3----4t-tT4-,-4-5---rR3-4-3---3-?,-5z-4-:-4-4---4---4--5---rR3-4t-5-?,-[-------------5z--5-:-5-5--4-4-:-5--5-4-----3--4t--tT4-.}
  {   Ál-dott az Úr, * Ö-rök-kön-ö-rök-ké.  A--ki  e-gye-dül mű-vel cso-dá-kat. Dicsőség_az_A-tyá-nak a Fi-ú-nak és a Szent-lé-lek-nek.}

% Benedictus
\defsheet[tone={6-a}]{BenedictusDeus}
  {<-34tg3r-3rR3d1-¨-3--4--5--4--3--3-?,}
  {  Ál-----dott   * Iz-ra-el Is-te-ne.}

% ············································································
% Hora Media
% ············································································

% Antiphona 1
\defsheet[tone={4-a}]{PraeceptumDomini}
  {<-3--1--2---3r-2---¨-eE2--1-:-3---4--5--4--4t--4--3-eE2-1w-2-?,}
  {  Az Úr tör-vé-nye * fény-lő, meg-vi-lá-go-sít-ja a sze-me-ket.}

% Antiphona 2
\defsheet[tone={8szo-c}]{DomineDeusMeusInTe}
  {<-4-3----5--7i-7--7----:-5--7---7---4y-4y-?,}
  {  U-ram, én Is-te-nem, * te-ben-ned bí-zom.}

% Antiphona 3
\defsheet[tone={8szo-c}]{DeusIudexIustus}
  {<-4--4---5z-5--4t-rR3-:-4-5---7--6----5----7--4-,-7--7---6--7---5----4t-:-2e--4t--4-?,}
  {  Is-ten i--gaz bí-ró * e-rős és hosz-szan tű-rő: va-jon ha-rag-szik-e    min-den nap.}

% ============================================================================
% FERIA III
% ============================================================================

% ············································································
% Invitatorium
% ············································································

\defsheet{RegemMagnum}
  {<-1-3----2e-1------,-3-3---4---5--rR3E2-2-.}
  {A nagy Ki-rályt, * i-mád-juk az U-----rat!}

% ············································································
% Officium Lectionis
% ············································································

% Hymnus
\defsheet{ConsorsPaterniLuminis}
  {<ô-2--2--2---1----2-4----3---2í-:-2--3----4---5--4----5---5--6x-,-6--6--5---6---4-----5-rR3-2í--:-2----1--2--4-----5---2---1--2í-?,-1e4R3E2W1-2í-.}
  {   Vi-lá-gos-ság, A-tyád-dal egy, te fény-nek fé-nye, nap-vi-lág, az éj-ből fel-száll é-ne--künk: jöjj el kö-zénk, így kér-le-lünk. Á---------men.}

% ············································································
% Laudes
% ············································································

% Hymnus
\defsheet{PergrataMundoNuntiat}
  {<--5---5---5---5---5---5z---5---5-:--5---5---5---4---3---4t---4---4-;-3---5---5---4---3---4---4---5-:---4---4---4---3---1---3---2---1--?,-1wW1-0\\q\\-.}
  {   Várt haj-nal ad már hírt ne-künk, a nap dár-dá-ja éjt sza-laszt, min-dent a-rany-ba önt a fény, s_a föld szí-nes ru-há-kat ölt.        Á----men.}

% Responsorium breve
\defsheet{DeusMeus}
  {<X-3--4--3--:-3--4t-tT4--,-4-5--rR3-4--3--3-?,-5--5---5z-4-:-5--rR3-4t--tT4-?,-[-------------5z--5-:-5-5--4-4-:-5--5-4-----3--4t--tT4-.}
  {   Is-te-nem, se-gí-tőm, * A-ki-ben re-mé-lek. Ol-tal-ma-zóm és meg-men-tőm.   Dicsőség_az_A-tyá-nak a Fi-ú-nak és a Szent-lé-lek-nek.}

% Benedictus
\defsheet[tone={7la-d}]{Erexit}
  {<-4---4-6---7--8---7---6--7i-8---:-9--ö--9---8---9-8--8--,-8--8-9----8--uU6-7--:-5--5---6---5--4--4-?,}
  {  Föl-e-mel-te né-künk az Is-ten * az üd-vös-ség e-re-jét, az ő szol-gá-já--nak, Dá-vid-nak há-zá-ban.}

% ············································································
% Hora Media
% ············································································

% Antiphona 1
\defsheet[tone={2-b}]{Beati}
  {<X-6---5z-4----:-£-----------------3----4z--4y-?,}
  {   Bol-do-gok, * kik_törvényed_sze-rint jár-nak.}

% Antiphona 2
\defsheet[tone={2-f}]{CantaboDomino}
  {<-1-3--2--3---2--1--1---:-3--1--qQ0-1-eE2--1--1-?,}
  {  É-ne-ke-lek az Úr-nak * ki jó-kat a-dott ne-kem.}

% Antiphona 3
\defsheet[tone={8szo-c}]{LaeteturIacob}
  {<-5-4---3---4t-4----:-5--7--uU6--tT4R3-5z-5--4-?,}
  {  Ö-rül-jön Já-kob, * és uj-jong-jon   Iz-ra-el.}

% ============================================================================
% FERIA IV
% ============================================================================

% ············································································
% Invitatorium
% ············································································

\defsheet{DominumQuiFecitNos}
  {<-3--3r-3---:-3-3--4--3--4t--5----,-5---3--4t---4-4---3-.}
  {  Az Is-tent, a mi al-ko-tón-kat, * jöj-je-tek, i-mád-juk!}

% ············································································
% Officium Lectionis
% ············································································

% Hymnus
\defsheet{RerumCreatorOptime}
  {<ô-4--4---4--4---4--4--3-2í--:-4----4-----4--4---1--2--1--0í---,-2-2---3---4--3---2---1---2eE2í-:-2---2---1---0---1--2--1--1í-?,--1wW1-0q-.}
  {   Vi-lág-te-rem-tő jó U-runk, nézz ránk, ve-zér-lő tá-ma-szunk! A-kik-nek al-vás eny-het ad,     bűn-ben tes-ped-ni so-se hagyd. Á----men.}

  % Antiphona 1
\defsheet[tone={6-a}]{DiligamTeDomine}
  {<-3---4---5---3r-3---eE2-1---:-3--4-4---4--3-?,}
  {  Sze-ret-lek té-ged U---ram * én e-rős-sé-gem.}

% Antiphona 2
\defsheet[tone={1re-a}]{FactusEst}
  {<-3--4tT4-5x--:-5--4---2e--4--eE2--1--0q-1í-?,}
  {  Se-gí---tőm * te let-tél ne-kem, Is-te-nem.}

% Antiphona 3
\defsheet[tone={4-a}]{RetribuitMihiDominus}
  {<-4---3--2---3--qQ0-1e-2--:-2-4---4t--rR3-2---2-?,}
  {  Meg-fi-zet ne-kem az Úr * i-gaz-vol-tom sze-rint.}

% ············································································
% Laudes
% ············································································

% Hymnus
\defsheet{NoxEtTenebrae}
  {<-3--3--4--5----4--4---3--2--:-4--2--3----1--0---1---3--3---,-3--3---4--5---4---4---3-2---:-4-----2---3---1----0--3---1--1-?,-02eE2W1Q0-1-.}
  {  Éj és ho-mály és fel-le-gek, ör-vé-nyek és köd-tor-la-szok, fe-hér az ég, már kél a fény: Krisz-tus van itt, tá-voz-za-tok! Á---------men.}

% Responsorium breve
\defsheet{InclinaCorMeumDeus}
  {<X-¢------------------4--3--4t--5------:-4--3--3-?,-5-5--5-5z-4-:-4--3----4t-5--?,--[-------------5z--5-:-5-5--4-4-:-5--5-4-----3--4t--tT4-.}
  {   Szívemet_tanításod fe-lé for-dítsd, * Is-te-nem. A te u-ta-don él-tess en-gem!   Dicsőség_az_A-tyá-nak a Fi-ú-nak és a Szent-lé-lek-nek.}

% ············································································
% Hora Media
% ············································································

% Antiphona 1
\defsheet[tone={1la-a}]{InTotoCordeMeo}
  {<X-1---0---3r--3---5---:-6--5---4---4--4----3t-4--,-4--4-----5--4---3-:-2--3---4---eE2-1-?,}
  {   Tel-jes szí-vem-ből * ke-res-lek té-ged, U--ram, ne hagyj el-tér-ni  pa-ran-csa-id--tól.}

% Antiphona 2
\defsheet[tone={7do-d}]{InclinaDomine}
  {<-8---9-----7i-8----:-8-6----7--6--4-,-6--6----6u---5---6---5--4-?,}
  {  For-dítsd ho-zám, * U-ram, fü-le-det és hall-gasd meg sza-va-mat.}

% Antiphona 3
\defsheet[tone={5-c}]{EgoAutem}
  {<-tT4-tT4-3---:-3-5---77i-7---7--6--7i--5-:-uU6-5---5-4-----tT4-3-?,}
  {  Én  pe--dig * i-gaz-ság-ban je-le-nek meg szí-ned e-lőtt, U---ram.}

% ============================================================================
% FERIA V
% ============================================================================

% ············································································
% Invitatorium
% ············································································

\defsheet{AdoremusDominum}
  {<-3-3---4---3--4t-5-----,-5-3--4--5----4---3-.}
  {  I-mád-juk az Is-tent, * ő al-ko-tott min-ket!}

% ············································································
% Officium Lectionis
% ············································································

% Hymnus
\defsheet{NoxAtraRerum}
  {<-5-5---5---5---5---5----4--4-,-5--5---5--4-----3--4---5--5-,-5-5---5--4-----3---4----3--2-,-3--4---4-3----4--5---4--4-?,-4tT4-3\\r\\-.}
  {  A föl-dön min-den szín-re még sö-tét fá-tyolt bo-rít az éj. I-gaz Bí-ránk, szí-vünk fe-léd ál-dón e-seng és kér-ve kér. Á----men.}

% ············································································
% Laudes
% ············································································

% Hymnus
\defsheet{SolEcceSurgit}
  {<ô-2---2---2-1----2-4---3--2í-:-2--3----4--5----4--5---5--6x-,-6--6---5--6---4--5--rR3-2í-:-2--1--2--4---5---2----1--2í-?,-1e4R3E2W1-2í-.}
  {   Már kél a nap, a csu-pa tűz, és ránk pi-rít, bá-nat-ra int; ki tud-na bűn-be es-ni  már, ha ez fé-nyé-vel ránk te-kint! Á---------men.}

% Responsorium breve
\defsheet{ClamaviInTotoCorderMeo}
  {<X-3---3---3---4---3---4--4t-tT4--:-5----4----3----rR3-3--,-5--5---5---5z-4-:-5--rR3-4t--tT4--?,--[-------------5z--5-:-5-5--4-4-:-5--5-4-----3--4t--tT4-.}
  {   Tel-jes szí-vem-ből ki-ál-tok, * Hall-gass meg, U---ram. Pa-ran-csa-i--dat én meg-tar-tom. Dicsőség_az_A-tyá-nak a Fi-ú-nak és a Szent-lé-lek-nek.}

% Benedictus
\defsheet[tone={7la-a}]{InSanctitate}
  {<ô-1----1---3----4--5z-5---:-77----zZ5z-5x-,-5--5---5---6--5---rR3-4-:-2--2---3--2--1í--1í-?,}
  {   Szol-gál-junk az Úr-nak * szent-ség--ben, és meg-sza-ba-dít min-ket el-len-sé-ge-ink-től.}

% ============================================================================
% FERIA VI
% ============================================================================

% ············································································
% Invitatorium
% ············································································

\defsheet{DominumDeumNostrum}
  {<-3--3r-3-:--3-3--4--3--4t--5----,-5---3--4t---4-4---3-.}
  {  Az U--rat, a mi Is-te-nün-ket, * jöj-je-tek, i-mád-juk!}

% ············································································
% Officium Lectionis
% ············································································

% Hymnus
\defsheet{TuTrinitasUnitas}
  {<ô-4-----4--4---4---4---4--3--2í--:-4--4---4---4--1-2----1--0í--,-2-----2---3--4---3--2-1--2eE2í-:-2--2---1--0--1----2----1---1í-?,-1wW1-0\\q\\-.}
  {   Szent-há-rom-ság egy Is-te-nünk, ha-tal-mas Úr a föld fe-lett, halld meg di-csé-rő é-ne-künk,   éb-red-ve me-lyet zeng szí-vünk. Á----men.}

% ············································································
% Laudes
% ············································································

% Hymnus
\defsheet{AeternaCaeliGloria}
  {<-3-3---4---5---4--4---3--2--:-4--2---3--1---0---1---3--3---,-3-3---4--5---4---4---3--2-:-4-2----3--1----0---3---1--1-?,-02eE2W1Q0-1-.}
  {  Ö-rök tün-dök-lő Nap-vi-lág, ha-lan-dó nép-nek egy re-mény, a Fön-sé-ges-nek egy Fi-a,  a tisz-ta Szűz-nek mag-za-ta:  Á---------men.}

% Responsorium breve
\defsheet{AuditamFacMihiMane}
  {<X-¢--------------------------4--3----4t--5--:--4--3--rR3-3-?,-5--5-5----4---5--4---4--:-4--3----4t-?,--[-------------5z--5-:-5-5--4-4-:-5--5-4-----3--4t--tT4-.}
  {   Add,_hogy_már_korán_reggel ta-pasz-tal-jam * Jó-sá-go--dat! Az u-tat, mer-re men-jek, mu-tasd meg.   Dicsőség_az_A-tyá-nak a Fi-ú-nak és a Szent-lé-lek-nek.}

% Benedictus
\defsheet[tone={5-c}]{VisitavitEtFecit}
  {<-3---5--7--7---uU6-¨-tT4-7---8---5---4--tT4-3--4--4t-3--3-?,}
  {  Meg-lá-to-gat-ta  * és  meg-vál-tot-ta az  Úr az ő  né-pét.}

% ============================================================================
% SABBATUM
% ============================================================================

% ············································································
% Invitatorium
% ············································································

\defsheet{PopulusDomini}
  {<-1--2--4t-5-:-5--4--5--6--5--4---5--5--5--,-6u--6--5----zZ5T4-5z--6---5-.}
  {  Az Úr né-pe, és le-ge-lő-jé-nek ju-ha-i, * jöj-je-tek, i-----mád-juk őt.}


% ============================================================================
% DOMINICA
% ============================================================================

% ············································································
% Vesperas I.
% ············································································

% Antiphona 3
\defsheet[tone={1la-a}]{DeditIlliDominus}
  {<X-1-1---1--0---34t-5---:-6-5----4--4--3t-4-,-4--2-4---tT4-3-:-3--3-----2--3--4-eE2-1--1-?,}
  {   Ö-rök di-cső-sé--get * a-dott ne-ki az Úr, és ö-rök ne--vet ha-gyott ne-ki ö-rök-sé-gül.}

% ············································································
% Officium lectionis
% ············································································

% Hymnus
\defsheet{MediaeNoctis}
  {<-1----0--1---3----4--3--2---1-:-4---4---3-4---5--3---2---1-,-1--1----2---0--ð----0----1--1--:-1-0----1----3--4-3----2--1-?,-1wW1-0q-.}
  {  Csak fé-lig múlt az éj-sza-ka, szó-lít a pró-fé-ták sza-va: di-csér-jük Is-tent szün-te-len, A-tyát s_Fi-út e-gyüt-te-sen, Á----men.}

% vel:
% TODO

% Antiphona 1
\defsheet[tone={3-c}]{Confessionem}
  {<-4---4ç6u-6--:-5--7---6--5---5--6--5--4--,-4-----5z-6----5--4---tT4-:-2e---rR3-2-?,}
  {  Fön-ség--be * és mél-tó-ság-ba öl-tö-zöl, fényt öl-tesz ma-gad-ra,   mint ru-hát.}

% Antiphona 2
\defsheet[tone={8szo-c}]{DomineDeusMeus}
  {<-rR3E2-1----3r-5--4----:-4--5----4---3----5u-6--4-?,}
  {  U-----ram, Is-te-nem, * ma-gasz-tos vagy i--ga-zán.}

% Antiphona 3
\defsheet[tone={tirr-e}]{QuamMagnificataSuntE}
  {<-2----2---1--2--1---:-1-0--1w-3--1---3-2-?,}
  {  Mily fön-sé-ge-sek * a te mű-ve-id, U-ram.}

% ············································································
% Laudes
% ············································································

% Hymnus
\defsheet{EcceIamNoctis}
  {<X-5--4----5--4---1í-:-4--4--3---4--5x-5x-,-5---7---8-7---5x-:-5--4---3---4--5x---5x-,-5---4---5----4---1í-:-4-4--3---4-tT4-3y-,-3---3---2-1í---1í-?,-1wW1-0\\q\\-.}
  {   Sá-pad, el-osz-lik  kö-de már az éj-nek, pir-kad a haj-nal, ró-zsa-szí-nű fény kél; for-dul-junk buz-gón  a vi-lág U-rá--hoz  han-gos i-mánk-kal,  Á----men.}

% Antiphona 1 (Nb41)
\defsheet[tone={8szo-c}]{DominusMihiAdiutor}
  {<-2--3---4--4--:-5--4--3--4t--eE2-1í-,-3-4----5---4--45uU67iI7-:-5---5---4---3--4---5--4x-4x-?,}
  {  Ve-lem az Úr * az én se-gít-sé--gem, a-zért nem fé-lek,        mit árt-hat né-kem az em-ber?}

% Antiphona 2
\defsheet[tone={8szo-c}]{TresExUno}
  {<-4---5---¨-4---5u---uU6Z5-:-iI7-iol7i-7---7iI77U6Z5-3-5--5---5--uU6--tT4-5-4-,-5--5uU6-tT4-3---3--:-4--5----6--5--4-?,}
  {  Hár-man * egy száj-jal     ki-ál-----tot-tak       a ke-men-ce láng-ja--i-ban és é----ne--kel-ték: ál-dott az Is-ten.}

% Antiphona 3
\defsheet[tone={8szo-c}]{LaudateDominum}
  {<-4--4----4--4---3t-7-6---¨-4----6--uj5-rR3-:-5--5--7--5--zZ5-4-?,}
  {  Di-csér-jé-tek az U-rat * nagy-sá-gá--nak   so-ka-sá-ga sze-rint.}

% Responsorium breve
\defsheet{ConfitebimurTibiDeus}
  {<X-3--3----3--3----4--3--:-4--4t-5-----,-4--4--5--rR3-4---3---3-?,-5--5--5z---4-:-5---5--4---3--4t-5-?,-5--5---5-.}
  {   Ma-gasz-ta-lunk té-ged, Is-te-nünk, * És ne-ve-det szó-lít-juk. El-be-szél-jük cso-da-tet-te-i--det. Di-cső-ség...}

% ············································································
% Vesperas II.
% ············································································

% Antiphona 1 (SzKD43)
\defsheet[tone={8szo-c}]{IuravitDominus}
  {<-4---6--8--8----7i-7---:-7--6---4---zZ5-4-,-44--eE2--1--3----4t-4x--4x-?,}
  {  Meg-es-kü-dött az Úr, * és meg nem bán-ja: pap vagy te mind-ö--rök-ké.}

% ============================================================================
% FERIA II
% ============================================================================

% ············································································
% Officium lectionis
% ············································································

% Hymnus
\defsheet{IpsumNuncNobis}
  {<-1----0---1--3-4---3-2--1-:-4--4-3---4--5--3--2---1--,-1---1--2--0---ð-0--1--1--:-1--0-----1--3---4---3-2--1-?,-1wW1-0q-.}
  {  Most itt az ó-ra, a-mi-kor az e-van-gé-li-um sze-rint meg-ér-ke-zik a Vő-le-gény és menny-or-szá-got a-la-pít. Á----men.}

% Antiphona 1
\defsheet[tone={8szo-c}]{InTuaIustitia}
  {<-7-7--7--7---5--4---45u-7---:-7---7--7----7---5--4----5-4-?,}
  {  A te ig-gaz-sá-god ál--tal * sza-ba-díts meg en-gem, U-ram!}

% Antiphona 2
\defsheet[tone={8szo-c}]{FaciemTuam}
  {<-5-3--5--7--6----4-5----:-7--7----7----7---6----7---5--4-?,}
  {  A te ar-co-dat, U-ram, * ra-gyog-tasd fel szol-gád fö-lött!}

% Antiphona 3
\defsheet[tone={2-f}]{DiligiteDominum}
  {<-1---1---1--0---1--eE2-1----:-3---3---3-----4--5-:-3--3---3--2-3--1---0---0q--3----3r-eE2-1-?,}
  {  Sze-res-sé-tek az U---rat, * min-den szent-je-i,  mi-vel az i-ga-za--kat meg-tart-ja az  Úr.}

% ············································································
% Laudes
% ············································································

% Hymnus
\defsheet{LucisLargitorSplendide}
  {<ô-4--4----4---4---4---4--3--2í-:-4--4-----4-4-----1--2---1--0í-,-2--2--3---4---3----2---1--2eE2í-:-2---2---1--0--1--2---1--1í-?,-1wW1-0q-.}
  {   Fé-nyes-ség-osz-tó, bő-ke-zű,  ki fényt á-raszt-va gaz-da-gon  az éj-nek gyá-szát osz-la-tod,    s_a nap vi-lá-ga tűz le ránk. Á----men.}

% Antiphona 1
\defsheet[tone={8szo-c}]{Sitivit}
  {<-4u---6--4---¨-6-7--5---4-tT4-3--:-5z-5-4-?,}
  {  Szom-ja-zik * u-tá-nad a lel-kem, én U-ram!}

% Antiphona 2
\defsheet[tone={3-c}]{Ostende}
  {<-4--5z---6---¨-7--tT4---5z-tT4-:-4çz-5---4---4--3----2--2-?,}
  {  Mu-tasd meg * ne-künk, U--ram,  ir--gal-mad fé-nyes-sé-gét!}

% Antiphona 3
\defsheet[tone={2-f}]{Opera}
  {<-3--eE2-1---¨-33r-3---3-:-3tT4t-4--eE2-0-äe----1-?,}
  {  Ke-zé--nek * mű--ve--it  hir---de-ti  a menny-bolt.}

% Responsorium breve
\defsheet{ExultateIustiInDomino}
  {<X-¢----------------4----3--4t-tT4--,-[---------------4---3---rR3-3-?,-5-5z--4-:-5--5----5-4--3---4t-5-?,-5--5---5-.}
  {   Örvendjetek,_iga-zak, az Úr-ban, * A_szentekhez_di-csé-ret il--lik. E-lőt-te  új dalt é-ne-kel-je-tek. Di-cső-ség...}

% Benedictus
\defsheet[tone={4-a}]{BenedictusDominusQuia}
  {<-2--1----wW1-0---:-2----ed1-2--3--4t-5--:-tT4-3---4t--4----2---2-?,}
  {  Ál-dott az  Úr, * mert meg-lá-to-ga-tott és  meg-vál-tott min-ket.}

% ············································································
% Hora Media
% ············································································

% Hymnus
\defsheet{DicamusLaudesDomino}
  {<-5--5----5---5---5---5--4-4-,--5---5--5----4---3---4---5--5--,-5--5-5---4-3---4---3--2-,-3-----4---4---3-4---5---4--4-?,-4tT4-3\\r\\-.}
  {  Di-csér-jük dal-lal az U-rat, buz-gó kész-ség-gel lel-ke-sen: az ó-ra, í-me, dél-re jár s_ben-nün-ket i-mád-ság-ra hív. Á----men.}
  
% ============================================================================
% FERIA III
% ============================================================================

% ············································································
% Officium lectionis
% ············································································

% Hymnus
\defsheet{NocteSurgentes}
  {<X-5---4----5---4--1í-:-4----4--3---4---5x--5x---,-5---7--8-7----5x-:-5---4--3--4---5x--5x-,-5----4---5---4-----1í-:-4-4--3---4---tT4--3y-,-3-3---2-1í---1í-?,-1wW1-0\\q\\-.}
  {   Kel-jünk fel éj-jel  s_kö-zö-sen vir-ras-szunk, egy-re a zsol-tár  sza-va-it fon-tol-va,  szár-nya-lón száll-jon  a di-csé-ret hang-ja   é-des U-runk-hoz.  Á----men.}

% Antiphona 1
\defsheet[tone={6-a}]{Revela}
  {<-3----4---5--4--3---:-4-3--3---?,}
  {  Tárd fel az Úr-nak * u-ta-dat.}

% Antiphona 2
\defsheet[tone={2-f}]{IuniorFui}
  {<-eE2-1--2e--2---¨-1--1---1--wW1-1--:-0q-3---3r--3---3-4--3---wW1-eE2--1-?,}
  {  If--jú vol-tam * és meg-vé-nül-tem, de nem lát-tam i-ga-zat el--hagy-va.}

% Antiphona 3
\defsheet[tone={6-a}]{SperaInDomino}
  {<-3---3---2--4--5---:-5--5-4---5---4--3-?,}
  {  Bíz-zál az Úr-ban * és a jót cse-le-kedd.}

% ············································································
% Laudes
% ············································································

% Hymnus
\defsheet{AeternaeLucisConditor}
  {<ô-2----2--2--1--2---4--3--2í--:-2-3---4--5--4-----5---5--6x-,-6--6---5---6----4-----5--rR3-2í--:-2---1--2---4---5---2---1--2í-?,-1e4R3E2W1-2í-.}
  {   Fény al-ko-tó-ja, Is-te-nünk, ö-rök ve-rő-fény, nap-vi-lág, va-lód-ból csak fényt is-me--rünk, nem ér fel hoz-zád vak ho-mály. Á---------men.}

% Antiphona 1
\defsheet[tone={6-a}]{Salutare}
  {<-3--4t-3---:-3--4---eE2-1---3--4--4--3-?,}
  {  Én or-cám * üd-vös-sé--ge, én Is-te-nem.}

% Antiphona 2
\defsheet[tone={3-c}]{CunctisDiebus}
  {<-4-4--5u---7---¨-5---7---zZ5-4-:-4-----3---4---5----rR3E2-2-?,}
  {  É-le-tünk-nek * min-den nap-ján ments meg min-ket, U-----runk!}

% Antiphona 3
\defsheet[tone={2-f}]{TeDecetHymnus}
  {<-3--3---3--2----3--1----:-1-1--1---1---0--1e-1-?,}
  {  Té-ged il-let, Is-ten, * a di-csé-ret Si-on-ban!}

% Responsorium breve
\defsheet{VocemMeamAudiDomine}
  {<X-¢----------------4---3---4t-5----,-4----4--4----3--4--5---4-ed1-3r-3-?,-[------------4--5--4--4-:-5---4---3---4t-5--?,-5--5---5-.}
  {   Tekints_könyörgé-sem-re, U--ram, * Mert na-gyon re-mé-lek i-gé--id-ben. Téged_szólít ki-ál-tá-som már haj-nal e--lőtt. Di-cső-ség...}

% Benedictus
\defsheet[tone={1re-a}]{DeManuOmnium}
  {<-3----5-4---5--3r-4----:-1-1---3---3---2---2e--3-:-4---4--4----2---eE2-1----0q-1-?,}
  {  Mind-a-zok ke-zé-ből, * a-kik gyű-löl-nek min-ket sza-ba-díts meg min-ket, U--runk.}

% ============================================================================
% FERIA IV
% ============================================================================

% ············································································
% Officium lectionis
% ············································································

% Hymnus
\defsheet{OSatorRerum}
  {<X-5---4--5---4---1í-:-4--4---3-4--5x-5x-,-5----7--8-----7-----5x-:-5--4---3--4---5x-5x-,-5---4--5--4---1í--:-4-4--3---4--tT4-3y-,--3----3-2---1í---1í-?,-1wW1-0\\q\\-.}
  {   Min-de-nek Aty-ja,  vi-lág ú-jí-tó-ja,  nagy Ki-rály, Krisz-tus, fé-lel-me-tes Bí-ró,  áld-va di-csé-rünk, e-se-dez-ve ké--rünk, nézz a szí-vünk-re!   Á----men.}

% Antiphona 1
\defsheet[tone={6-a}]{UtNonDelinquam}
  {<-3--4---5---3---4----eE2-1----:-3---4----4-3-?,}
  {  Ne vét-kez-zem nyel-vem-mel, * add meg, U-ram!}

% Antiphona 2
\defsheet[tone={7do-a}]{NonSisMihi}
  {<ô-5--3----5z-5---:-3--4t--wW1-1í-,-4--4--4--2r-3y-:-3-2---1---2--4---3---1í--1í-?,}
  {   Ne légy né-kem * fé-lel-mem-re,  én re-mé-sé-gem  a szo-ron-ga-tás-nak nap-ján.}

% Antiphona 3
\defsheet[tone={4-a}]{Expectabo}
  {<-3--1----2--3---¨-4--3---2---1e-2--,-eE2W1-23r-eE2-?,}
  {  Re-mény-ke-dem * ne-ved-ben U--ram, mert  jó  az.}

% ············································································
% Laudes
% ············································································

% Hymnus
\defsheet{FulgentisAuctor}
  {<-5--5--5--5--1---3--2--1í--:-5--5--6---4----5-----7----6--5x-,-5-----6---7--8----7----5---7--zZ5T4x-:-5---5--6---4--2---4---5--5x-?,-5zZ5-4\\t\\-.}
  {  Ég Al-ko-tó-ja, Is-te-nünk, ki éj-jel hold-fényt adsz ne-künk s_nap-pal na-pot, mely kör-be fut,     s_e-lő-írt út-ján cél-ba jut.  Á----men.}

% Antiphona 1
\defsheet[tone={8szo-c}]{DeusInSanctoViaTua}
  {<-4u-5----¨-4-----5-rR3-4t-4-,-6u-8-7----5----7--6--:-4----6u-5--4--4--4-?,}
  {  Is-ten, * szent a te  U--tad ki o-lyan nagy is-ten, mint a  mi Is-te-nünk?}

% Antiphona 2
\defsheet[tone={8szo-c}]{Exsultavit}
  {<-7--7---6u-5---:-5--5--5---4---3--4t-4-?,}
  {  Ör-ven-de-zik * az én szí-vem az Úr-ban.}

% Antiphona 3
\defsheet[tone={3-c}]{DominusRegnavit}
  {<-4--5u-uU6-tT4-5---¨-4--3---4t--4---rR3-2-?,}
  {  Az Úr or--szá-gol * ör-ven-dez-zen a --föld!}

% Responsorium breve
\defsheet{BenedicamDominum}
  {<X-3--4---3--4t-tT4-,-5---rR3-4-3--3-?,-5--5---4--4-:-5----5--5--5---4---3--4t-tT4-?,-5--5---5-.}
  {   Ál-dom az U--rat * Min-den i-dő-ben. Di-csé-re-tét szün-te-le-nül zen-gi aj-kam.   Di-cső-ség...}

% Benedictus
\defsheet[tone={8szo-c}]{Liberati}
  {<-4---4---5--5z--zZ5-:-5----7---uU6--5--zZ5-4t---7--6--5----5-:-5z----5---4-?,}
  {  Meg-sza-ba-dul-ván * szol-gál-junk az Úr--nak, Is-te-nünk-nek szent-ség-ben.}

% ============================================================================
% FERIA V
% ============================================================================

% ············································································
% Officium lectionis
% ············································································

% Hymnus
\defsheet{AlesDieiNuntius}
  {<-3-3---4---5----4----4---3--2-:-4--2---3--1----0----1--3-3---,-3-3---4---5--4----4--3--2-:-4-----2----3---1---0---3--1--1-?,-0\\2eE2W1Q0-1\\-.}
  {  A nap-pal szár-nyas hír-nö-ke  je-len-ti már, hogy jő a fény; a lel-kek éb-resz-tő-je is, Krisz-tus, ben-nün-ket él-ni hív: Á-----------men.}

% Antiphona 1
\defsheet[tone={8do-c}]{Extende}
  {<-5---3r----4---:-4t-4-----rR3-4---4t-3--,-5--¦-----------------6--7--5--:-3--4--5z---5---4-?,}
  {  Ter-jeszd ki, * U--runk, ha--tal-ma-dat, és szabadíts_meg_lel-ke-in-ket, el ne pusz-tul-junk.}

% Antiphona 2
\defsheet[tone={2-c}]{Confundantur}
  {<-4---5---7---7--7i--8----:-7-8---9--iI7-8--8---8--,-8--8--5--8---7---6---tT4-5--:-4--3-4----5--7--4--4--?,}
  {  Szé-gye-nül-je-nek meg, * a-kik en-gem ül-döz-nek, és én ne szé-gye-nül-jek meg, én U-ram, én Is-te-nem.}

% Antiphona 3
\defsheet[tone={2-f}]{MementoDomine}
  {<-1wW1Q0-1---eE2-3rR3-2eE2W1-:-3--3--4----5---4---3---4--3r-4-:-3-1----qQ0-3-eE2-1-?,}
  {  Em-----lék-ezz U----ram,   * és mu-tasd meg nek-ünk ma-ga-dat a szük-ség i-de--jén!}

% ············································································
% Laudes
% ············································································

% Hymnus
\defsheet{IamLucisOrto}
  {<-7---7---7-7----7---7----6--5-:-7-7----7--7---4---5--4-3--,-5---5---6--7---6--5----4--5zZ5-:-5--5----4--3--4---5-4--4-?,--4tT4-3\\r\\-.}
  {  Már kél a fény-nek csil-la-ga: e-seng-ve kér-jük az U-rat, jár-jon ma min-de-nütt ve-lünk,  ne ront-sa ár-tás é-le-tünk. Á----men.}

% Antiphona 1
\defsheet[tone={2-d}]{TuamDomine}
  {<-7---7i-----8----:-7--8---9--iI7--8-8---,-5i-uU6--tT4-:-5----5---7---5--4----4t--5-?,}
  {  Ser-kentsd fel, * ha-tal-ma-dat, U-runk, és jöjj el,   hogy meg-sza-ba-díts min-ket!}

% Antiphona 2
\defsheet[tone={8szo-c}]{ConversusEst}
  {<-4---4---4----5-4--5--4--3----:-3-----5--7----7---4---3r-4-?,}
  {  Meg-for-dult a te ha-ra-god, * s_meg-vi-gasz-tal-tál en-gem.}

% Antiphona 3
\defsheet[tone={1re-a}]{ExsultateDeo}
  {<-5--5---5---5--4---5--4---rR3--:-2-3--4--eE2-1----1-?,}
  {  Ör-ven-dez-ze-tek Is-ten-nek, * a mi se-gí--tőnk-nek!}

% Responsorium breve
\defsheet{InMatutinisDomine}
  {<X-3---4--3----4t-5----,-5--5---4--3---rR3-3-?,-5z-4-:-5--5---4--3---4t--tT4-?,-5--5---5-.}
  {   Haj-na-lig, U--ram, * Ró-lad el-mél-ke--dem. Mi-vel vé-del-me-zőm let-tél.   Di-cső-ség...}

% Benedictus
\defsheet[tone={7do-d}]{DaScientiam}
  {<-6--7-----8--7--6----7-6----,-5--rR3-4--5--7iI7-:-7--5--4--4-?,}
  {  Ta-nítsd né-pe-det, U-ram; * bo-csá-na-tá-ra     bű-ne-ik-nek.}

% ============================================================================
% FERIA VI
% ============================================================================

% ············································································
% Officium lectionis
% ············································································

% Hymnus
\defsheet{GalliCantu}
  {<X-3---4---4---1---3-4--4---3--:-6--5--4--5---3r--5-4--,-5--6--7--4---6---5--4--3-:-6----5---4--3---5-zZ5-4-,-3--4--4---1---3---4---4--3---:-6-----5---4--5-----3r---5--4-?,--4tT4-3r-.}
  {   Fel-har-san már a ka-kas-szó, éj sö-té-tet old e jel, és az éj-fél vak ré-me-it  mint fel-le-get ű-zi  el. Jó-sá-gos Úr, hoz-zád té-rünk, s_buz-gón ké-rünk, légy ve-lünk. Á----men.}

% Antiphona 1
\defsheet[tone={8do-c}]{NeInIraTua}
  {<-7---7-uU6-5---:-7--7---7----8--8---7-7-?,}
  {  Har-a-god-ban * ne bün-tess en-gem U-ram.}

% Antiphona 2
\defsheet[tone={2-d}]{ConfundanturEtRevereantrur}
  {<-4---5---7---7--7i--8---:-7i-9---iI7-8--,-8-5---8---7---6--tT4-5--:-3r---5--7--5---5-?,}
  {  Szé-gye-nül-je-nek meg * és fél-je--nek, a-kik lel-kem ke-re--sik, hogy el-ra-gad-ják.}

% Antiphona 3
\defsheet[tone={8szo-c}]{NeDerelinquasMe}
  {<-4--4-----5--7--7---¨-7--7-7----7--7--7--7--:-7--7--zZ5--zZ5-4--4-?,}
  {  Ne hagyj el en-gem * én U-ram, én Is-te-nem, ne tá-vozz-el  tő-lem.}

% ············································································
% Laudes
% ············································································

% Hymnus
\defsheet{DeusQuiCaeliLumen}
  {<X-3--3----4--5-----4---4--3---2-:--5--2---3---1----0---1---3--3-,-3-3------4-5---4----4----3--2-:-4--2-----3--1---0---3----1--1-?,-0\\2eE2W1Q0q\\-.}
  {   Is-ten, ki menny-nek fé-nye vagy és min-den fény-nek kút-fe-je, A-tyánk, e-get tart fenn ka-rod és menny-or-szá-got nyit ke-zed, Á----men.}

% Antiphona 1
\defsheet[tone={1la-a}]{AmpliusLavaMe}
  {<-0----1---3--2----1-qQ0--:-3---4---tT4-3--,-3----4----2---3--2--1--1-?,}
  {  Moss meg en-gem, U-ram, * tel-jes-ség-gel, tisz-títs meg bű-ne-im-től!}

% Antiphona 2
\defsheet[tone={tirr-e}]{DomineAudivi}
  {<-2-1w---3---1---1---:-0-1--3---3--2---1--2---2-?,}
  {  U-ram, hal-lot-tam * a te sza-va-dat és fél-tem.}

% Antiphona 3
\defsheet[tone={tirr-e}]{Lauda}
  {<-2--1------¨-1--0--1---2e---1--3-2-?,}
  {  Di-csérd, * Je-ru-zsá-lem, az U-rat!}

% Responsorium breve
\defsheet{ClamaboAdDominumAltissimum}
  {<X-¢-------------4---3---4--4t-5----,-£--------5--rR3-4---3--3-?,-5--5----5--5z-4--:-5--5---5---4--3---4t-5-?,-5--5---5-.}
  {   A_fölséges_Is-ten-hez ki-ál-tok, * Aki_jóra vi-szi sor-so-mat. El-küld az ég-ből, és meg sza-ba-dít en-gem. Di-cső-ség...}

% Benedictus
\defsheet[tone={1re-b}]{PerViscera}
  {<ô-1--2---2z-6---:-6----6--5--4-5z-6-,-6---6--6--5--4----4t--4x-:-4-3--4---5---rR3-2--1w-2y-?,}
  {   Ir-gal-má-nak * mély-sé-ge-i ál-tal meg-lá-to-ga-tott min-ket  a ma-gas-ság-ból tá-ma-dó.}

% ············································································
% Vesperas
% ············································································

% Hymnus
% TODO

% Antiphona 1
\defsheet[tone={1re-a}]{InclinavitDominus}
  {<-3---3----4--tT4-5-:-2-3---4--eE2-1---1---?,}
  {  Meg-hall-ja az  Úr  i-mád-sá-gom sza-vát.}

% Antiphona 2
\defsheet[tone={2-f}]{AuxiliumMeum}
  {<-3--eE2W1-1-:-0-1--3--3---0e-1----?,}
  {  Az Úr-tól    a mi se-gít-sé-günk.}

% vel ad libitum:
% TODO

% Antiphona 3
\defsheet[tone={8szo-c}]{InConsilioIustorum}
  {<-£---------5--4--5---4-:-3t-7---7--6--5--4u-zZ5-:-4--4----5--5--4--4--4-?,}
  {  Az_igazak ta-ná-csá-ban és gyü-le-ke-ze-té-ben   na-gyok az Úr mű-ve-i.}

% ============================================================================
% SABBATO
% ============================================================================

% ············································································
% Officium lectionis
% ············································································

% Hymnus
\defsheet{LuxAeterna}
  {<-3-4---4--1----3----4--4--3---:-6----5--4---5---3r--5--4--,-5-6---7----4---6--5---4---3-:-6-5-----4--3--5---zZ5-4--,-3---4---4---1--3---4--4---3-:-6--5----4---5---3r--5--4-.}
  {  Ö-rök su-gár, szép ve-rő-fény, fogy-ha-tat-lan nap-vi-lág, a gyá-szos éjt te győ-zöd le, a fényt új-ra meg-ho--zod, pok-lun-kat te rom-bo-lod le, el-ménk ben-ned föl-de-rül.}

% Antiphona 1 (Extra tempus per annum)
\defsheet[tone={8do-c}]{VisitaNosDomine}
  {<-7--7--5----4---45u-7-:-7--5---4--5---4-?,}
  {  Lá-to-gass meg min-ket üd-vös-sé-ged-del!}

% Antiphona 2 (Extra tempus per annum)
\defsheet[tone={1la-a}]{OblitiSuntDeum}
  {<-0--1--3---3---3---4t-4---:-4--2---3r--2---2e-1--1-?,}
  {  El-fe-lej-tet-ték Is-tent, ki meg-men-tet-te ő-ket!}

% Antiphona 3 (Extra tempus per annum)
\defsheet[tone={1la-a}]{SalvosNosFacDomine}
  {<-0-----1---2e--4----eE2-1--:-3------eE2-4t--eE2-qQ0-1-1---eE2-1---0q-1-?,}
  {  Ments meg min-ket, U---ram, gyűjts ösz-sze min-ket a nem-ze--tek kö-zül!}

% Antiphona 1 (Tempore per annum)
\defsheet[tone={3-c}]{ConfiteminiDomino}
  {<-4--5--5---5u-6---5--4--tT4-:--2e---4t-2---2-?,}
  {  Ad-ja-tok há-lát az Úr-nak, * mert jó hoz-zánk.}

% Antiphona 2 (Tempore per annum)
\defsheet[tone={3-c}]{Quoniam}
  {<-7-uU6-tT4-5--4--:-4--4---4---rR3-4t-rR3-2--2-?,}
  {  Ö-rök-ké--va-ló * ir-gal-mas-sá--ga az  Úr-nak.}

% Antiphona 3 (Tempore per annum)
\defsheet[tone={1re-a}]{TuEsDeus}
  {<-5--4----5--4--3----:-3--2---3r-eE2-1---0q--1-?,}
  {  Te vagy az Is-ten, * ki cso-dá-kat cse-lek-szel.}

% ············································································
% Laudes
% ············································································

% Hymnus
\defsheet{DieiLuceReddita}
  {<-7-7---7----7--7---7---6--5---:-7----7---7---7--4--5-4--3---,-5--5----6---7--6---5----4--5zZ5-:-5-----5---4--3---4-----5--4--4-?,-4tT3-2\\e\\-.}
  {  A nap-fény új-ból fel-ra-gyog, csen-dül-jön há-la-é-ne-künk, di-csér-jük Is-ten nagy ne-vét,   Krisz-tus ke-gyel-mét, ér-de-mét. Á----men.}

% Antiphona 1
\defsheet[tone={8szo-c}]{QuamMagnificataSunt}
  {<-7----6---7--5--4---:-6-7--6--5--4---5-4-?,}
  {  Mily fön-sé-ge-sek * a te mű-ve-id, U-ram.}

% Antiphona 2
\defsheet[tone={6-a}]{DateMagnitudinem}
  {<-3--4----4---eE2-1---:-3--4--4---3-?,}
  {  Ma-gasz-tal-já--tok * Is-te-nün-ket.}

% Antiphona 3
\defsheet[tone={1re-a}]{QuamAdmirabileEst}
  {<-0----1---1tz-5--tT4-:-3-4--5--rR3-3r-4-:--3--4-eE2--1----1-?,}
  {  Mily cso-dá--la-tos * a te ne-ved U--runk az e-gész föld-ön.}

% Responsorium breve
\defsheet{ExultabuntLabiaMea}
  {<X-3--3---4--3---4t--tT4--,-4-4--4---5--rR3-4-3--3-?,-[-----------5z--4--:-[-----------4--3--4t--5--?,-5--5---5-.}
  {   Aj-kam ör-ven-dez-zen, * A-mi-kor ne-ked é-ne-kel. Nyelvem_nap-hos-szat igazságodat em-le-ges-se!. Di-cső-ség...}

% Benedictus
\defsheet[tone={1re-a}]{InViamPacis}
  {<-3--3-5---rR3-4--4--5--:-3--2----4---eE2-1-1-?,}
  {  Az i-gaz-ság út-já-ra * ve-zess min-ket U-runk.}


% ============================================================================
% DOMINICA
% ============================================================================

% ············································································
% I. Vesperas
% ············································································

% Hymnus
% Ut in hebdomada I

% Antiphona 1 (ÉE789)
\defsheet[tone={7do-a}]{SitNomenDomini}
  {<ô-4--4----5zZ5-4----:-4--4--4--2--3----2-1---1-?,}
  {   Le-gyen ál---dott * az Úr ne-ve mind-ö-rök-ké.}

% ············································································
% Officium lectionis
% ············································································

% Hymnus
% Ut in hebdomada I

% Antiphona 1
\defsheet[tone={8szo-c}]{EcceDeusMeus}
  {<-4-rF1-3--4t-4--4---:-4--tT4-3---4t--4-,-5-7---7--4--eE2-3-:-1--2--3---4t-5---4-?,}
  {  Ő az  én Is-te-nem * és di--csé-rem őt, a-tyá-im Is-te--ne  és di-cső-í--tem őt.}

% Antiphona 2
\defsheet[tone={1re-a}]{RegnumTuumDomine}
  {<X-3-4--5--6---5---4-5----:-5-4---2--3--4--eE2-1--?,}
  {   A te or-szá-god U-ram, * ö-rök-ké va-ló or--szág.}

% Antiphona 3
\defsheet[tone={8do-c}]{InAeternum}
  {<-7---uU6-5---:-rR3-4----5-6---5---5-4t--5-?,}
  {  Min-den-kor * és  mind-ö-rök-kön ö-rök-ké.}

% ············································································
% Laudes
% ············································································

% Hymnus
% Ut in hebdomada I

% Antiphona 1
\defsheet[tone={4-a}]{RegnavitDominusDecorem}
  {<-2--1---2---4--4t--:-5-5---4---3--4--3--2-?,}
  {  Or-szá-gol az Úr, * é-kes-ség-be öl-tö-zött.}

% Antiphona 2
\defsheet[tone={7la-d}]{Benedicite}
  {<-4uU6-7--iI7-8--8o-8--:-89pP9-8--7---5i--7--6--5--6--5z-4--4--?,}
  {  Áld--já-tok az U-rat * az    Úr min-den mű-ve-i, al-le-lu-ja.}

% Antiphona 3
\defsheet[tone={1re-a}]{Laudate}
  {<-5--5----5--4---:-2--3-4---3-eE2-1-?,}
  {  Di-csér-jé-tek * az U-rat a men-nyből.}

% ············································································
% Hora Media
% ············································································

% Antiphona 1 (Nb41)
\defsheet[tone={8szo-c}]{InTribulatione}
  {<-4-4---3---2--3r--4---:-£----------5---4---5--4-3y-,-3t-7---7----7--5----7--7c-:-8--7--5---4---3r-4í-?,}
  {  A szo-ron-ga-tás-ban * segítségül hív-tam az U-rat, és meg-hall-ga-tott en-gem, és tá-gas-ság-ra vitt.}

% Antiphona 2 (Nb42)
\defsheet[tone={2-g}]{Circumdantes}
  {<ô-4--3---2--2---:-4--4---2---4---5z-5x-,-5--5--6--5--rR3-2-:-2--1----ñ---1w-2-?,}
  {   Kö-rül-vé-vén * kö-rül-vet-tek en-gem, de az Úr ne-vé--ben el-vesz-tém ő--ket.}

% Antiphona 3 (Nb40)
\defsheet[tone={2-g}]{ODomine}
  {<ô-4--rR3E2-2----:-2--2---2--1---2r-4x-,-4-5zZ5-4--:-3--4---5--4-3--2--1w-2-?,}
  {   Ó, U-----ram, * ad-jál üd-vös-sé-get, ó U----ram, ad-jál jó e-lő-me-ne-telt.}

% ············································································
% II. Vesperas
% ············································································

% Antiphona 1 (NZS14)
\defsheet[tone={1la-a}]{DixitDominus}
  {<-11--0--1--3---wW1-q\\Q0\\-:-3---4--tT4-4x-3r-2e\\-2-1í--1í-?,}
  {  Ülj az én job-bom-ra,     * mon-dá az  Úr az én   U-ram-nak.}

% Antiphona 2 (NZS14)
\defsheet[tone={4-a}]{Fidelia}
  {<-2-4t-5---:-44---2--3-4---eE2-1e-3y-,-3---3---3---01e-3y-:-3-2--3---4---3---2-1w--2-?,}
  {  E-rő-sek * mind az ő vég-zé--se-i;   szi-lár-dan áll-nak, i-ga-zak min-den i-dők-re.}

% ============================================================================
% FERIA II
% ============================================================================

% ············································································
% Officium lectionis
% ············································································

% Hymnus
% Ut in hebdomada I

% Antiphona 1 (PSI86)
\defsheet[tone={8do-c}]{DeusDeorum}
  {<-7--uU6-5--4---5--4--4---:-4t-zZ5-4---3r--4-?,}
  {  Az is--te-nek Is-te-ne, * az Úr  meg-szó-lalt.}

% Antiphona 2
\defsheet[tone={tirr-a}]{ImmolaDeo}
  {<X-5--4----4--5--4---4---:-3-4--5---6---4--6--6--5--?,}
  {   Ál-dozd az Is-ten-nek * a di-csé-ret ál-do-za-tát!}

% Antiphona 3
\defsheet[tone={7do-d}]{DixiIniquis}
  {<-7---7----8---6-7--8oO8-7----7----:-5--7--6----4t-4---6u-6----5--5---rR3-3t--4--4-?,}
  {  Azt mond-tam a go-no---szok-nak: * ne be-szél-je-tek go-nosz-sá-got Is--ten el-len.}

% ············································································
% Laudes
% ············································································

% Hymnus
% Ut in hebdomada I

% Antiphona 1
\defsheet[tone={8szo-c}]{BeatiQuiHabitant}
  {<-4---4u--5---rR3-4t--5--4--:-3--4--5---6---tT4-?,}
  {  Bol-dog-ok, a---kik az Úr * há-zá-ban lak-nak.}

% Antiphona 2
\defsheet[tone={2-g}]{Venite}
  {<-4---rR3-2----:-4---4----4---2--4--5--5z--5-:-[------------6---5-rR3-2---:-¡-------------------1--ñ-1w-2-?,}
  {  Jöj-je--tek, * men-jünk fel az Úr he-gyé-re, mert_Sionból jön a tör-vény, és_Jeruzsálemből_az Úr i-gé-je.}

% Antiphona 3 (NZS124)
\defsheet[tone={8do-c}]{CantateDomino}
  {<-¦--------------45u-6----,-7--7---7--6---5--4-5--4-?,}
  {  Énekeljetek_az Úr--nak, * és áld-já-tok az Ő ne-vét.}

% ============================================================================
% FERIA III
% ============================================================================

% ············································································
% Officium lectionis
% ············································································

% Hymnus
% Ut in hebdomada I

% Antiphona 1
\defsheet[tone={7la-d}]{ExsurgatDeus}
  {<-4---6---7--iI7oO8-8---:-8--9--8---uU6-4t--6--7---5--4--4-?,}
  {  Fel-kél az Is-----ten * és el-szé-led-nek el-len-sé-ge-i.}

% Antiphona 2
\defsheet[tone={1la-a}]{DeusNoster}
  {<X-1-0--3--4t-5----¨-5-6---5--4--5--eE2-1-?,}
  {   A mi Is-te-nünk * a sza-ba-dí-tó Is--ten.}

% Antiphona 3
\defsheet[tone={8szo-c}]{InEcclesiis}
  {<-3---3--3--4--4tT4-3---:-5---7--tT4-3--4t-4-?,}
  {  Gyü-le-ke-ze-tek--ben * áld-já-tok az Is-tent!}

% ············································································
% Laudes
% ············································································

% Hymnus
% Ut in hebdomada I

% Antiphona 1
\defsheet[tone={6-a}]{Benedixisti}
  {<-3--4---5---3r---¨-eE2-1--:-3-4--4---3--3-?,}
  {  Meg-ál-dot-tad, * U---ram, a te föl-de-det.}

% Antiphona 2
\defsheet[tone={7do-d}]{DeNocte}
  {<-ik68o-8----:-8oO8I7-7iI7-:-4u-7i--uj5--zZ5-4--4-?,}
  {  Éj----jel  * lel----kem    te-rád vár, Is--te-nem.}


% Antiphona 3
\defsheet[tone={tirr-e}]{IlluminaDomine}
  {<-2--2--1--2-----3--1---:-0-1--3--2--1----2-2-?,}
  {  Vi-lá-go-sítsd re-ánk * a te or-cá-dat, U-runk!}

% ============================================================================
% FERIA IV
% ============================================================================

% ············································································
% Officium lectionis
% ············································································

% Hymnus
% Ut in hebdomada I

% Antiphona 1
\defsheet[tone={8szo-c}]{DeusQuiGlorificatur}
  {<-4u-5-----¨-4---5--4---3-4t-4-:-7-8----8---7--5--7---7--,-©-------------8oO8-uU67iI7-:-[-----------------------7--7---4--4?,}
  {  Is--ten, * kit di-cső-í-te-nek a szen-tek ta-ná-csá-ban, nagy_és_rette-ne---tes       mindazok_fölött,_kik_őt kö-rül-ve-szik.}

% Antiphona 2
\defsheet[tone={8szo-c}]{DeFructu}
  {<-4u--zh4-¨-6u--tT4-tT4-3-:-7--6--7---5---4---4-?,}
  {  Tes-ted * gyü-möl-csé-ből ül-te-tek szé-ked-re.}

% Antiphona 3
\defsheet[tone={8szo-c}]{Potens}
  {<-5--4---3---¨-5u-zh4--6---7--tT4-tT4-3--,-£---------3---5---7--6--6u--5--4--4-?,}
  {  Ha-tal-mas * lé-szen mag-va a   föl-dön, az_igazak nem-zet-sé-ge meg-ál-da-tik.}

% ············································································
% Laudes
% ············································································

% Hymnus
% Ut in hebdomada I

% Antiphona 1
\defsheet[tone={2-f}]{InclinaDomine}
  {<-1---3-----2--3----2-1----:-eE2-3-2---1--0----1e---2---1--1-?,}
  {  For-dítsd ho-zám, U-ram, * fü-le-det és hall-gass meg en-gem.}

% Antiphona 2
\defsheet[tone={8szo-c}]{BeatiQuiAmbulant}
  {<-5---3r-4----:-5èu-7---6u-5----4---5----4---3----4t-4-?,}
  {  Bol-do-gok, * kik tör-vé-nyed sze-rint jár-nak, U--ram.}


% Antiphona 3
\defsheet[tone={tirr-a}]{QuiaMirabila}
  {<X-5----4---5--4---:-3---4--5--6----4z-5-?,}
  {   Mert cso-dá-kat * cse-le-ke-dett az Úr.}

% ============================================================================
% FERIA V
% ============================================================================

% ············································································
% Officium lectionis
% ············································································

% Hymnus
% Ut in hebdomada I

% Antiphona 1
\defsheet[tone={1la-a}]{OmnesInimiciMei}
  {<X-1--1---3--1--0--3----4t--5---:-6---5---4---4--4--4----3t-4x-,-4----4---2---4---tT4-3-:-3----2---3---4--3---2---1í-1í-?,}
  {   El-len-sé-ge-im mind-nyá-jan * hal-lot-ták ba-jo-mat, U--ram, s_ör-ven-dez-tek raj-ta, hogy ezt cse-le-ked-ted ve-lem.}

% Antiphona 2
\defsheet[tone={6-a}]{BenedictusDominus}
  {<-3r-5----eE2-1í-:-3----4-4---3y-?,}
  {  Ál-dott az  Úr * mind-ö-rök-ké.}

% Antiphona 3
\defsheet[tone={1la-a}]{ConvertereDomine}
  {<X-1w--3----4---2-----eE2-1ed1qQ0-:-qQ0-3r---5--4---3--4--3-:-2-3--4----eE2-1--1-?,}
  {   For-dulj hoz-zánk, U---runk,   * és  légy kö-nyö-rü-le-tes a te szol-gá--id-hoz.}

% ············································································
% Laudes
% ············································································

% Hymnus
% Ut in hebdomada I

% Antiphona 1
\defsheet[tone={6-a}]{FundamentaEius}
  {<-3--4--5-3r--eE2-1-:-3-4-----4--3---3y-?,}
  {  Si-on a-lap-ja--i * a szent he-gye-ken.}

% Antiphona 2
\defsheet[tone={2-f}]{EcceDominusNoster}
  {<-1-1---3-2--0q-1---:-1--1--3---2---0---0-,-2--0--1----3---eE2-1-?,}
  {  Í-me, a mi U-runk * ha-ta-lom-mal jön el, és ma-gasz-tos kéz-zel.}


% Antiphona 3
\defsheet[tone={tirr-e}]{IustusEtSanctus}
  {<-2-2---1--wW1---:-0-1--2--2--2-?,}
  {  I-gaz és szent * a mi Is-te-nünk.}

% ============================================================================
% FERIA VI
% ============================================================================

% ············································································
% Officium lectionis
% ············································································

% Hymnus
% Ut in hebdomada I

% Antiphona 1
\defsheet[tone={1la-a}]{ZelusDomusTuae}
  {<X-1-1--1--3----1--0--3---4t-5-:-6---5-4----4----3t-4--,-4-----4--4--2--4---5--4--3-:-2---3---4---eE2--1-1-?,}
  {   A há-za-dért va-ló buz-gó-ság meg-e-mész-tett en-gem, s_gya-lá-zó-id-nak át-ka-i   rám-hul-lot-tak, U-ram. }

% Antiphona 2
\defsheet[tone={8do-c}]{ConsolantemMe}
  {<-7--uU6-5----¨-¦-------------zZ5-4-:-4--2---3--4---4--,-4--4-4t--3--5---7--uU6-5-4--5--tT4-:-2--2----3r-2--eE2-1-:-1-2---3---4-3---4t--4--4-?,}
  {  Ke-res-tem, * aki_megvígasz-tal-na  és nem ta-lál-tam: és e-pét ad-tak ne-kem e-le-de-lül   és szom-jú-sá-gom-ban e-cet-tel i-tat-tak en-gem.}

% Antiphona 3
\defsheet[tone={8szo-c}]{QuaeriteDominum}
  {<--4--4---4--3---5--7-uU6-¨-6--7--5-5--6---5--4-?,}
  {   Ke-res-sé-tek az U-rat * és él a ti lel-ke-tek.}

% ············································································
% Laudes
% ············································································

% Hymnus
% Ut in hebdomada I

% Antiphona 1
\defsheet[tone={8szo-c}]{TibiSoli}
  {<-2--3r-4x---:-3t---¦-----------5---6uc-:-5--5---5----6---5----4-4-?,}
  {  El-le-ned, * csak ellened_vét-kez-tem,  kö-nyö-rülj raj-tam, U-ram.}

% Antiphona 2
\defsheet[tone={8szo-c}]{NonNosDerelinquas}
  {<-7--6-----tT4-5---4----:-3-4t---6--5--4-?,}
  {  Ne hagyj el  min-ket, * U-runk so-ha-se.}

% Antiphona 3
\defsheet[tone={tirr-a}]{ServiteDomino}
  {<X-3----5---6--5---4--4--3---:-3--3---4t--5-?,}
  {   Szol-gál-ja-tok az Úr-nak * ví-gas-ság-gal.}

% ============================================================================
% SABBATO
% ============================================================================

% ············································································
% Officium lectionis
% ············································································

% Hymnus
% Ut in hebdomada I

% Antiphona 1
\defsheet[tone={4-a}]{ConfiteanturDomino}
  {<-1--2--3---4--2e--4--3--2----:-0--1---2---3--eE2-2-?,}
  {  Ad-ja-nak há-lát az Úr-nak, * ir-gal-mas-sá-gá--ért.}

% Antiphona 2
\defsheet[tone={8do-c}]{DeNecessitatibusNostris}
  {<-7---7---7--5--4--7---7---:-¦-------------8---8----7-7-?,}
  {  Szo-ron-ga-tá-sa-ink-ból * szabadíts_meg min-ket, U-runk.}

% Antiphona 3
\defsheet[tone={1la-a}]{QuisIntelleget}
  {<-1--qQ0-34t-tT4--:-3--3---3---4--2--3-2----1-1-?,}
  {  Ki ér--ti  meg, * ir-gal-mas-sá-ga-i-dat, U-runk?}

% ············································································
% Laudes
% ············································································

% Hymnus
% Ut in hebdomada I

% Antiphona 1
\defsheet[tone={3-b}]{InMatutinisDomineMeditabor}
  {<XB-3---4--3--4z--6c---:-zZ5-4x---3--2---3---4---eE2W1-1í-?,}
  {    Már ko-ra reg-gel, * U---ram, ró-lad gon-dol-ko----zom.}

% Antiphona 2
\defsheet[tone={2-c}]{MecumSitDomine}
  {<-7--6----7--5----tT4-6----:-6-7--5---4----6u-5-?,}
  {  Le-gyen ve-lem, U---ram, * a te böl-cses-sé-ged.}

% Antiphona 3
\defsheet[tone={4-a}]{VeritasDomini}
  {<-2--1--2-4---4t-5--:-5-4---3--4---3--2-?,}
  {  Az Úr i-gaz-sá-ga * ö-rök-ké meg-ma-rad.}
% ============================================================================
% DOMINICA
% ============================================================================

% ············································································
% Vesperas I.
% ············································································

% Hymnus
% Ut in hebdomada II

% Antiphona 1
\defsheet[tone={1re-a}]{FiatPax}
  {<-5--4t---5--4---:-2-3--4-eE2-1-?,}
  {  Le-gyen bé-ke, * a te e-rőd-ben.}

% Antiphona 2 (SzKD66)
\defsheet[tone={8szo-c}]{Sustinuit}
  {<-4--4u---5--4---3-4t--4-:-4--5--4-5--rR3-3y-,-3--5--7---7---6u-5--4--3--4t-4-?,}
  {  Re-mény-ke-dik a lel-kem az Úr i-gé-jé--ben, az én lel-kem ő--rá vá-ra-ko-zik.}

% Antiphona 3
% Ut in hebdomada II

% ············································································
% Officium lectionis
% ············································································

% Hymnus
% Ut in hebdomada II

% Antiphona 1
\defsheet[tone={8szo-c}]{DominiEstTerra}
  {<-4--4u-tT4R3-4t-4-----¨-7--8---7--5u--6--,-7-8----7--8oO8-uU67iI7-:-4--5-----5-5---7--7---4---4-?,}
  {  Az Ú--ré    a  föld, * és ben-ne min-den, a föld-ke-rek--ség       és mind, a-kik ab-ban lak-nak.}

% Antiphona 2
\defsheet[tone={6-a}]{BenediciteGentes}
  {<-3---4--5----rR3-4--1----:-3--4--4---3-?,}
  {  Áld-já-tok, nem-ze-tek, * Is-te-nün-ket.}

% Antiphona 3
\defsheet[tone={2-d}]{Introibo}
  {<-5--7i-8---:-7-8--9--8---7---8-8--,-5i-uU6-5---4--tT4-:-3r----5----7--5---5-?,}
  {  Be-lé-pek * a te há-zad-ba, U-ram, és i---mád-ko-zom   szent temp-lo-mod-ban.}

% ············································································
% Laudes
% ············································································

% Antiphona 1 (D39)
\defsheet[tone={7la-a}]{OmnisSpiritus}
  {<ô-1---3---5z-5--:-4--2----3--2--1-1í-?,}
  {   Min-den é--lő * di-csér-je az U-rat.}

% ············································································
% Hora Media
% ············································································

% Antiphona 3 (vö. GH363)
\defsheet[tone={7do-a}]{Vovete}
  {<ô-5--5---3---5--6u-zZ5-5---:-5--3--4t-wW1-1-,-4--4-4--2r--3-:-2--1--2rR3-1---1-?,}
  {   Te-gye-tek fo-ga-dal-mat * és ad-já-tok meg ti U-ra-tok-nak Is-te-ne---tek-nek.}

% ············································································
% Vesperas
% ············································································

% Antiphona 2 (NZS14)
\defsheet[tone={8do-c}]{InMandatisEius}
  {<-7---6---7--5x---:-4--3--5--5--6---5----4---4--4x-?,}
  {  Bol-dog em-ber, * ki az Úr pa-ran-csát ked-ve-li.}

% ============================================================================
% FERIA II
% ============================================================================

% ············································································
% Laudes
% ············································································

% Antiphona 1 (D60)
\defsheet[tone={6-a}]{DomineRefugium}
  {<X-3-4----5--3---3---:-3-1--3--4--4--3?,}
  {   U-ram, te let-tél * a mi me-ne-dé-künk.}

% ············································································
% Hora Media
% ············································································

% Antiphona 2 (Nok145)
\defsheet[tone={8do-c}]{IusteIudicate}
  {<-7-7--7---5-4---45u-7----:-7--7--5--4---5--4-4x-?,}
  {  I-ga-zul í-tél-je--tek, * ti em-be-rek fi-a-i!}

% Antiphona 3 (NZS256)
\defsheet[tone={2-g}]{Clamavi}
  {<ô-4--4--rR3E2-2y---:-2--1---2----4--4----1r-2y-?,}
  {   Ki-ál-tot---tam, * és meg-hall-ga-tott en-gem.}

% ············································································
% Vesperas
% ············································································

% Antiphona 1 (D184)
\defsheet[tone={3-c}]{QuoniamInAeternum}
  {<-7-uU6-tT4-5--4x-:-4--4---4---rR3-4t-rR3-2y-2y-?,}
  {  Ö-rök-ké--va-ló * ir-gal-mas-sá--ga az  Úr-nak.}

% ============================================================================
% FERIA III
% ============================================================================

% ············································································
% Officium lectionis
% ············································································

% Antiphona 1 (NZS150)
\defsheet[tone={1re-a}]{ClamorMeus}
  {<-5--4--5--rR3-:-2---3---rR3-4t---eE2-1í-1í-?,}
  {  Ki-ál-tá-som * jus-son e---léd, ó   Is-ten!}

% Antiphona 3 (SzKD87)
\defsheet[tone={8szo-c}]{IustiConfitebuntur}
  {<-£---------5--4---4--3---:-5-7--7--uU6-tT4--5-4-,-£----------------eE2-1-:-3---3---3-4t---4x--4x-?,}
  {  Az_igazak há-lát ad-nak * a te ne-ved-nek, U-ram s_az_egyenes_szí-vű--ek  szí-ned e-lőtt lak-nak.}

% ············································································
% Vesperas
% ············································································

% Antiphona 1 (D184)
\defsheet[tone={8szo-c}]{HymnumCantate}
  {<-4-4--4---5--4---tT4-3y----:-3--5--7--7--7-6u--5--4x-?,}
  {  É-ne-kel-je-tek him-nuszt * Si-on da-la-i-ból ne-künk.}

% Antiphona 2 (D193)
\defsheet[tone={5-c}]{InConspectu}
  {<-5--5---4---3---5---7--7iI7U6-5----,-5----5--5--4---5--4----3-3y-?,}
  {  An-gya-lok-nak szí-ne e------lőtt * zsol-tá-ro-zok né-ked, U-ram.}

% ============================================================================
% FERIA IV
% ============================================================================

% ············································································
% Officium lectionis
% ············································································

% Antiphona 1 (Nok213)
\defsheet[tone={1re-a}]{Benedic}
  {<-5---4----¨-2e-4---3----2--1-1-?,}
  {  Áld-jad, * én lel-kem, az U-rat!}

% ============================================================================
% SABBATO
% ============================================================================

% ············································································
% Hora Media
% ············································································

% Antiphona 1
\defsheet[tone={1la-a}]{FiatManusTuaDomine}
  {<X-0-1---1tu-5----5--rR3--4t-5----:-3----2---3---4--3---2---1--1--,-3----2-1--1--2---3---4tT4-3-:-2e-4----eE2-1-?,}
  {   A-zon le--gyen ke-zed, U--ram, * hogy meg-sza-ba-dít-son en-gem; mert a te pa-ran-csa-i----dat vá-lasz-tot-tam.}

% Antiphona 2
\defsheet[tone={6-a}]{Eructavit}
  {<X-3--4--5---3---:-4--eE2-1--3---3---4---4-3-?,}
  {   Az én szí-vem * ün-ne--pi szó-zat-tól á-rad.}

% Antiphona 3
\defsheet[tone={4-a}]{VidiIerusalemNovam}
  {<-2---2---2--1--3--2--qQ0-1w-2----:-1----2-3---4t-4---4---5-4--3--4----4----rR3-2-?,}
  {  Lát-tam az új Je-ru-zsá-le-met, * mint a fér-jé-nek föl-é-ke-sí-tett meny-as--szonyt.}

\makeatother
% ============================================================================
% IUNIUS
% ============================================================================

% ············································································
% Die 24: In Nativitate Sanctis Ioannis Baptistae
% ············································································

% --- Laudes ---

% Antiphona 3
\defsheet[tone={7la-a}]{TuPuer}
  {<ô-1---1t---5----:-5-5--6---7---zZ5-5-:-5---3--4--5--wW1-1---,-£--------------2r---3---:-2--1--2---2--2--4--3-1--1-?,}
  {   Te, gyer-mek, * a Ma-gas-ság-be--li  pró-fé-tá-ja le--szel, az_Úr_orcája_e-lőtt mégy, el-ké-szí-te-ni az ő út-ját.}

% ············································································
% Die 27: S. Ladislai
% ············································································

% --- Officium Lectionis ---

% Antiphona 1
\defsheet[tone={1la-a}]{HicIniquisNonConsesnsit}
  {<-0--1---1tu-5---¨-tT4-3---4tT4-eE2W1-,-1--0--34t-5---44-3--44--1---,-44t-4-eE2-1---2----0--1---1-,-5--4---5---7i---78o-uU6Z5-4t--5-,-5-4---55--3-----4-3---4---5---,-5-4--eE2-1----2ed1-0q--1---1-?,}
  {  Go-no-szok-kal * nem cim-bo---rált,   út-ja-i---kon so-ha-sem járt. Úr  i-gá--ját váll-ra vet-te; Is-ten ked-vét: azt ke----res-te. A-mit mun-kált, a-mit gon-dolt, é-le-te  mind ta---ní--tás volt.}

% Antiphona 2
\defsheet[tone={2-f}]{RexAChristoConstitutusEst}
  {<-0--1--12ed1-1---¨-qQ0---ð---0q--1-,-0--1--23r-rR3E2W1-2--0--1---1-,-3-4--4----3--4--3---4t--5-,-5-4-----eE2-1----qQ0-ð---0q--1-?,}
  {  Őt ki-rály--lyá * Krisz-tus tet-te, tá-mo-gat-ta,     se-gí-tet-te, e-ré-nyek-be öl-töz-tet-te, a szent he--gyen vé--del-mez-te.}

%% Antiphona 3
%\defsheet[tone={4-a}]{LaudateTeRexGloriae}
%  {<-3---4z--3--rR3--1w-3-3---,-3r--5?,}
%  {  Áld-jon té-ged, ég U-ra, * Egy-há-zunk szent szó-za-ta, ün-ne-pel-jen zeng-ve ma é-des é-nek dal-la-ma.}

% ============================================================================
% IULIUS
% ============================================================================

% ············································································
% Die 3: S. Thomae
% ············································································

% --- Laudes ---

% Antiphona 1
\defsheet[tone={8szo-b}]{StetitIesus}
  {<-4---3--5----7--7---:-4-6--7---5--4----tT4-3---,-4---4---4---3--5----7--7--7--:--4--6u-tT4-3---4t-5--4x-4x-?,}
  {  Meg-ál-lott Jé-zus * a ta-nít-vá-nyok kö--zött: "Bé-kes-ség le-gyen ve-le-tek!" Al-le-lu--ja, al-le-lu-ja.}

% Antiphona 2
\defsheet[tone={1la-a}]{ThomasQuiDiciturDydimus}
  {<X-1--1----1--3-1---0---3---4t-5----:-5---5----5u-5--:-5-4--3---2e-4----eE2-1--,-1------1---3--1--0-3--4---3t-5---:-5---4---4t-eE2-qQ0--1--eE2-1í-1í-?,}
  {   Ta-más, ki I-ker-nek mon-da-tik, * nem volt ve-lük, a-mi-kor el-jött Jé--zus. S_mond-ták ne-ki a ta-nít-vá-nyok: lát-tuk az U---rat, al-le--lu-ja.}

% Antiphona 3
\defsheet[tone={4-a}]{MitteManumTuam}
  {<ô-1-------2--4-4--5z-5----:-4-5------6-5---4---5--5---,-5--2--5----4--3---2--1---wW1-:-0--1---2-–2e--4--3--2--2-?,}
  {   Nyújtsd ki a ke-ze-det, * é-rintsd a sze-gek he-lyét, és ne légy im-már hi-tet-len,  ha-nem hí-vő, al-le-lu-ja.}

% Benedictus
\defsheet[tone={8szo-c}]{QuiaVidistiMe}
  {<-4--4---3--5---7--7----:-4--6u---5--4---tT4-3--,-4---4--4----3-5---7---7---7--:-4--6u---5----6--5z-4x-4x-?,}
  {  Mi-vel lá-tál en-gem, * Ta-más, te-hát hit-tél, bol-do-gok, a-kik nem lát-nak, és hisz-nek, al-le-lu-ja.}

% --- Vesperas ---

% Magnificat
\defsheet[tone={8szo-c}]{Misi}
  {<-4--4----4u--5tT4-¨-3-4t-5--4-:-7-8---uj5-7--7---7-:-7--8--7--8oO8-uU67iI77-:-5--5-7--7--4x-4x-:-4-----eE2-3--:-1--2-3----4--4--rR3--4t-5--4x-4x-?,}
  {  Ki-nyúj-tot-tam  * a ke-ze-met a szö-gek he-lyé-be, és uj-ja-i----mat        az ő ol-da-lá-ba;  s_mon-dot-tam: én U-ram, Is-te-nem, al-le-lu-ja.}

% ············································································
% Die 11: S. Benedicti
% ············································································

% --- Officium Lectionis ---

% Hymnus
\defsheet{InterAeternas}
  {<X-5-4--5--4----1í--:-4--4---3---4--5x--5x-,-5---7--8----7---5x-:-5--4-3---4--5x-5x-,-5---4--5-------4---1í-:-4--4--3---4--tT4-3y-,-3--3--2--1í-1í-?,-1wW1-0\\q\\-.}
  {   É-gi fé-nyek közt, ki-ket áld-va hir-det  har-ci győz-tes-ként di-a-dal-mi é--nek, ott ra-gyogsz, bol-dog  Be-ne-dek su-gár-zó   ér-de-me-id-del.  Á----men.}

% --- Laudes ---

% Hymnus
\defsheet{LegiferPrudens}
  {<X-5-----4--5--4-----1í--:-4--4--3----4--5x-5x-,--5-----7--8----7--5x-:-5-4--3---4--5x-----5x--,-5---4--5--4---1í-:-4-4--3----4--tT4--3y-,-3--------3--2---1í--1í-?,-1wW1-0\\q\\-.}
  {   Bölcs ta-ní-tónk, Szent Be-ne-dek, vi-lá-gíts, Krisz-tu-sunk fé-nyét a vi-lág-ba hintsd szét, hí-ven is-mer-nünk e hi-tet, se-gíts meg, s_töltsd be szí-vün-ket!  Á----men.}

% Benedictus
\defsheet[tone={8do-c}]{FuitVir}
  {<-7-5u-7--¨-7iI7-uU6u-uU6-5--6u--5--4---,-rR3-4--3--5--uU67iI7-:-5u-6---5---4---3r-4-?,}
  {  É-le-te * tisz-te---let-re-mél-tó volt, ne--ve Be-ne-dek:      ke-gye-lem-mel ál-dott.}

% ············································································
% Die 15: S. Bonaventurae
% ············································································

% Benedictus
\defsheet[tone={2-f}]{SpirituSapientiae}
  {<-1-1äe--1---¨-1---1----qQ0-1--1--2--3---2---0--0-:-2--4---3---1--3r-eE2-1-?,}
  {  A hasz-nos * böl-cses-ség és ér-te-lem lel-ké-vel be-töl-töt-te őt az  Úr.}

% ============================================================================
% AUGUSTUS
% ============================================================================

% ············································································
% Die 8: S. Dominici
% ············································································

% --- Laudes ---

% Hymnus
\defsheet{NovusAthletaDomini}
  {<ô-2--2---2--1---2---4---3--2í-:-2--3---4----5--4---5---5--6x-,-6--6--5---6---4--5--rR3-2í-:-2--1----2-4--5---2---1--2í-?,-1e4R3E2W1-2í-.}
  {   Di-cső ne-vet mél-tán vi-selt Do-mon-kos, Is-ten baj-no-ka,  az Úr-hoz fűz-te őt ne--ve,  és lett a jó hír har-co-sa.   Á---------men.}

% Benedictus
\defsheet[tone={1re-a}]{Beatus}
  {<-0q--1tz-5--5z-4tT4---:-4-5--6--5--rR3-3rR3-1---:-4--5--4--3---eE2--1---2---0-,-0--1--3-:-3-----3--4t--wW1-3--2---1--1-?,}
  {  Bol-dog az a  szent, * a-ki az Úr-ban-bí---zott, az Úr pa-ran-csát hir-det-te, és az ő   szent he-gyé-re  ál-lít-ta-tott.}

% ············································································
% Die 24: S. Bartholomaei
% ············································································

% --- Laudes ---

% Hymnus
\defsheet{RelucensInterPrincipes}
  {<-5--5--5--5--1---3--2--1í--:-5--5--6---4----5-----7----6--5x-,-5-----6---7--8----7----5---7--zZ5T4x-:-5---5--6---4--2---4---5--5x-?,-5zZ5-4\\t\\-.}
  {  Is-ten nagy há-za né-pe közt ki-rá-lyi fény-ben tün-dö-kölsz, Szent Ber-ta-lan, fi-gyelj re-ánk, ké-rünk, míg így ma-gasz-ta-lunk:  Á----men.}

% ============================================================================
% BEATEAE MARIAE VIRGINIS
% ============================================================================

% ············································································
% Officium Lectionis
% ············································································

% Antiphona 1
\defsheet[tone={1la-a}]{QuisAscendet}
  {<-0--1--1tu-5---:-[-----------tT45uj5T4-:-5--4--3---3r--3-:-2e-4-eE2---1--1---,-eE2-1--2---3---4tT4-3-:-2e-4----eE2-1---1-?,-------3--3----3---2---3---4--eE2-1--1-?,}
  {  Ki me-het fel * az_Úr_hegyé-re,         és ki áll-hat meg az ő szent he-lyén? Az  ár-tat-lan ke---zű  és tisz-ta  szí-vű. <T.P.> és tisz-ta  szí-vű, al-le--lu-ja.}

% Antiphona 2
\defsheet[tone={6-a}]{Sanctificavit}
  {<-3---4----5---4--¨-4t-rR3-3--4-5---4--3-?,--------4--4t-3---3-?,}
  {  Meg-szen-tel-te * az Úr  az ő haj-lé-kát. <T.P.> Al-le-lu--ja.}

% Antiphona 3
\defsheet[tone={5-c}]{Gloriosa}
  {<-¦----------7iI7U6u-:-5---7--6---4t-rR3-4t---4--3--3-?,------4--4t-3--3-?,}
  {  Dicsőséges-nek       mon-da-nak té-ged Szűz Má-ri-a. <T.P.> Al-le-lu-ja.}

% ============================================================================
% APOSTOLORUM
% ============================================================================

% ············································································
% Invitatorium
% ············································································

\defsheet{RegemApostolorum}
  {<-0-0ãwe-2--:-2--2-2---2--3r--rf2-3-1wW1Q0\\-,-0wår-4--4t---5-tT4R3rR3E2-2-.}
  {  U-run--kat, az a-pos-to-lok Ki-rá-lyát,    * jöj--je-tek, i-mád--------juk!}

% ············································································
% Officium Lectionis
% ············································································

% Hymnus
\defsheet{OSempiternaeGloriae}
  {<ô-4--4-4--4--4---4--3---2í-:-4-4---4--4--1-----2----1--0í-,-2--2---3--4----3--2----1--2eE2í-:-2---2---1--0---1----2----1--1í-?,--1wW1-0\\q\\-.}
  {   Az é-gi ud-var dí-sze-i,   ö-rök Ki-rá-lyunk nagy-ja-i,   ki-ket ne-künk ne-velt az Úr,     s_a-pos-to-lok-ként ránk ha-gyott, Á----men.}

% Antiphona 1
\defsheet[tone={2-f}]{InOmnemTerram}
  {<-1e--2---3r---4--:-4--3---1--3-4---1--1--,-1-3---2--3---4---4---3--3--1e-2e-4-1--1-?,--------0--0w-1--1-?,}
  {  Min-den föld-re * el-jut az ő zen-gé-sük, a föl-ke-rek-ség szé-lé-ig az ő  i-gé-jük. <T.P.> Al-le-lu-ja.}

% Antiphona 2
\defsheet[tone={8szo-c}]{AnnuntiaveruntOpera}
  {<-3---5u-zh4-:-6--7---tT4-tT4-3-,-4--3--5-7--6--4--6u--5--4-?,--------6u--5--4t---6--5z-4--4-?,}
  {  Hir-de-tik * Is-ten te--te--it  és az ő te-te-it meg-ér-tik. <T.P.> meg-ér-tik, al-le-lu-ja.}

% Antiphona 3
\defsheet[tone={8szo-c}]{AnnuntiaveruntIustitiam}
  {<-rR3d1-34t-4---:-4--6i-7-5---77-6--,-88-7---8oO8-uU67i7-:-5---5--4--5--5èuj5-zZ5-4-?,--------5--5èuj5-zZ5-4t---6--5--4--4-?,}
  {  Hir---de--tik * az ő  i-gaz-sá-gát, és min-den  nép      meg-lá-ta di-cső---sé--gét. <T.P.> di-cső---sé--gét, al-le-lu-ja.}

% ············································································
% Laudes
% ············································································

% Antiphona 1
\defsheet[tone={1la-a}]{HocEstPraeceptumMeum}
  {<-1--0--1--3--wW1-qQ0---:-3---4---5--4---tT4-3---:-3-3r---33-2---3---4--eE2-1--1--1-?,}
  {  Ez az én pa-ran-csom: * sze-res-sé-tek egy-mást, a-mint én sze-ret-te-lek ti-te-ket.}

% Antiphona 2
\defsheet[tone={1la-a}]{MaioremCaritatem}
  {<X-1--0-----3---4--3t-5--:-6---5--4---4-----3t-4--,-4----4--4-2--4---tT4-3-:-2--3-4--eE2-1--0q-1-?,}
  {   Na-gyobb sze-re-te-te * sen-ki-nek sincs an-nál, mint ki é-le-tét ad--ja  az ő ba-rá--ta-i--ért.}

% Antiphona 3
\defsheet[tone={1la-a}]{VosAmiciMeiEstis}
  {<X-1--1--3--1--0--34t--5----:-4--6---5--4---4---3t---4--,-4-4---4--4---2----4---tT4-3--:-2----3--4--eE2-1-----1-?,}
  {   Ti ba-rá-ta-im vagy-tok, * ha meg-te-szi-tek mind-azt, a-mit pa-ran-csol-tam nek-tek, mond-ja az Úr  Krisz-tus.}

% Responsorium breve
\defsheet{ConstituesEosPrincipes}
  {<X-¢------------4--3----4t-tT4-,-5--4-3----rR3-3-?,-5--5--5----5z-5--:-5---4---4-,-5---4---3--4t--tT4-?,-[-------------5z--5-:-5-5--4-4-:-5--5-4-----3--4t--tT4-.}
  {   Fejedelmekké te-szed ő--ket * Az e-gész föl-dön. Ne-ve-det, U--ram, min-den-kor hir-det-ni fog-ják.   Dicsőség_az_A-tyá-nak a Fi-ú-nak és a Szent-lé-lek-nek.}

% Benedictus
\defsheet[tone={2-c}]{DumSteteris}
  {<-5-5--tG2-4t-5z-5----¨-5--5-5---6uU6-55T4-5-7----5z-6---;-6--8---8---9--iI7-uU6-5--5-:-7---uU6-tT4-5--6u-5----5-,-5----7---7-6--tT4-55--4-:-4-5---77-6u-8o--uU6-5----4t-5-?,}
  {  A-mi-kor ki-rá-lyok * és e-löl-já---rók  e-lőtt ál-ltok, ne ter-vez-zé-tek e---lő-re, mit fog-tok vá-la-szol-ni: mert meg-a-da-tik nék-tek a-zon ó-rá-ban, mit mond-ja-tok.}

% ············································································
% Vesperas II.
% ············································································

% Responsorium breve
\defsheet{AnnuntiateInterGentes}
  {<X-¢-----------------4--3---4t-tT4--,-5--5--4--3---rR3-3-?,-[-----------5z-4-,-4---5---rR3-4t-tT4-?,-[-------------5z--5-:-5-5--4-4-:-5--5-4-----3--4t--tT4-.}
  {   Hirdessétek_a_nem-ze-tek kö-zött * Az Úr di-cső-sé--gét. És_csodatet-te-it  min-den nép kö-zött!  Dicsőség_az_A-tyá-nak a Fi-ú-nak és a Szent-lé-lek-nek.}

% ============================================================================
% UNIUS MARTYRIS
% ============================================================================

% ············································································
% Officium Lectionis
% ············································································

% Antiphona 1
\defsheet[tone={1la-a}]{SiQuisPreseveraverit}
  {<-qQ0-3--rR3-4t--5---:-6----5--4---3t-4---:-4-2---4---tT4-3-:-2--3r--eE2-1-?,--------0--0q-1--1-?,}
  {  A---ki a---zon-ban * mind-vé-gig ki-tart, é-let-ben ma--rad és nem hal meg. <T.P.> Al-le-lu-ja.}

% Antiphona 2
\defsheet[tone={8szo-g}]{Existimo}
  {<ô-1r-2---wW1Q0-1w-2---1---:-5----4--2--4--4-:-[----------------------------6---5---rR34tT4-:-¡------------------4--4---1---1-?,---------1--1--]ñ-ð-:-0--1w-1--1-?,}
  {   En-nek az    é--let-nek * szen-ve-dé-se-i   nem_mérhetők_az_eljövendő_di-cső-ség-hez,      amely_majd_megnyil-vá-nul raj-tunk. <T.P.> Al-le-lu--ja, al-le-lu-ja.}

% Antiphona 3
\defsheet[tone={7do-a}]{TamquamAurumInFornace}
  {<ô-5----5-5--3----5-6--5--5----:-[-------------------3---4--5--wW1-1-:-2--4-4--4--4--2r--3-:--¡--------------4----3-1---1-?,}
  {   Mint a-ra-nyat a ko-hó-ban, * megvizsgálta_válasz-tot-ta-it az  Úr, és é-gő-ál-do-zat-ként elfogadta_őket mind-ö-rök-ké. }

% ············································································
% Laudes
% ············································································

% Antiphona 1
\defsheet[tone={1re-a}]{QuiaMeliorEst}
  {<-1--qQ0-3r-3---4t--:-5u---5--4--5----4--rR3-4t-,-1----0--1----eE2-4---3--1---0q-1-?,}
  {  Mi-vel ir-gal-mad * több-et ér mint az é---let  hadd ma-gasz-tal-jon té-ged aj-kam.}

% Antiphona 2
\defsheet[tone={tper-g}]{MartyresDomini}
  {<-1-0----1---3--3r-4--:-5---5--4---3--2-1---3----4t-4---4--?,}
  {  U-runk vér-ta-nú-i, * áld-já-tok az U-rat mind-ö--rök-ké.}

% Antiphona 3
\defsheet[tone={1re-a}]{QuiVicerit}
  {<-0q-1tz--5-----:-56u-5--XtT45zZ5T4-3r-4t-1-:--2----3--4t--4tT4R3-2e---3rR3E2-0q-1-?,}
  {  A  győz-test, * osz-lo-pom--------má te-szem temp-lo-mom-ban    mond-ja     az Úr.}

% Benedictus
\defsheet[tone={3-c}]{QuiOdit}
  {<-4-4--456u-5--5--:-7-uU6-5z--5-4---5-uj5-6u-6--,-6u--8--7--uU6-tT4-:-tT4R3-4---rR3-2---2-?,}
  {  A-ki gyű--lö-li * é-le--tét e-zen a vi--lá-gon, meg-őr-zi azt az    ö-----rök é---let-re.}

% ============================================================================
% SANCTORUM VIRORUM
% ============================================================================

% ············································································
% Invitatorium
% ············································································

\defsheet{MirabilemInSanctisSuisDominum}
  {<-0ã2e-2---:-2-2--23r-rR3-4---4t--rR3---1w-2e-2----,-1t--5z-5----5-tT4R3rR3E2-2-?,}
  {  Is---tent, a-ki cso-dá--la--tos szent-je-i--ben, * jöj-je-tek, i-mád--------juk!}

% ············································································
% Officium Lectionis
% ············································································

% Antiphona 1
\defsheet[tone={8szo-c}]{Vitam}
  {<-3-4t-4---¨-4t---7--5--:-7--7---7--tT4-5z-5---tT4-4--,-4--7---7i-6---7--5----4-4t--rR3-3---3---4t--uj5-4-,-4--eE2-1--3--4--5---uU6-5z--7---5--4--4-?,}
  {  É-le-tet * kért tő-led, és meg-ad-tad ne-ki, U---ram, di-cső-sé-get és nagy é-kes-sé--get tet-tél rá--ja, fe-jé--re ko-ro-nát rak-tál drá-ga-kő-ből.}

% Antiphona 2
\defsheet[tone={7do-d}]{IustumDeduxit}
  {<-8--8-6--8---9-8---8---8-7i-8---:-8--6---7--uU6-4-,-¦-------------5u-6-:-5--4---5u-6---4-?,}
  {  Az i-ga-zat e-gye-nes u-ta-kon * ve-zet-te az Úr, és_megmutatta ne-ki  Is-ten or-szá-gát.}

% Antiphona 3
\defsheet[tone={8szo-c}]{IustusUtPalma}
  {<-[-------------4-tT4-3--:-3tè7i-uU6-7iiI7-,-7----7-7--6--5---uU6-tT4-tT4-:-3r---tT4-4-?,}
  {  Az_igaz,_mint a pál-ma * vi----rág-zik,    mint a Li-ba-non céd-ru--sa    nagy-ra  nő.}

% ············································································
% Laudes
% ············································································

% Antiphona 1
\defsheet[tone={1la-a}]{Iucunditatem}
  {<-2--0---1tu-5---:-[----------tT45uj5T4-:-3r--5---eE2-1-----,-eE2-1-1---1--2---3--4tT4-3---:-3-3---3rR3-0-:-2-3--4--eE2-1-?,}
  {  Vi-gas-sá--got * és_örvende-zést        hal-moz rá--ja/juk, és  ö-rök ne-vet ad ne---ki(k) ö-rök-sé---gül a mi Is-te--nünk.}

% Antiphona 2
\defsheet[tone={tper-g}]{ServiDomini}
  {<-0--1--3----3r-4-:--5--4----5--4---3--2-1---3----4t-4---4-?,}
  {  Az Úr szol-gá-i, * di-csér-jé-tek az U-rat mind-ö--rök-ké.}

% Antiphona 3
\defsheet[tone={4-a}]{Exsultabunt}
  {<-2--2----2--1äeE2s0-:-2-2e---4---3-2--2e--4---4--,-4--4--3--eE2-1--1---3----1--3--4---2-?,}
  {  Uj-jong-ja-nak     * a szen-tek a di-cső-ség-ben, és vi-ga-doz-za-nak nyug-vó-he-lyü-kön!}

% Responsorium breve
\defsheet{LexDeiEius}
  {<X-33r-3-:-4---4t-tT4-,-5--4-3---rR3-3-?,-5z--4-,-5--5--5---4---3---4t-tT4-?,-[-------------5z--5-:-5-5--4-4-:-5--5-4-----3--4t--tT4-.}
  {   Is--ten tör-vé-nye * Él a szí-vé--ben. Nem is  té-to-váz-nak lép-te-i. Dicsőség_az_A-tyá-nak a Fi-ú-nak és a Szent-lé-lek-nek.}

% ============================================================================
% SANCTORUM MULIERUM
% ============================================================================

% ············································································
% Invitatorium
% ············································································

% Ut supra, in Communi sanctorum virorum

% ············································································
% Officium Lectionis
% ············································································

% Hynus (Pro una sancta muliere)
\defsheet{HaeFeminaeLaudabilis}
  {<-1-----3---eE2---4---ed1--2-4--3\\-:-3--1--2--3---2e--1---2--1í-,-1----3----eE2-4-----tT4--2-4--3\\-:-3---1--2---3---2e---1--2--1í-?,-2eE2-1\\q\\-.}
  {  Szent asz-szony-ról szól é-ne-künk, ki ér-de-mek-ben gaz-da-gon, mert tisz-ta, szent volt é-le-te,   már an-gya-lok közt ün-ne-pel.  Á----men.}

% ············································································
% Laudes
% ············································································

% Hymnus (Pro una sancta muliere)
\defsheet{NobilemChristiFamulam}
  {<X-5-----4---5---4-1í-:-4--4---3----4--5x-5x-,-5----7--8---7-5x-:-5----4--3-4--5x--5x-,-5-----4-5---4---1í----:-4--4--3--4-tT4-3y-,-3--3--2--1í-1í-?,-1wW1-0\\q\\-.}
  {   Krisz-tus-nak é-kes, ne-mes szol-gá-ló-ját  zeng-ze-tes é-nek  hang-ja-i da-lol-ják, szent e-rős asz-szonyt, ki-re az Í-rás ád   is-te-ni ál-dást. Á----men.}

% Responsorium breve (Pro una sancta, Extra tempore paschali)
\defsheet{AdiuvabitEam}
  {<X-3--3---3--4---3---4t-5--:-4--5--rR3-4--3---3---,-[--------------5z--4-:-5----5--5---4--3---4t--tT4-?,-[-------------5z--5-:-5-5--4-4-:-5--5-4-----3--4t--tT4-.}
  {   Is-ten ol-tal-maz-ta őt * És re-á   te-kin-tett. Nem_fog_megren-dül-ni, mert Is-ten la-kik ben-ne.    Dicsőség_az_A-tyá-nak a Fi-ú-nak és a Szent-lé-lek-nek.}

% Benedictus
\defsheet[tone={8szo-c}]{SimileEst}
  {<-4--5z--6--¨-7-7----5----6--5---4-,-4-4i-8---8--6--7--6---6--:-5-4--5-6---7----5----6--5--4---,-rF1-45z--6u--8--7--tT4-4--,-4--4i-8--7---5--uU6-:-4--5---5èuj5-zZ5-4-?,}
  {  Ha-son-ló * a meny-nyek or-szá-ga  a ke-res-ke-dő em-ber-hez, a-ki i-gaz-gyön-gyöt ke-re-sett; ta--lált egy ér-té-ke--set, el-ad-ta min-de-nét,  és meg-vet---te  azt.}

\defsheet{DeusInAduitorium}
  {<ô---£---------------------3r--4-,------£----------------3r-4-?,-£------------------------------------------3r--4--,-£---------------------------------------3r-4--,-4--5--rR3-3-.}
  {<V.> Istenem,_jöjj_segítsé-gem-re! <R.> Uram,_segíts_meg en-gem! Dicsőség_az_Atyának,_a_Fiúnak_és_a_Szentlé-lek-nek, miképpen_kezdetben,_most_és_mindörökké. Á--men. Al-le-lu--ja.}

\defsheet{BenedicamusDomino}
  {<----£---------------2--1w-2--,------£-----------2----1w-2-.}
  {<V.> Mondjunk_áldást az Úr-nak! <R.> Istennek_le-gyen há-la.}

% ============================================================================
% TONI INVITATORII
% ============================================================================

\defsheet{InvitatoriumMi}
  {<-5---4--5---:-[------------4--6--5--:-4--[----------------------4--4z-5--5---,-5---4----[---------4--3--4t--5--:-5--rR3--23r-eE2-1-:-3---3---4-5---rR3E2-2----?,-----5----4--[----------4----6--5--:-4--[--------------------4-----6--5---,-4-[-------------------------4---6---5--:-4--[----------------4---6-----5-,-5-4--3-4t--5--:-2e-4-eE2-1--:-1-¢-----------------------4--5--rR3E2-2e--?,-5---4--[------------------------4----6-5---5-:-4----[-----------------------4----6-5---,-4----5-5-4--3--4t-5---:-5--rR3-2e-4--eE2-1-:-1--3--4--5---rR3E2-2----?,----5---4----[------------------4-6---5--5-,--4--[-----------4----6---5--5-:-t4---[---------------------4---3-4----5--5--:-5-4---[--------4---3--2---3---4---eE2-1-:-1--¢--------------------------4---5---4---rR3E2-2e--?,-5----4---[---------------------4-6---5--5--:-5---4----5----[---------------4--6---5-,-5---4--5---5---5---4-3--4t-5--:-5-4----[----------4--3----2e--4--eE2-1--:-3---1---¢---------------4---5---rR3E2-2-----?,----5--4---[-----------------------------------4---6--5---5--:-3--4---[----------------rR3-2----3--4---eE2-1-:-1--¢-----------3-4---5---rR3E2-2e-?,-.}
  {  Jöj-je-tek,  örvendezzünk az Úr-nak, és ujjongjunk_üdvünk_szik-lá-ja e--lőtt! Lép-jünk színe_elé há-la-dal-lal, ma-gasz-tal-juk őt  han-gos é-nek-szó---val! ANT. Mert az Úr_valóban nagy Is-ten, az összes_isteneknél_di-csőbb ki-rály. A föld_mélységei_az_ő_kezé-ben van-nak,  és övé_a_hegyek_min-den or----ma. Ö-vé a ten-ger, az ő mű--ve,  a szárazföld_is_az_ő_keze al-ko-tá----sa. *  Jöj-je-tek,_boruljunk_le,_hódol-junk e-lőt-te, hull-junk_térdre_Urunk,_Alko-tónk e-lőtt, mert ő a mi Is-te-nünk, mi meg há-za né--pe  és ke-zé-nek nyá---ja. ANT. Bár-csak hallgatnátok_ma_az ő sza-vá-ra: „Ne legyetek_ke-mény szí-vű-ek, mint Meribánál,_Massza_nap-ján a pusz-tá-ban, a-hol atyáitok meg-kí-sér-tet-tek en-gem, és próbára_tettek,_bár_látták szá-mos cso-dá----mat. * Negy-ven évig_terhemre_volt_ez a nem-ze-dék, azt mond-tam: Ez_a_nép_tévely-gő szí-vű. Nem is-mer-ték meg u-ta-i--mat, e-zért haragomban ki-mond-tam es-kü--met: Nyu-gal-mam_helyére_nem jut-hat-nak   el!” ANT. Di-cső-ség_az_Atyának_és_Fiúnak_és_Szentlé-lek Is-ten-nek, mi-kép-pen_kezdetben_va-la, most és min-den-kor és mindörökkön ö-rök-ké! Á-----men.     * ANT.}

\defsheet{InvitatoriumFa}
  {<X-5---6--¦-----------------5--6--5--:-3r-[----------------------6--5--4-3---,-5---6----¦------------5--6---5--:-3--4----[----------------6-5---4---3------?,-5z---¦-------------5----6--5--:-3r-[--------------5---6--5-----4--3---,-5z-¦------------------------5---6---5--:-3r-[------------6---5---4--3-,-5-6--¦------------5--6-5-:-3r-[-----------------------4--5--6--5--?,-5---6--¦---------------------7---5----6-6---5-:-3----4----[----------------6--5----4-3---,-5z---7-7-5--6--6--5---:-3--4---[----------------6-5---4---3---?,---5---6----¦------------------5-6---6--5-:--3--4--5---5---6--5----4---3--3-,-5z---¦---------------------------------5--:-¦-------------------------5---6--5--:-3r-[--------------------------4---5---6---5--5---?,-5----6---¦---------------------5-6---6--5--:-3r--[--------------------6----5--4---3-,-5z--¦-----------------5--:-¦-----------------------5---6--6--5--:-3---4---[---------------6---5---4---3----?,---5--6---¦--------------5--6--5-5-:-3r-5-----5--6---5--4---3--,-5--6---¦------------------------5--6---6---5-:-3r-[-------------4---5---6-5----?,-.}
  {   Jöj-je-tek,_örvendezzünk az Úr-nak, és ujjongjunk_üdvünk_szik-lá-ja e-lőtt. Lép-jünk színe_elé_há-la-dal-lal, ma-gasz-taljuk_őt_hangos é-nek-szó-val! ANT. Mert az_Úr_valóban nagy Is-ten, az összes_istenek-nél di-csőbb ki-rály. A  föld_mélységei_az_ő_kezé-ben van-nak, és övé_a_hegyek min-den or-ma. Ö-vé a_tenger,_az ő mű-ve, a  szárazföld_is_az_ő_keze al-ko-tá-sa. * Jöj-je-tek,_boruljunk_le,_hó-dol-junk e-lőt-te, hull-junk térdre_Urunk,_Al-ko-tónk e-lőtt, mert ő a mi Is-te-nünk, mi meg háza_népe_és_kez-é-nek nyá-ja. ANT. Bár-csak hallgatnátok_ma_az ő sza-vá-ra: „Ne le-gye-tek ke-mény szí-vű-ek, mint Meribánál,_Massza_napján_a_pusztá-ban, ahol_atyáitok_megkísértet-tek en-gem, és próbára_tettek,_bár_látták szá-mos cso-dá-mat. * Negy-ven évig_terhemre_volt_ez a nem-ze-dék, azt mondtam:_Ez_a_nép_té-vely-gő szí-vű. Nem ismerték_meg_utai-mat, ezért_haragomban_kimond-tam es-kü-met: Nyu-gal-mam_helyére_nem jut-hat-nak el!” ANT. Di-cső-ség_az_Atyának és Fi-ú-nak és Szent-lé-lek Is-ten-nek, mi-kép-pen_kezdetben_vala,_most és min-den-kor és mindörökkön_ö-rök-ké! Á-men. * ANT.}

\defsheet{InvitatoriumSol}
  {<ô-5---4--5----[------------4--4t-5--:-4--[-----------------4----4z-6--5-5---,-5---4----[------------4--4t--5--:-4--[---------------------6-5---4---eE2---?,--5----4--[----------4----4t-5--:-4--[-----------------4--4z----5--5---,-5--4----[-------------------4---4t--5-:--4--[--------4----4z--6---5--5-,-5-4--[------------4-4t-5-:-4--[-----------------------6--5--4--eE2-?,-5---4--[------------------------4----4t-5---5-:-4----[------------------4--4z-6----5-5---,-5----4-5-4--4t-5--5---:-4--[--------------------6-5---4---eE2-?,---5---4----[------------------4-4t--5--5-:--4--5--5---5---4--4z---6---5--5-,-5----4--[-------------------------------4--:-4-[---------------------4---4t--5--5--:-4--[--------------------------6---5---4---4--eE2-?,-5----4---[---------------------4-4t--5--5--:-4---[-----------------4--4z---6--5---5-,-5---4--[---------------4--:-4-[----------------------4---4t-5--5--:-4---[------------------6---5---4---eE2--?,---5--4---[-----------------4--4t-5-:-5--5-----5--4z--6--5---5--,-5--4---5---5---4---4t--5--5-:-4----4z-6---5---5-,-4--[-------------6---5---4-eE2--,-.}
  {   Jöj-je-tek, örvendezzünk az Úr-nak, és ujjongjunk_üdvünk szik-lá-ja e-lőtt. Lép-jünk színe_elé_há-la-dal-lal, ma-gasz_taljuk_őt_hangos é-nek-szó-val! ANT. Mert az Úr_valóban nagy Is-ten, az összes_isteneknél di-csőbb ki-rály. A  föld mélységei_az_ő_kezé-ben van-nak, és övé_a_he-gyek min-den or-ma. Ö-vé a_tenger,_az ő mű-ve, a  szárazföld_is_az_ő_keze al-ko-tá-sa. *  Jöj-je-tek,_boruljunk_le,_hódol-junk e--lőt-te, hull-junk_térdre_Urunk, Al-ko-tónk e-lőtt, mert ő a mi Is-te-nünk, mi meg_háza_népe_és_kez-é-nek nyá-ja. ANT. Bár-csak hallgatnátok_ma_az ő sza-vá-ra: „Ne le-gye-tek ke-mény szí-vű-ek, mint Me-ribánál,_Massza_napján_a_pusztá-ban, a-hol_atyáitok_megkísér-tet-tek en-gem, és próbára_tettek,_bár_látták szá-mos cso-dá-mat. * Negy-ven évig_terhemre_volt_ez a nem-ze-dék, azt mondtam:_Ez_a_nép té-vely-gő szí-vű. Nem is-merték_meg_utai-mat, e-zért_haragomban_kimond-tam es-kü-met: Nyu-galmam_helyére_nem jut-hat-nak el!” ANT. Di-cső-ség_az_Atyának_és Fi-ú--nak és Szent-lé-lek Is-ten-nek, mi-kép-pen kez-det-ben va-la, most és min-den-kor és mindörökkön_ö-rök-ké! Á-men. * ANT.}

% ============================================================================
% ALLELUIA (IN TONIS DIVERSIS)
% ============================================================================

\defsheet[tone={1re-a}]{AlleluiaDorian}
  {<X-3--3--4t-5---:-5--4--tT4-3r--4--eE2-1--1-?,}
  {   Al-le-lu-ja, * al-le-lu--ja, al-le--lu-ja!}

\defsheet[tone={2-g}]{AlleluiaHypodorian}
  {<ô-2--2--3--wW1-2--rR3-2--ñ---:-3--4--3--2-?,}
  {   Al-le-lu-ja, al-le--lu-ja, * al-le-lu-ja!}

\defsheet[tone={3-c}]{AlleluiaPhrygian}
  {<-4--5--7--7---:-uU6-5--4t-4-:-rR3-4tT4-2--2-?,}
  {  Al-le-lu-ja, * al--le-lu-ja, al--le---lu-ja!}

\defsheet[tone={4-a}]{AlleluiaHypophrygian}
  {<-2--1--2--1---:-0q-eE2-1w-2-?,}
  {  Al-le-lu-ja, * al-le--lu-ja!}

\defsheet[tone={5-c}]{AlleluiaLydian}
  {<-3--4--4--3---:-5--7--7i-uU6-:-5--4t-3--3-?,}
  {  Al-le-lu-ja, * al-le-lu-ja,   al-le-lu-ja!}

\defsheet[tone={6-a}]{AlleluiaHypolydian}
  {<X-3--4--5--3rf2ed1-:-3r-zZ5-rR3-3-?,}
  {   Al-le-lu-ja,     * al-le--lu--ja!}

\defsheet[tone={7do-d}]{AlleluiaMixolydian}
  {<-6--7--8o-8---:-iI7-5u-zZ5-4-:-5-5uU6-4--4-?,}
  {  Al-le-lu-ja, * al--le-lu--ja, al-le--lu-ja!}

\defsheet[tone={8szo-c}]{AlleluiaHypomixolydian}
  {<-4--4--5--rR3-:-5--7--6u-tT4R3-4t-5--4--4-?,}
  {  Al-le-lu-ja, * al-le-lu-ja,   al-le-lu-ja!}

\defsheet[tone={tper-g}]{TonusPeregrinus}
  {<X-0--1--3r-4---:-5--6--5--rR3-4t-5--4y-4y-?,}
  {   Al-le-lu-ja, * al-le-lu-ja, al-le-lu-ja!}

% ============================================================================
% TONUS DE CANTU LECTIONUM
% ============================================================================

%\defsheet{TonusDeCantuLectionum}
%  {<-------¦------------------------------7--5-.}
%  {<flexa> Boldog_az_ember,_ha_fenyíti_az Is-ten; ne vesd meg azért a Mindenható feddő szavát! Ha sebet üt rajtad, majd be is kötözi, ha szétzúz is, meggyógyít a keze. Ki látta át Isten gondolatait?}

% ============================================================================
% TE DEUM
% ============================================================================

\defsheet{TeDeum}
  {<-2--4t---5--5---5--56u-zZ5---:-2--4t--5--5z--4--5z-uU6Z5-4--?,-7--7----7-5---zZ5T4¨5z-5--:-2---5---5---5z-4--5z--uU6Z5-4--?,-¦--------------5-zZ5T4-5zZ5-5-:-2--5---5---5z-4--5z-uU6Z5-4-?,-¦------------6---5--uU6Z5T4-5zZ5-5-:-2---5---5--5z--4-5z-uU6Z5-4-?,-2tT4--5z-5-?,-2tT4--6-?,-8-----5--6--5z---4--5z--uU6Z5-4-?,--¦---------------5----zZ5T4¨5z-5-:--2---5--5z--4--5z--uU6-5---4-?,-7--7---7-5---zZ5T4¨5z-5-:-5---5z-4---5z--uU6-5----4-?,-7--7---7-----5---zZ5T4¨5z-5-:-5--5z--4---5z--uU6-5----4-?,-¦------------5----zZ5T4-5zZ5-5-:--2---5-5---5z--uU6-5---4-?,-¦-----------5---zZ5T4¨5z-5-:--5-5---5z----4---5z--uU6Z5-4-?,--2tT4-5z-5-:-7--7---7---uU6-4-5z-uU6Z5-4-?,-¦------------5--zZ5T4¨5z-5--:-2-5z---4---5z--uU6-5-4-?,-2tT4--5z--5--:-7---7---4---5z-4----3--2-?,-5---5-----5----4----5--6---4--2--?,-[---------56u-zZ5-:-2----5-5---4--5z--4--2-?,-[----------------------6---7---6---5-:-2--[------------4--5---6---4-3--2-?,-[----------6----7---6----5-:-2---[------------------4--5----6---4---2-?,-[-------------6---7----6--5-:-2-5---5---4--6---4---3---2-?,-7--5-:---6--5--5---4--6--4--2-?,-[--------------6----7--6--5---:-2----[--------------¨:-[--------------4---5---6---4---2-?,-3----2---3--1-:-2-3---£-----------5---3---rR3E2-2-?,--0--1e-¢------------------1-3---wW1Q0-0--:-£------------------3--4t-rR3E2-2-?,-0--1e----3---3---3--:-¢--------------------2--1---4---5-?,-5---7i--8--8---8----89p-oO8-:-8--7--8----9--7---7i-?,-5---7i-©---------------9-pP9-8-:-8----8-8---7--9--7----7i-?,-5---7i----©--------------9---ö-9---8-:-8--8--7---9---7---7i-?,-5----7i-©-------------------9--ö--9---8--:-8----7--8---9---7--7i-?,-5--7i---8--88(---8--8---8---8--9---ö--9---8---:-5----©-----------7----9--7---5-?,-X6---6--4---6--6-----7i-uU6Z5-5--:-3-X4z---6-4---6--7iI7-zZ5T4-5-?,-X4uU667iI7U6Z5T4-5-.}
  {  Is-ten, té-ged di-csé-rünk, * té-ged Úr-nak mi is val---lunk. Té-ged, ö-rök Is-------ten, min-den föl-di ál-lat di----csér. Tenéked_mennye-i an----gya--lok és men-nye-i  fe-je-del---mek, tenéked_kéru-bim és szé-----ra---fim min-den-ko-ron é-ne-kel---nek: szent Is-ten! Szent Úr!  Szent Úr Is-ten, di-cső Ki----rály! Teljes_menynyor-szág és       föld fel-sé-ges di-cső-sé--ged-del. Té-ged a-pos-to-------lok tár-sa-ság-gal di--csér-nek. Té-ged szent pró-fé-------ták vi-gas-ság-gal di--csér-nek. Téged_kínval-lott már---tí---rok, kik é-ret-ted meg-hol-tak, téged_ez_vi-lág sze------rint A-nya-szent-egy-ház di----csér, fel--sé-ges ma-gas-ság-nak ő Is-te----nét, mi_tisztelen-dő U--------runk e-gyet-len egy Fi--a-dat, Szent-lel-ket, min-ket meg-vi-gasz-ta-lót. Te, Krisz-tus, vagy di-cső ki-rály. Te_Atyais-ten-nek   vagy ö-rök-sé-ges Fi-a.   Te_megváltásra_embersé-get fel-ven-ni  Má-riának_szent mé-hét nem u-tá-lád. Te_örök_ha-lált meg-győz-vén men-nyei_szent_kapukat ne-künk meg-nyi-tál. Te_ülsz_Isten-nek jobb-já-ra, A-tyá-nak di-cső-sé-gé-ben.   Hi-szünk el-jö-ven-dő bí-ró-nak. Téged_azért,_U-ram, mi ké-rünk, hogy minket_segélj, @  kiket_szent_vé-red-del meg-vál-tál. Szen-tid kö-zé  ö-rök dicsőségben min-ket szám--lálj. Üd-vö-zítsed_te_népedet, U-ram Is----ten, és_megáldjad_te_ma-ra-dé-ki----dat. És tarts meg min-ket, és_elfelvígy_az_örök di-cső-ség-be!  Min-den na-pon-ként té--ged,  Úr Is-ten, di-csé-rünk. És, U--ram,_te_nevedet ö-rök-ké, mind-ö-rök-ké di-csér-jük.  Mél-tólj, Uram_Isten_min-ket ő-riz-ned ez na-pon bűn nél-kül!  Légy ir-galmas,_Uram_Isten, bű-nö-sök-nek, légy ir-gal-mas ne-künk! Le-gyen te szent ir-gal-mas-sá-god mi-hoz-zánk, mert minekünk_te vagy ol-tal-munk. Te-ben-ned bí-zunk, Úr Is----ten, a--zért ö-rök-ké mi   é-----lünk. Á---------------men.}

% ============================================================================
% MELODIAE HYMNORUM
% ============================================================================

% Hymnus
\defsheet{NepekMegvalto}
  {<ô-2--2--2---1---2---4---3--2í--:-2--3----4--5------4--5---5--6x--,-6----6----5---6--4----5--rR3-2í-:-2--1----2--4--5----2-1--2í-?,-1e4R3E2W1-2í-.}
  {   El-ső nap min-den nap kö-zött, me-lyen ég s_föld al-kot-ta-tott, s_me-lyen fel-tá-madt Al-ko--tónk le-győz-te ér-tünk a ha-lált, Á---------men.}

% Hymnus
\defsheet{IttAzAlkalmas}
  {<ô-4--4----4---4---4---4--3--2í-:-4--4-----4-4-----1--2---1--0í-,-2--2--3---4---3----2---1--2eE2í-:-2---2---1--0--1--2---1--1í-?,-1wW1-0\\q\\-.}
  {   Fé-nyes-ség-osz-tó, bő-ke-zű,  ki fényt á-raszt-va gaz-da-gon  az éj-nek gyá-szát osz-la-tod,    s_a nap vi-lá-ga tűz le ránk. Á----men.}

% Hymnus
\defsheet{ANapvialg}
  {<-5--5--5--5--1---3--2--1í--:-5--5--6---4----5-----7----6--5x-,-5-----6---7--8----7----5---7--zZ5T4x-:-5---5--6---4--2---4---5--5x-?,-5zZ5-4\\t\\-.}
  {  Ég Al-ko-tó-ja, Is-te-nünk, ki éj-jel hold-fényt adsz ne-künk s_nap-pal na-pot, mely kör-be fut,     s_e-lő-írt út-ján cél-ba jut.  Á----men.}

% Hymnus
\defsheet{IttJelenVagyon}
  {<X-5--4----5--4---1í-:-4--4--3---4--5x-5x-,-5---7---8-7---5x-:-5--4---3---4--5x---5x-,-5---4---5----4---1í-:-4-4--3---4-tT4-3y-,-3---3---2-1í---1í-?,-1wW1-0\\q\\-.}
  {   Sá-pad, el-osz-lik  kö-de már az éj-nek, pir-kad a haj-nal, ró-zsa-szí-nű fény kél; for-dul-junk buz-gón  a vi-lág U-rá--hoz  han-gos i-mánk-kal,  Á----men.}

% Hymnus
\defsheet{IttAzAlkalmasSzentIdo}
  {<-7---7---7-7----7---7----6--5-:-7-7----7--7---4---5--4-3--,-5---5---6--7---6--5----4--5zZ5-:-5--5----4--3--4---5-4--4-?,--4tT4-3\\r\\-.}
  {  Már kél a fény-nek csil-la-ga: e-seng-ve kér-jük az U-rat, jár-jon ma min-de-nütt ve-lünk,  ne ront-sa ár-tás é-le-tünk. Á----men.}

% Hymnus
\defsheet{KiMindenFenynek}
  {<-3--3---4---5---4---4-----3--2---:-4--2-----3----1---0--1----3--3--,-3----3---4--5----4---4--3-2---:-4--2---3--1---0---3----1--1-?,-02eE2W1Q0-1-.}
  {  Üd-vös-ség nap-ja, Krisz-tu-sunk, ha-sítsd szét vak-sá-gunk kö-dét, hogy fel-ra-gyog-jon az e-rény, mi-dőn re-ánk nap-palt ho-zol. Á---------men.}

% Hymnus
\defsheet{IstenTeMindentAlkoto}
  {<-5--5---5--5---5--5--4--4-,-5--5--5----4-3---4--5--5-,--5-5---5-4--3---4---3--2----,-3-------4--4--3--4---5-4--4-?,-4tT4-3\\r\\-.}
  {  Vi-lág te-rem-tő Is-te-ne, ki Úr vagy é-jen és na-pon, i-dőt i-dő-vel vál-to-gatsz, s_nincs mű-ve-id-ben u-na-lom, Á----men.}

% ============================================================================
% DEFTONES
% ============================================================================

\deftone{1la-a}{<-[-4-3-rR3-1-.}{1. lá tónus}
\deftone{1re-a}{<-[-4-3-4t-4-.}{1. ré tónus}
\deftone{1re-b}{<-¥-5-4-5z-5-.}{1. ré tónus}

\deftone{2-f}{<-¢-0q-1-.}{2. tónus}
\deftone{2-g}{<-£-1w-2-.}{2. tónus}
\deftone{2-b}{<-¥-3r-4-.}{2. tónus}
\deftone{2-c}{<-¦-4t-5-.}{2. tónus}
\deftone{2-d}{<-©-5z-6-.}{2. tónus}

\deftone{3-b}{<-¥-4z-4-3-.}{3. tónus}
\deftone{3-c}{<-¦-5u-5-4-.}{3. tónus}

\deftone{4-g}{<-£-3-4-3-1-.}{4. tónus}
\deftone{4-a}{<-[-4-5-4-2-.}{4. tónus}

\deftone{5-c}{<-¦-8-6-7-5-.}{5. tónus}
\deftone{5-g}{<-£-5-3-4-2-.}{5. tónus}

\deftone{6-a}{<-[-4-3-4-3-.}{6. tónus}

\deftone{7do-a}{<-[-6-5-rR3-4-.}{7. dó tónus}
\deftone{7do-d}{<-©-9-8-uU6-7-.}{7. dó tónus}
\deftone{7la-a}{<-[-6-5-rR3-2-.}{7. lá tónus}
\deftone{7la-d}{<-©-9-8-uU6-5-.}{7. lá tónus}

\deftone{8szo-g}{<-£-3-4-2-1-.}{8. szó tónus}
\deftone{8szo-b}{<-¥-5-6-4-3-.}{8. szó tónus}
\deftone{8szo-c}{<-¦-6-7-5-4-.}{8. szó tónus}
\deftone{8do-c}{<-¦-5-7-8-7-.}{8. dó tónus}

\deftone{8szo-a-c}{<-2e-[-…-[-4-5-6-‡-5-,-[-4-5-3-„-2-.}{8. szó tónus}

\deftone{tirr-e}{<-¡-3-1-3-2-.}{ton.irreg.}
\deftone{tirr-a}{<-[-6-4-6-5-.}{ton.irreg.}
\deftone{tirr-b}{<-¥-7-5-7-6-.}{ton.irreg.}

\deftone{tper-g}{<-£-5-1-eE2-1-.}{ton.per.}

\title{Tábori zsolozsmáskönyv}
\subtitle{902. Kucsera Ferenc cserkészcsapat - 2026. nyári tábor}

\begin{document}

\maketitle

\section*{Állandó részek}

\textcolor{red}{(}Nyisd meg, Uram, ajkamat neved dicséretére; tisztítsd meg szívemet minden
hiába\-való és vétkes gondolattól: világosítsd meg értelmemet, gyújtsd lángra szívemet, hogy ezt a
zsolozsmát méltón, figyelmesen, áhitattal végezzem, és meghallgatásra méltó legyek isteni fölséged
színe előtt. Krisztus, a mi Urunk által. \textcolor{red}{)}

\subsubsection*{Imádságra hívás}

\versicle{Nyisd meg Uram, ajkam\underline{a}t,}{Hogy dicséretedet hirdesse szav\underline{a}m!}

\redtext{4. tónus:}

\printsheet{InvitatoriumMi}

\redtext{6. tónus:}

\printsheet{InvitatoriumFa}

\subsubsection*{Olvasmányos imaóra}

\parttitle{Imaóra kezdete}

\printsheet{DeusInAduitorium}

\rubric{A fentiek elmaradnak, ha az Imádságra hívás az Olvasmányos imaórát előzi meg.}

\parttitle{Himnusz}

\parttitle{Zsoltározás}

\parttitle{Átvezető vers}

\parttitle{Olvasmányok}

\parttitle{Himnusz, Te Deum}

\rubric{Vasárnapokon, főünnepeken és ünnepeken a Második olvasmány és Válaszos éneke után imádkozzuk
a következő hinnuszt:}

\label{TeDeum}

\printsheet{TeDeum}

\parttitle{Könyörgés}

\subsubsection*{Reggeli dicséret}

\parttitle{Imaóra kezdete}

\rubric{Mint az Olvasmányos imaórában.}

\parttitle{Imaóra befejezése}

\printsheet{BenedicamusDomino}

\subsubsection*{Befejező imaóra}

\parttitle{Imaóra kezdete}

\rubric{Mint az Olvasmányos imaórában.}

\parttitle{Lelkiismeret-vizsgálat}

\rubric{Dicséretes dolog most a lelkiismeret-vizsgálat, amely közös végzés esetén a mise bűnbánati
formája szerint történhet.}

\parttitle{Himnusz}

\printsheet{TeLucisAnteTerminum}

\begin{hymnstanzas}
    Álmodjék rólad a szívünk, \\
    álmunkban is te légy velünk, \\
    téged dicsérjen énekünk, \\
    midőn új napra ébredünk!
\stanzabreak
    Adj nekünk üdvös életet, \\
    szítsd fel a szív fáradt tüzét, \\
    a ránk törő rút éjhomályt \\
    világosságod rontsa szét!
\end{hymnstanzas}
\centerstanza[0.5]{%
    Kérünk, mindenható Atyánk, \\
    Úr Jézus Krisztus érdemén, \\
    ki Szentlélekkel és veled \\
    uralkodik, s örökre él. Ámen.
}

\rubric{Vagy azonos dallamon:}

\begin{hymnstanzas}
    Krisztus, tündöklő nappalunk,\\
    az éj árnyát eloszlatod,\\
    kit fényből fénynek vall a hit\\
    és fénnyel áldod szentjeid.
\stanzabreak
    Esengünk hozzád, szent Urunk, \\
    míg tart az éj, vigyázz reánk, \\
    találjunk benned pihenést, \\
    adj nyugodalmas éjszakát.
\stanzabreak
    Mikor lecsukjuk két szemünk, \\
    a szívünk virrasszon veled, \\
    őrködjék mindig jobb kezed \\
    hűséges híveid felett.
\stanzabreak
    Védelmezőnk, tekints le ránk, \\
    igázd le ellenségeink, \\
    vezesd szolgáidat, kikért \\
    véred váltsága volt a bér.
\end{hymnstanzas}
\centerstanza{
    Legyen most néked, Krisztusunk, \\
    és szent Atyádnak tisztelet, \\
    s a Szentléleknek is veled \\
    zengjen dicséret szüntelen. Ámen.
}

\parttitle{Zsoltározás}

\parttitle{Rövid olvasmány}

\parttitle{Rövid válaszos ének}

\printsheet{InManusTuasDomine}

\parttitle[Lk 2,29-32]{Evangéliumi kantikum}

\printsheet{SalvaNos}

\parttitle{Könyörgés}

\parttitle{Imaóra befejezése}

\versicle{A nyugodalmas éjszakát és a jó halál kegyelmét adja meg nekünk a mindenható Isten!}{Ámen.}

\parttitle{Szűz Mária záróantifónája}

\printsheet{SalveRegina}

\section*{Változó részek}

\subsection*{Évközi 14. vasárnap - Előtábor 2. nap}

\label{D/Compl1}

\subsubsection*{Befejező imaóra az I. Esti dicséret után}

\parttitle{Zsoltározás}

\redtext{1. antifóna}

\printsheet{MiserereMihiDomine}

\begin{psalmsheet}[tone=8, continuous, number={4. zsoltár}, title={Hálaadás}, motto={Az Úr csodálatra méltóvá tette azt, akit feltámasztott a holtak közül (Szt. Ágoston).}]
  \psalmtext{
    1. Halld meg kiáltásomat, igazságos Istenem, * szorult helyzetemben segíts meg engem. \\
    2. Könyörülj rajtam, * és hallgasd meg imámat! ---

    3. Meddig lesz még kemény a szívetek, emberek? * Miért szeretitek a haszontalanságot, meddig törekedtek még hazugságra? \\
    4. Tudjátok meg: az Úr csodásan megmenti hívét; * meghallgat az Úr, ha hozzá kiáltok. ---

    5. Bár haragusztok, többé ne vétkezzetek! † Szálljatok magatokba szívből, * nyugovóra térve teremtsetek csendet. \\
    6. Igaz áldozatot hozzatok, * és bízzatok az Úrban! ---

    7. Azt mondják sokan: „Ki részesít jóban minket?” * Ragyogtasd ránk, Uram, arcod fényességét! \\
    8. Szívembe nagyobb örömet öntöttél, * mint azoknak, akik dúskálnak búzában meg borban! \\
    9. Alig térek nyugovóra, békében elalszom, * mert biztonságot egyedül te adsz nekem, Uram.
}
\end{psalmsheet}

\redtext{2. antifóna}

\printsheet{InNoctibus}

\begin{psalmsheet}[tone=4, continuous, number={133. zsoltár}, title={Esti ima a templomban}, motto={Dicsérjétek Istenünket, ti, összes szolgái, és akik félve tisztelitek őt, kicsinyek és nagyok (Jel 19, 5).}]
  \psalmtext{
    1. Áldjátok az Urat mindnyájan, ti, az Úr szolgái, * akik éjjelente az Úr házában álltok. \\
    2. Tárjátok ki kezeteket szentélye felé, * és áldjátok az Urat! ---

    3. Áldjon meg téged az Úr Sionból, * aki az eget és a földet alkotta!
}
\end{psalmsheet}

\parttitle[MTörv 6, 4-7]{Rövid olvasmány}
Halld, Izrael! Az Úr, a mi Istenünk az egyetlen \underline{Ú}r! † Szeresd Uradat, Istenedet szíved, lelked mélyéből, m\underline{i}nden erőddel! * Ezeket a parancsokat, amelyeket ma szabok neked, őrizd meg szívedben és vésd gyermekeidnek is az eszébe; beszélj róluk, amikor otthon tartózkodsz és amikor úton vagy, amikor lefekszel s amik\underline{o}r f\underline{ö}lkelsz!

\parttitle{Könyörgés}
Látogass meg minket, Istenünk, ezen az éjszakán, hogy holnap reggel tőled erőt nyerve ébredjünk, és
Krisztus feltámadásának örvendhessünk. Aki él és uralkodik mindörökkön-örökké.

\subsubsection*{Imádságra hívás}

\label{D/Inv}

\printsheet{VeniteExsultemusDomino}

\subsubsection*{Olvasmányos imaóra}

\parttitle{Himnusz}

\printsheet{MediaeNoctis}

\begin{hymnstanzas}
    s a Lelket, aki egy velük\\
    és egyesíti lényegük:\\
    a Hármas-Egy élő Csodát,\\
    kit énekünk örökre áld.
\stanzabreak
    Félelmes éji óra ez,\\
    öldöklő angyal jár körül,\\
    Egyiptomnak halála lett:\\
    elsőszülöttek végzete.
\stanzabreak
    Az igazaknak üdvöt ad:\\
    az angyal büntető keze\\
    nem mert reájuk sújtani,\\
    a vér jelétől megriadt.
\stanzabreak
    Egyiptom sírt keservesen\\
    tömérdek bú és gyász között.\\
    De vígan ujjong Izrael:\\
    a bárány vére őrzi őt.
\stanzabreak
    Urunk, az igaz Izrael\\
    mi vagyunk: benned örülünk.\\
    Az ellenség nem érhet el,\\
    megóvja véred életünk.
\stanzabreak
    Országodhoz, áldott Király,\\
    ha eljössz, méltónak találj,\\
    hadd zengjünk akkor majd neked\\
    meg nem szűnő dicséretet. Ámen.
\end{hymnstanzas}

\parttitle{Zsoltározás}

\redtext{1. antifóna}

\printsheet{Confessionem}

\begin{psalmsheet}[tone=3, continuous, number={103 zsoltár}, title={Himnusz a teremtő Istenhez}, motto={Mindenki, aki Krisztusban van, új teremtmény. A régi megszűnt, és valami új valósult meg (2 Kor 5, 17).}, numeral={I.}]
    \psalmtext{
        1. Mondj áldást, lelkem, az Úrnak: * csodálatos a te nagyságod, Uram, Istenem!\\
        2. Fölséget és pompát öltöttél magadra, * a fényesség a te palástod. ---

        3. Kifeszíted az eget, mint egy sátort, * vizek fölött áll a palotád.\\
        3. Fölszállsz a felhők fogatára, * és a szelek szárnyán utazol.\\
        4. Követségbe küldöd a szélvihart, * és teneked szolgál a lobogó villám. ---

        5. Szilárd alapokra ágyaztad a földet, * nem fog az sohasem meginogni.\\
        6. Az őstenger köntösként takarta, * vizek borították még a hegyeket is. ---

        7. Ám feddő szavadra megfutottak, * félve menekültek mennydörgő hangod elől.\\
        8. Magasba szöktek a hegyek, behorpadtak a völgyek * ott, ahol kijelölted helyüket.\\
        9. Határt vontál eléjük, és át nem lépik azt, * nem nyelik el újra a föld színét.

        10. Patakokba gyűjtöd a forrásvizeket, * ott csörgedeznek a hegyek között.\\
        11. Belőlük isznak a mező vadjai, * szomjukat oltják az állatok.\\
        12. Partjukon fészkelnek a dalos madarak, * az ágak közt énekük hangzik.
}
\end{psalmsheet}

\redtext{2. antifóna}

\printsheet{DomineDeusMeus}

\begin{psalmsheet}[tone=8, continuous, numeral={II.}]
    \psalmtext{
        13. A hegyeket égi palotádból öntözöd, * alkotásod gyümölcsével jól tartod a földet.\\
        14. Füvet sarjasztasz az állatoknak, * és növényeket az ember hasznára,\\
        15. hogy kenyeret teremjen a föld, * és bort, hadd vigadjon tőle az emberi szív.\\
        15. Olajjal örömet derítesz az arcra, * a kenyér meg erőt önt az ember szívébe. ---

        16. Teleisszák magukat az Isten erdői * és Libanon cédrusai, melyeket ő ültetett.\\
        17. A madarak fészküket odarakják, * tetejükben tanyázik a gólya.\\
        18. A magas hegy a zergék hazája, * a sziklák meg a borznak kínálnak menedéket. ---

        19. Holdat készítettél, hogy mérje az időt, * nyugtának idejét jól tudja a nap.\\
        20. Kibocsátod az árnyakat, leszáll az éj: * benne az erdő vadjai zsákmány után osonnak.\\
        21. Az oroszlánkölykök prédáért üvöltenek, * és betevő falatot kérnek Istentől.\\
        22. Amikor a nap fölkel, visszatérnek, * és elhevernek vackaikon.\\
        23. Az ember meg munka után lát, * és dolgozik, míg nem jön az este.
}
\end{psalmsheet}

\redtext{3. antifóna}

\printsheet{QuamMagnificataSuntE}

\begin{psalmsheet}[tone=tonus-irregularis, continuous, numeral={III.}]
    \psalmtext{
        24. Mennyi mindent alkottál, Uram! † Bölcsességgel rendeztél el mindent, * betöltik a földet teremtményeid.\\
        25. Ott van a nagy és széles tenger, † nyüzsög benne a számtalan tengeri állat: * parányi életek és hüllőóriások.\\
        26. Hajók siklanak rajta mindenfelé, * s a Leviatán, amelyet azért alkottál, hogy játszadozzál vele. ---

        27. Tőled várják mindahányan, * hogy enniük adj kellő időben.\\
        28. Hogyha te adsz nekik, összeszedik, * ha megnyitod kezedet, eltelnek jókkal.\\
        29. De ha elfordítod arcod, megrettennek, † leheletüket ha megvonod, elhullanak, * és visszatérnek a porba.\\
        30. Ám, ha újra kiárasztod a lelked, életre kelnek, * s így megújítod a föld színét. ---

        31. Legyen örök az Úr dicsősége, * lelje örömét műveiben!\\
        32. Rátekint a földre, és az beléremeg, * megérinti a hegyeket, és füstölni kezdenek.\\
        33. Az Úrnak dalolok egész életemben, * zsoltárt zengek Istennek, amíg csak élek.\\
        34. Bárcsak fogadná tetszéssel dalomat, * hisz gyönyörűségemet őbenne találom.\\
        35. A bűnösök vesszenek a földről, † istentelenek ne legyenek többé! * Mondj áldást, lelkem, az Úrnak!
}
\end{psalmsheet}

\versicle{A ti szemetek boldog, mert l\underline{á}t.}{A fületek is az, mert h\underline{a}ll.}

\parttitle[12, 1-25]{Első olvasmány}
Sámuel második könyvéből

\begin{lectio}{Dávid bűnbánata}
\noindent Azokban a napokban az Úr elküldte Nátán prófétát Dávidhoz. El is ment hozzá, és így
beszélt: „Egy városban élt két ember. Az egyik gazdag volt, a másik szegény. A gazdagnak rengeteg
juha, kecskéje és barma volt. A szegénynek azonban nem volt egyebe, csak egy kisbáránya, vette
magának. Etette, úgyhogy megnőtt, vele magával és gyermekeivel. Evett a kenyeréből, és ivott a
poharából. Az ölében aludt, s olyan volt, mintha a lánya lett volna. Egyszer vendég érkezett a
gazdag emberhez. De nem tudta rászánni magát, hogy a juhai vagy barmai közül egyet is elvegyen, és a
vendégnek, aki érkezett, elkészítse. Ezért elvette a szegény ember báránykáját, és azt készítette el
vendégének.” Dávid nagyon megharagudott az emberre, és azt mondta Nátánnak: „Amint igaz, hogy az Úr
él: az az ember, aki ilyet tesz, az a halál fia. A bárányt négyszeresen vissza kell fizetnie, amiért
így tett, és nem volt kímélettel.”

Erre Nátán így szólt Dávidhoz: „Magad vagy ez az ember! Ezt mondja hát az Úr, Izrael Istene:
Fölkentelek Izrael királyává, és kiszabadítottalak Saul kezéből. Neked adtam urad házát és kebledre
urad asszonyait, rád bíztam Izrael házát és Júda házát. S ha ez kevés lett volna, még megtetéztem
volna ezzel meg azzal. Miért vetetted hát meg az Urat, s tettél olyat, ami gonosznak számít az Úr
szemében? Kard által elveszítetted a hetita Úriját, hogy feleségül vedd a feleségét. Igen,
elveszítetted az ammoniták kardja által! Nos, a kard sose fordul el házadtól, amiért megvetettél, és
elvetted a hetita Úrijától a feleségét, hogy a te asszonyod legyen. Ezt mondja az Úr: Nos, a saját
házadból hozok rád romlást. Elveszem a szemed láttára asszonyaidat, és odaadom őket más valakinek,
hogy napvilágnál asszonyaiddal háljon. Te titokban tetted, de én egész Izrael szeme láttára,
napvilágnál viszem ezt végbe!”

Erre Dávid így szólt Nátánhoz: „Vétkeztem az Úr ellen!” Nátán ezt válaszolta Dávidnak: „Így az Úr is
megbocsátja bűnödet, s nem halsz meg. De mert ezzel a tetteddel káromoltad az Urat, a fiú, aki
született neked, meghal.” Ezzel Nátán elment haza.

Az Úr megverte a gyermeket, akit Úrija felesége szült Dávidnak, úgyhogy súlyosan megbetegedett.
Dávid az Úrhoz fordult a fiú miatt; szigorú böjtöt tartott, s amikor hazament, éjszaka a puszta
földön aludt zsákkal takarózva. Udvarából a főemberek eléje járultak, és igyekeztek rábeszélni, hogy
keljen föl a földről, de nem akart, s nem is evett velük. A hetedik nap a gyermek meghalt. Dávid
udvari emberei nem merték neki jelenteni, hogy a gyerek meghalt. Azt mondták magukban: „Már akkor
sem hallgatott ránk, amikor még élt a gyerek, s a lelkére beszéltünk. Hát hogy mondhatnánk akkor meg
neki: a fiú halott? Még kárt tenne magában!” De Dávid észrevette, hogy udvari emberei suttogtak
egymás közt. Ebből megértette, hogy a gyerek meghalt. Ezért Dávid megkérdezte udvari embereitől:
„Meghalt a gyerek?” „Meg” – felelték.

Erre Dávid fölkelt a földről, megmosdott, bekente magát olajjal, és más ruhát öltött. Aztán elment
az Úr házába, és leborult. Ezután ételt tettek eléje, és evett. Erre udvari emberei így szóltak
hozzá: „Mit csinálsz? Amikor még élt a gyerek, böjtöltél, és siránkoztál. Most meg, hogy a gyerek
meghalt, fölkelsz, és ételt veszel magadhoz.” Ezt válaszolta: „Amíg a gyermek még élt, böjtöltem és
sírtam, mert azt gondoltam, ki tudja, hátha megkönyörül rajtam az Úr, és életben marad a gyerek. De
most, hogy meghalt, minek böjtöljek tovább? Hát visszahozhatom? Én utána megyek, de ő nem tér hozzám
vissza.”

Dávid Batsebát is vigasztalta. Elment hozzá, és együtt hált vele. (Batseba) fogant, és fiút szült. A
Salamon nevet adta neki. Az Úr kedvét találta benne, s ezt tudtul is adta Nátán próféta által, aki
az Úr szavára Jedidjának nevezte el.
\end{lectio}

\parttitle[Manassze imája 9. 10. 12; Zsolt 50, 5. 6]{Válaszos ének}

\versicle{%
    Megsokasodtak vétkeim, és bűneimnek száma, mint a teng\underline{e}r hom\underline{o}kja; †
    bűneim sokasága miatt nem vagyok méltó, hogy lássam a magas eget, mert haragra
    inger\underline{e}lt\underline{e}lek, * És ami színed előtt gonosz, \underline{o}lyat
    t\underline{e}ttem.
}{%
    Gonoszságomat bel\underline{á}tom, † bűnöm előttem lebeg szünt\underline{e}len, * egyedül csak
    ellened v\underline{é}tk\underline{e}ztem.
}

\parttitle{Második olvasmány}
Szent Ágoston beszédeiből

{\centering\redtext{(Serm. 19, 2-3: CCL 41, 252-254)}}

\begin{lectio}{Istennek tetsző áldozat a megtört szív}
\noindent \textit{Gonoszságomat belátom} (Zsolt 50, 5), mondja Dávid. Ha én belátom, akkor te ne
ródd föl azt nekem, Istenem. Még ha becsületesen éltünk is, ne gondoljuk, hogy nincs bűnünk. Életünk
annyira lesz dicséretes, amennyire Isten bocsánatát kérjük. Kétségbeejtő az olyan ember, aki nem
annyira saját bűneit nézi, de annál inkább vájkál a másokéban. Szimatolja, de nem azért, hogy
javítson rajta, hanem mert marni akarja. Mivel magát nem tudja tisztára mosni, ezért inkább másokat
igyekszik bemocskolni. Mi ne így tegyünk, amikor Dávid példájára imádkozunk, és elégtételt akarunk
nyújtani Istennek. Dávid azt mondta: \textit{Gonoszságomat belátom, bűnöm előttem lebeg szüntelen}
(Zsolt 50, 5). Nem figyelt ő más bűnére, hanem saját magát vádolta. Nemcsak a felszínen
tapogatódzott, hanem mélyen magába szállt. Saját magának nem kegyelmezett, hogy a maga számára ne
meggondolatlanul kérje a megkegyelmezést.

Ki akarsz engesztelődni Istennel? Gondolkozz: mit kell tenned, hogy Isten megbékéljen veled. Figyeld
meg, mit olvasunk ugyanebben a zsoltárban: \textit{Te a véres áldozatot nem kedveled,
égőáldozataimban nem leled kedved} (Zsolt 50, 18). Hát akkor ne is legyen már áldozatod? Ne ajánlj
fel neki semmit? Hát semmiféle áldozattal sem lehet Istent kiengesztelned? Hogy is mondtad a
zsoltárban? Te a véres áldozatot nem kedveled, égőáldozataimban nem leled kedved. Folytasd csak
tovább, és halld: \textit{Isten előtt kedves áldozat a megtört lélek, te nem veted meg a töredelmes
és alázatos szívet} (Zsolt 50, 19). Vesd el, amit hajdan áldoztál, hiszen megtaláltad, hogy ezután
mit áldozzál. Hajdan atyáid állatokat öltek le az oltáron, és azt tartották áldozatnak. Te a véres
áldozatot nem kedveled. Tehát azokat már ne keresd, hanem keresd a te áldozatodat.

Égőáldozataimban – mondja – nem leled kedved. Tehát ha égőáldozataimra rá sem nézel, vajon akkor
semmiféle áldozat sem jó neked? Dehogynem! Isten előtt kedves áldozat a megtört lélek, te nem veted
meg a töredelmes és alázatos szívet. Van tehát, mit áldozzál. Hát akkor ne nyájak körül nézelődj, de
hajóútra se készülj, hogy távoli országokból drága illatszereket hozzál. Saját szívedben keressed,
mi volna kedves Istennek. Szívedet kell megtörnöd. Mit félsz attól, hogy elvész, ha összetöröd? Hisz
ott áll a zsoltárban megírva: \textit{Tiszta szívet teremts bennem, Istenem!} (Zsolt 50, 12) Hogy
tiszta szívet lehessen benned teremteni, ahhoz a tisztátalant össze kell törnöd.

Ne legyen tetszésünkre az, ha vétkezünk, mert a bűn nem tetszik Istennek. Mivel nem vagyunk bűn
nélkül, legalább abban hasonlítsunk Istenhez, hogy nekünk se legyen a tetszésünkre az, ami neki nem
tetszik. Legalább annyira egyesülj Isten akaratával, hogy neked sem tetszik tenmagadban, amit utál
az, aki téged teremtett.
\end{lectio}

\parttitle[Zsolt 50, 12]{Válaszos ének}

\versicle{%
    Bűneim, Uram, mint a nyilak átszög\underline{e}ztek \underline{e}ngem, † de mielőtt sebeim
    elmérges\underline{e}d\-n\underline{é}nek, * Uram, Istenem gyógyíts meg engem a bűnbocsánat
    gy\underline{ó}gyszer\underline{é}vel.
}{%
    Tiszta szívet teremts bennem, Ist\underline{e}nem; * új és erős lelket \underline{ö}nts
    b\underline{e}lém.
}

\rubric{Te Deum, \pageref{TeDeum}.}

\parttitle{Könyörgés}

\rubric{Mint a Reggeli dicséretben.}

\subsubsection*{Reggeli dicséret}

\parttitle{Himnusz}

\printsheet{EcceIamNoctis}

hogy mindnyájunkon könyörüljön Isten,\\
aggodalmunkat szent keze elűzze,\\
s jó Atyánk lévén, tegye, hogy bejussunk\\
égi hazánkba.\\

Adja meg mindezt az isteni Fölség:\\
az Atyaisten s örök Egyszülöttje,\\
és a Szentlélek, akit áld az ének\\
mind e világon. Ámen.

\parttitle{Zsoltározás}

\redtext{1. antifóna}

\printsheet{DominusMihiAdiutor}

\begin{psalmsheet}[tone=8, continuous, number={117. zsoltár}, title={Az ujjongás és szabadulás éneke}, motto={Krisztus a kő, amelyet ti, az építők elvetettetek, mégis szegletkővé lett (ApCsel 4, 11).}]
    \psalmtext{
        1. Adjatok hálát az Úrnak, mert jó, * mert örökké szeret minket.

        2. Mondja Izrael háza: Jó az Úr, * mert örökké szeret minket.\\
        3. Mondja Áron papi háza is, * hogy örökké szeret minket.\\
        4. Hirdessék azok is, akik félik az Urat, * hogy örökké szeret minket. ---

        5. Az Úrhoz fohászkodtam a szenvedés idején, * meghallgatott, és tágas útra vezetett.\\
        6. Velem van az én Uram, Istenem, * mitől féljek, ember mit árthat nekem?\\
        7. Velem van az Úr, aki mindenben segítőm, * és megláthatom, hogyan szégyenül meg ellenségem. ---

        8. Jobb az Úrnál keresni oltalmat, * mint emberekben bizakodni.\\
        9. Jobb az Úrnál keresni oltalmat, * mint hatalmasokban reménykedni.\\
        10. Körülvettek mind a nemzetek, * de az Úr nevében porba sújtom őket.\\
        11. Szorongatnak minden oldalról, * de az Úr nevében porba sújtom őket.\\
        12. Körülvettek, mint a méhrajok, † tüzük felém csapott, mint a rőzseláng, * de az Úr nevében porba sújtom őket. ---

        13. Nekem jöttek, ütlegeltek, már-már elestem, * de az Úr felkarolt, és megsegített.\\
        14. Az Úr erősségem és dicsőségem, * ő lett az én szabadítóm.\\
        15. Örömujjongástól és énektől hangosak * az igazak sátrai. ---

        16. Az Úr jobbja tett nagy dolgokat, † az Úr jobbja emelt föl engem, * az Úr jobbja művelt csodákat!\\
        17. Nem halok meg, hanem élek, * hogy hirdessem az Úr nagy tetteit.\\
        18. Megfenyített az Úr és megpróbált, * de nem adott halálra engem. ---

        19. Tárjátok ki az igazak kapuit, * bemegyek, hogy áldjam az Urat.\\
        20. Az Úr kapuja, íme, itt van, * de csak az igazak mennek át rajta. ---

        21. Hálát adok, mert meghallgattál engem, * te lettél szabadítóm.\\
        22. A kő, amelyet az építők félredobtak, * az lett a háznak szegletköve.\\
        23. Mindezt az Úr vitte végbe, * a csodát saját szemünkkel látjuk.\\
        24. Ez az a nap, amelyet az Úr adott, * hogy ujjongjunk és örvendezzünk. ---

        25. Adj nekünk, Uram, szabadulást, * adj nekünk, Uram, jólétet!\\
        26. Áldott, aki az Úr nevében jön! * Megáldunk az Úr házából titeket.\\
        27. Isten az Úr, ő ragyog felettünk; * álljatok sort lombos ágakkal egészen az oltár szarváig! ---

        28. Hálát adok neked, te vagy Istenem, * magasztallak, Istenem, téged!\\
        29. Adjatok hálát az Úrnak, mert jó, * mert örökké szeret minket!
}
\end{psalmsheet}

\redtext{2. antifóna}

\parttitle[Dán 3,52-57]{Kantikum}

\printsheet{TresExUno}

\begin{psalmsheet}[tone=8, continuous, title={Minden teremtmény dicsérje az Urat!}, motto={A Teremtő… áldott legyen mindörökké! (Róm 1, 25)}]
    \psalmtext{
        52. Áldott vagy, Urunk, atyáink Istene, * dicséretre méltó és magasztos örökké. ---

        52. És áldott a te dicsőséges szent neved, * dicséretre méltó és magasztos mindörökkön örökké. ---

        53. Áldott vagy templomod dicsőséges szent helyén, * nagyon-nagyon dicséretre méltó és magasztos örökké. ---

        54. Áldott vagy királyi trónodon, * nagyon-nagyon dicséretre méltó és magasztos örökké. ---

        55. Áldott vagy, aki a mélységekbe tekintesz, † és a kerubok fölött trónolsz, * dicséretre méltó és magasztos örökké. ---

        56. Áldott vagy az égbolt fölött, * dicséretre méltó és magasztos örökké! ---

        57. Áldja az Urat az Úr minden műve, * dicsérje és magasztalja örökkön örökké!
}
\end{psalmsheet}

\redtext{3. antifóna}

\printsheet{LaudateDominum}

\begin{psalmsheet}[tone=8, continuous, number={150. zsoltár}, title={Dicsérjétek az Urat!}, motto={Daloljon a lelketek, daloljon a szívetek is, vagyis dicsőítsétek meg Istent lelketekben és testetekben egyaránt! (Hesychius)}]
    \psalmtext{
        1. Dicsérjétek az Urat szentélyében, * dicsérjétek a hatalmas égboltozaton!\\
        2. Dicsérjétek hatalmas tetteiért, * dicsérjétek, mert nagysága végtelen. ---

        3. Dicsérjétek őt harsonaszóval, * dicsérjétek lanttal, citerával.\\
        4. Dicsérjétek dobbal és körtánccal, * dicsérjétek citerával és fuvolával. ---

        5. Dicsérjétek őt zengő cimbalommal, † dicsérjétek búgó cimbalom hangjával! * Minden lélek az Urat dicsérje!
}
\end{psalmsheet}

\parttitle[Ez 36,25-27]{Rövid olvasmány}
Tiszta vizet hintek rátok, és megtisztultok minden tisztátalanságtól, minden bálványotoktól
megtisztítal\underline{a}k. † És új szívet adok nektek, új lelket öntök belétek. Kőszíveteket
elveszem testetekből, és hússzívet \underline{a}dok nektek. * Lelkemet árasztom belétek; ezt teszem,
hogy törvényeim szerint járjatok, hogy ítéleteimet megtartsátok, és
megt\underline{e}gy\underline{é}tek.

\parttitle{Rövid válaszos ének}

\printsheet{ConfitebimurTibiDeus}

\parttitle[Lk 1,68-79]{Evangéliumi kantikum}

\printsheet{ConfiteorTibiPater}

\parttitle{Fohászok}

\begin{preces}{Adjunk hálát Üdvözítőnknek, aki e világra jött, hogy velünk legyen, mint Istenünk.
Szólítsuk őt hangos szóval:}{Krisztus, dicsőség Királya, légy világosságunk és örömünk!}
  Krisztus Urunk, égből jött fényesség, az eljövendő feltámadás zsengéje,

  \reddash ne engedd, hogy a halál árnyékában üljünk, hanem add, hogy téged kövessünk és az élet világosságában járjunk!

  Mutasd meg nekünk jóságodat, amely bőséggel árad minden teremtményedre,

  \reddash hogy mindenütt szemléljük dicsőségedet!

  Ne engedd, Urunk, hogy a mai napon legyőzzön minket a Gonosz,

  \reddash hanem segíts, hogy legyőzzük a rosszat a jóval!

  Te megkeresztelkedtél a Jordán vizében, és fölkent a Szentlélek,

  \reddash add, hogy a mai napon a Szentlélek kegyelme vezessen minket!
\end{preces}
Mi Atyánk...

\parttitle{Könyörgés}

\label{DXIV/L/Coll}

Istenünk, te szent Fiad megaláztatása által fölemelted az elesett emberiséget. Adj szent örömet
nekünk, akiket kiragadtál a bűn szolgaságából, és fogadj be egykor az örök boldogságba. A mi Urunk.

\label{D/Compl2}

\subsubsection*{Befejező imaóra a II. Esti dicséret után}

\parttitle{Zsoltározás}

\redtext{Antifóna}

\printsheet{QuiHabitat}

\begin{psalmsheet}[tone=8, continuous, number={90. zsoltár}, title={A Fölséges oltalmában}, motto={Hatalmat adtam nektek, hogy kígyókon és skorpiókon járjatok (Lk 10, 19).}]
    \psalmtext{
        1. Aki a Fölséges oltalmában lakik, * a Mindenható árnyékában pihen,\\
        2. ezt mondja az Úrnak: † „Te vagy menedékem és erősségem, * Istenem, benned bízom.” ---

        3. Megment ő téged a vadászok tőrétől * és a gonosz beszédtől.\\
        4. Szárnyaival oltalmaz, * tollai alatt menedéket találsz.\\
        5. Pajzs és páncél az ő hűsége; * nem kell félned az éjjeli rémtől,\\
        6. a nappal szárnyaló nyíltól, † a sötétben ólálkodó ragálytól, * a fényes nappal pusztító dögvésztől. ---

        7. Essenek el bár ezren oldaladon, † és tízezren jobbod felől, * nem érhet téged semmi bántódás.\\
        8. Saját szemeddel láthatod majd, * látni fogod a bűnösök bűnhődését.\\
        9. Mert az Úr a te védelmező várad, * a Magasságbelit választottad oltalmadul. ---

        10. Így nem ér semmi baj, * otthonodhoz nem közelíthet semmi csapás,\\
        11. mert megparancsolja angyalainak, * hogy őrizzenek minden utadon.\\
        12. Tenyerükön fognak hordozni, * hogy kőbe ne üsd lábadat.\\
        13. Áspiskígyón és viperán lépdelsz, * kígyót és oroszlánt eltiporsz. ---

        14. Hűséges hozzám, azért megmentem, * oltalmazom, mert nevemet ismeri.\\
        15. Ha kiált hozzám, meghallgatom, † minden szükségben közel vagyok hozzá, * megszabadítom, és dicsőséget szerzek neki.\\
        16. Hosszú élettel áldom meg, * és megmutatom neki üdvösségemet.
}
\end{psalmsheet}

\parttitle[MTörv 6, 4-7]{Rövid olvasmány}
Látni fogják az Úr arcát, és a homlokukon lesz a nev\underline{e}. † Nem lesz többé éjszaka, és nem
szorulnak rá a lámpa világítására, sem \underline{a} nap fényére. * Az Úristen ragyogja be őket, és
uralkodni fognak örökkön-\underline{ö}r\underline{ö}kké.

\parttitle{Könyörgés}
Istenünk, a mai napon megünnepeltük Urunk feltámadását. Hallgasd meg alázatos kérő szavunkat, hogy
minden rossztól mentesen a te békédben nyugodjunk, és dicséretedben örvendezve ébredjünk. Krisztus,
a mi Urunk által.

\subsection*{Hétfő, 14. hét - Előtábor 3. nap}

\label{f2/Inv}

\subsubsection*{Imádságra hívás}

\printsheet{IubilemusDeo}

\subsubsection*{Olvasmányos imaóra}

\parttitle{Himnusz}

\printsheet{IpsumNuncNobis}

\begin{hymnstanzas}
    Az okos szüzek csapata\\
    már készenlétben várja őt;\\
    míg égő mécset tart kezük\\
    nagy vígan ujjong örömük.
\stanzabreak
    A balgák késnek, mécsesük\\
    nem fénylik: lángja kialudt.\\
    Az ajtót hasztalan verik,\\
    bezárult már az ég nekik.
\stanzabreak
    Mi csak virrasszunk éberen,\\
    szívünk világa most a mécs,\\
    hogy amikor megérkezik,\\
    méltón fussunk az Úr elé.
\stanzabreak
    Országodhoz, áldott Király,\\
    ha eljössz méltónak találj,\\
    hadd zengjünk akkor teneked\\
    meg nem szűnő dicséretet. Ámen.
\end{hymnstanzas}

\parttitle{Zsoltározás}

\redtext{1. antifóna}

\printsheet{InTuaIustitia}

\begin{psalmsheet}[tone=8, continuous, number={30. zsoltár, 2-17. 20-25}, title={Üldözött ember bizakodó imádsága}, motto={Atyám, kezedbe ajánlom lelkemet (Lk 23, 46).}, numeral={I.}]
    \psalmtext{
        2. Te vagy, Uram, én reményem, † ne hagyj soha szégyent érnem, * igazságodban szabadíts meg engem!\\
        3. Fordítsd felém füledet, * siess, Uram, ments meg engem!\\
        3. Te légy oltalmazó kősziklám, és megerősített házam, * hogy megszabadíts engem! ---

        4. Erősségem és menedékem valóban te vagy, * neved miatt vezess és irányíts engem!\\
        5. Kiszabadítasz a hálóból, amit titkon vetettek nekem, * mivel te vagy az én erősségem! ---

        6. Uram, kezedbe ajánlom lelkemet; * megváltottál engem, hűséges Istenem!\\
        7. Te utálod a hazug bálványok tisztelőit, * én azonban az Úrban remélek! ---

        8. Ujjongok és örvendezem jóságod miatt, * mert tekintetre méltattad alázatosságomat.\\
        9. Ismered lelkem gyötrődéseit, † nem adtál ellenségeim kezébe, * lábamat tágas helyre állítottad.
}
\end{psalmsheet}

\redtext{2. antifóna}

\printsheet{FaciemTuam}

\begin{psalmsheet}[tone=8, continuous, numeral={II.}]
    \psalmtext{
        10. Szorongatnak, Uram, könyörülj rajtam, † szemem a bánattól elhomályosult, * sorvadozik testem-lelkem.\\
        11. Életem már elernyed a kíntól, * sóhajtozásban fogynak el éveim.\\
        11. Erőm megtört a nyomorúságban, * csontjaim is sorvadoznak. ---

        12. Gyalázattá lettem minden ellenségem előtt, † szomszédaim és ismerőseim félnek tőlem, * ha meglátnak az utcán, elkerülnek.\\
        13. Szívükben elfelednek, mintha meghaltam volna, * olyanná lettem, mint az összetört edény.\\
        14. Sokak gyalázkodó szavát hallom; * rettegés vesz körül engem,\\
        14. mert ezek összefognak ellenem, * azon tanakodnak, hogy életemet kioltsák. ---

        15. Én pedig benned remélek, Uram, † azt mondtam: „Te vagy Istenem, * kezedben van a sorsom.”\\
        16. Ragadj ki az ellenség kezéből, * és azoktól, akik üldöznek engem.\\
        17. Szolgád fölé ragyogtasd fel arcodat, * irgalmasságodban szabadíts meg engem!
}
\end{psalmsheet}

\redtext{3. antifóna}

\printsheet{DiligiteDominum}

\begin{psalmsheet}[tone=2, continuous, numeral={III.}]
    \psalmtext{
        20. Mily nagy a te jóságod, Uram, * amelyet a téged félőknek tartogatsz,\\
        20. amelyet azoknak készítettél, akik benned bíznak, * minden ember szeme láttára. ---

        21. Arcod védő árnyékában rejted el őket * az emberek zaklatása elől.\\
        21. Sátradban adsz nekik biztos helyet, * a gonosz nyelvektől távol. ---

        22. Legyen áldott az Úr, * aki csodás szeretettel véd meg erős városában.\\
        23. Én pedig ijedtemben már azt mondtam: * „Szemed elől elvetettél engem!”\\
        23. De te meghallgattad könyörgő szavamat, * amikor hozzád kiáltottam. ---

        24. Szeressétek az Urat, ti, összes szentjei, † a hűségeseket megvédi az Úr, * de bőven megfizet a gőgösöknek.\\
        25. Legyetek bátrak, és erősödjék szívetek, * mind, akik az Úrban reménykedtek.
}
\end{psalmsheet}

\versicle{Igazságod szerint vezess, taníts engem, Ur\underline{a}m,}{Mivel te vagy az én megmentő Isten\underline{e}m.}

\parttitle[Zsolt 40,10; Mk 14,18]{Válaszos ének}

\versicle{%
    Még b\underline{a}rát\underline{o}m is, † akiben ped\underline{i}g b\underline{í}ztam, * Aki
    együtt ette kenyerét velem, sarkát emelt\underline{e} ell\underline{e}nem.
}{%
    Egyiketek elárul \underline{e}ngem, * egy, aki esz\underline{i}k v\underline{e}lem.
}

\parttitle[1 Kor 9,19. 22; Jób 29,15-16]{Válaszos ének}

\versicle{%
    Bár mindenkitől független voltam, mégis mindenkinek szolg\underline{á}ja l\underline{e}ttem. † A
    gyöngék közt gyöng\underline{e} l\underline{e}ttem. * Mindenkinek mindene lettem, hogy mindenkit
    \underline{ü}dvöz\underline{í}tsek.
}{%
    A vaknak úgy szolgáltam, hogy a szeme v\underline{o}ltam;; * a szegény embereknek
    atyj\underline{u}k v\underline{o}ltam.
}

\rubric{Könyörgés, mint vasárnap, a Reggeli dicséretben, \pageref{DXIV/L/Coll}.}

\subsubsection*{Reggeli dicséret}

\parttitle{Himnusz}

\printsheet{LucisLargitorSplendide}

\begin{hymnstanzas}
    Igaz fényesség csak te vagy,\\
    nem hajnal törpe csillaga,\\
    mely a jövendő napsugárt\\
    szegényes fénnyel jelzi csak,
\stanzabreak
    hanem napnál is teljesebb\\
    verőfény, tűző napvilág;\\
    szívünk legmélyebb rejtekén\\
    sugárzó fényed átragyog.
\stanzabreak
    Szív tisztasága győzze le\\
    lázongó vérünk vágyait;\\
    a tiszta testnek temploma\\
    őrizze lelkünk kincseit.
\stanzabreak
    Krisztus, kegyelmes nagy Király,\\
    adassék néked hódolat,\\
    Atyád áldjuk s a Lelket is\\
    idők végéig szüntelen. Ámen.
\end{hymnstanzas}

\parttitle{Zsoltározás}

\redtext{1. antifóna}

\printsheet{Sitivit}

\begin{psalmsheet}[tone=8, continuous, number={41. zsoltár}, title={Sóvárgás az Úr és az ő temploma után}, motto={Aki szomjazik, jöjjön. Aki kívánja az élet vizét, ingyen igyék (Jel 22, 17).}]
    \psalmtext{
        2. Mint szarvas sóvárog a forrásvízre, * úgy áhít a lelkem, téged, Istenem.\\
        3. Isten után szomjazik a lelkem, az élő Isten után: * mikor mehetek, hogy megjelenjek az ő színe előtt? ---

        4. Éjjel-nappal könny a kenyerem, * mert azt hallom naponta: „Hol a te Istened?”\\
        5. Csak arra gondolok, elmém azon mereng, * hogyan vonultam be a csodás hajlékba Isten házáig,\\
        5. miközben ujjongott és énekelt * az ünneplő sokaság. ---

        6. Miért csüggedsz, lelkem, * miért vagy nyugtalan?\\
        6. Bízzál az Úrban: fogom még áldani * arcom szabadítóját, Istenemet. ---

        7. Szomorkodik bennem a lélek, † hozzád röpül ezért gondolatom * a Jordán és a Hermon távoli hegyes tájairól.\\
        8. Harsognak zuhatagjaid, örvény kiált az örvénynek; * minden kavargó hullámod átzuhog fölöttem.\\
        9. Árassza rám az Úr irgalmát naphosszat, † neki zeng éjjel háladalom, * és áldom éltető Istenemet. ---

        10. Így szólok az Úrhoz: „Oltalmazóm vagy, † miért feledkeztél meg rólam? * Miért kell bánatban élnem, ellenségeimtől gyötörten?”\\
        11. Keserűség marja bensőmet, † amikor kínzóim gúnyolnak, * s naponta azt mondják: „Hol a te Istened?” ---

        12. Miért csüggedsz, lelkem, * miért vagy nyugtalan?\\
        12. Bízzál az Úrban: fogom még áldani * arcom szabadítóját, Istenemet.
}
\end{psalmsheet}

\redtext{2. antifóna}

\parttitle[Sir 36,1-7. 13-16]{Kantikum}

\printsheet{Ostende}

\begin{psalmsheet}[tone=3, continuous, title={Imádság a szent városért, Jeruzsálemért}, motto={Az az örök élet, hogy ismerjenek téged, egyedül igaz Istent, és akit küldtél, Jézus Krisztust (Jn 17, 3).}]
    \psalmtext{
        1. Könyörülj rajtunk, tekints ránk, mindenség Istene, * szereteted fényét mutasd meg nekünk!\\
        2. Bocsásd félelmedet a nemzetekre, * amelyek nem ismernek téged.\\
        2. Hadd tudják meg, hogy rajtad kívül nincs más Isten, * és hirdessék nagy tetteidet. ---

        3. Emeld föl kezedet az idegen népek ellen, * hogy lássák hatalmadat.\\
        4. Amint szemük láttára szent voltál velünk, * úgy a mi szemünk láttára légy dicső fölöttük,\\
        5. hogy úgy ismerjenek, ahogy mi ismerünk, * mert rajtad kívül, Urunk, nincs más Isten!\\
        6. Újítsd meg a jeleket, és tégy csodákat újból, * dicsőítsd meg kezedet, erős karodat! ---

        13. Gyűjtsd össze Jákob minden törzsét, † hadd tudják meg, hogy rajtad kívül nincs más Isten, * és mondják el nagy tetteidet.\\
        13. Tekintsd őket örökségednek, * mint a régmúlt időkben. ---

        14. Könyörülj népeden, amelyet terólad neveztek el, * Izraelen, amelyet elsőszülöttednek tekintettél,\\
        15. és könyörülj szent városodon, * nyugalmad városán, Jeruzsálemen!\\
        16. Áraszd el Siont kimondhatatlan igéiddel, * és népedet dicsőségeddel!
}
\end{psalmsheet}

\redtext{3. antifóna}

\printsheet{Opera}

\begin{psalmsheet}[tone=2, continuous, number={18 A. zsoltár, 2-7}, title={A teremtő Isten dicsérete}, motto={Eljön hozzánk felkelő Napunk a magasságból… hogy lépteinket a békesség útjára vezérelje (Lk 1, 78. 79).}]
    \psalmtext{
        2. Isten dicsőségét hirdetik az egek, * kezének művéről vall a mennybolt.\\
        3. A nappalok ezt zengik egymásnak, * erre oktatja éj az éjszakát. ---

        4. Nem olyan hang ez, nem olyan beszéd, * hogy meg ne értenéd szavát.\\
        5. Minden földre elhat szózatuk, * a földkerekség végéig eljut igéjük. ---

        6. Az égen vert Isten sátort a napnak, † és az, mint nászházból kilépő vőlegény, * mint ujjongó dalia fut neki útjának. ---

        7. Az ég egyik pereméről indul, † és másik széléig ér a pályája: * tüze elől rejtőzni nem lehet.
}
\end{psalmsheet}

\parttitle[Jer 15, 16]{Rövid olvasmány}
Ha szavaid elém kerültek, csak úgy nyeltem ők\underline{e}t; † kedvem telt szavadban és
ör\underline{ö}me szívemnek. * Mert a te nevedet viselem, Uram, Seregek
\underline{I}st\underline{e}ne!

\parttitle{Rövid válaszos ének}

\printsheet{ExultateIustiInDomino}

\parttitle[Lk 1,68-79]{Evangéliumi kantikum}

\printsheet{BenedictusDominusQuia}

\parttitle{Fohászok}

\begin{preces}{Adjunk hálát Üdvözítőnknek, aki e világra jött, hogy velünk legyen, mint Istenünk.
Szólítsuk őt hangos szóval:}{Krisztus, dicsőség Királya, légy világosságunk és örömünk!}
  Krisztus Urunk, égből jött fényesség, az eljövendő feltámadás zsengéje,

  \reddash ne engedd, hogy a halál árnyékában üljünk, hanem add, hogy téged kövessünk és az élet világosságában járjunk!

  Mutasd meg nekünk jóságodat, amely bőséggel árad minden teremtményedre,

  \reddash hogy mindenütt szemléljük dicsőségedet!

  Ne engedd, Urunk, hogy a mai napon legyőzzön minket a Gonosz,

  \reddash hanem segíts, hogy legyőzzük a rosszat a jóval!

  Te megkeresztelkedtél a Jordán vizében, és fölkent a Szentlélek,

  \reddash add, hogy a mai napon a Szentlélek kegyelme vezessen minket!
\end{preces}
Mi Atyánk...

\rubric{Könyörgés, mint vasárnap, a Reggeli dicséretben, \pageref{DXIV/L/Coll}.}

\subsubsection*{Befejező imaóra}

\parttitle{Zsoltározás}

\redtext{Antifóna}

\printsheet{IntendeVoci}

\begin{psalmsheet}[tone=8, continuous, number={85. zsoltár}, title={Bajba jutott szegény ember imája}, motto={Áldott legyen az Isten, aki megvigasztal minket minden szomorúságunkban (2 Kor 1, 3. 4).}]
    \psalmtext{
        1. Fordulj hozzám, és hallgass meg, Uram, * mert gyámoltalan és szegény vagyok.\\
        2. Végy oltalmadba, mert neked élek; * szabadítsd meg, Istenem, szolgádat, aki benned bízik. ---

        3. Könyörülj rajtam, Uram; * látod, egész nap hozzád kiáltok.\\
        4. Vigasztald meg szolgád szívét, * hiszen lelkem hozzád emelem, Uram,\\
        5. mert te, Uram, jóságos vagy és szelíd, * és nagy irgalmú azokhoz, akik hozzád kiáltanak. ---

        6. Imádságomat hallgasd meg, Uram, * könyörgésem szavára figyelj!\\
        7. Szorongatásom napján hozzád kiáltok, * mivel te meghallgatsz engem. ---

        8. Az istenek között nincs hozzád hasonló, Uram, * és műveidhez semmi sem fogható.\\
        9. Térjen hozzád minden nemzet, melyet csak alkottál, † boruljanak le előtted, Uram, * és magasztalják nevedet,\\
        10. mert nagy vagy, és csodákat művelsz: * te vagy Isten egyes-egyedül. ---

        11. Mutasd meg nekem utadat, Uram, * hogy igazságod szerint járjak.\\
        11. Tedd egyszerűvé szívemet, * hogy félje nevedet.\\
        12. Magasztallak, Uram, Istenem, teljes szívemből, * és nevedet dicsőítem mindörökké,\\
        13. mert annyira szeretsz engem, * hogy a sír mélyéből is kimentesz. ---

        14. Istenem, kevélyek kelnek fel ellenem, † hatalmaskodók serege tör vesztemre, * de nem számolnak veled.\\
        15. Pedig te, Uram, irgalmas és jóságos Isten vagy, * hosszan tűrő, könyörülő és hűséges.\\
        16. Tekints rám, és szánj meg engem, † adj erőt szolgádnak, * és szabadítsd meg szolgálód fiát! ---

        17. Mutasd meg rajtam jóságod jelét; † hadd lássák gyűlölőim, és szégyenkezzenek, * mert te, Uram, megsegítesz és megvigasztalsz engem.
}
\end{psalmsheet}

\parttitle[1 Tessz 5, 9-10]{Rövid olvasmány}
Isten az üdvösség elnyerésére szánt minket Urunk, Jézus Krisztus ált\underline{a}l, † aki meghalt
értünk, hogy akár élünk, \underline{a}kár meghalunk, * vele együtt elnyerjük az
\underline{é}l\underline{e}tet.

\parttitle{Könyörgés}
Istenünk, testünknek adj üdvös pihenést, és mai munkálkodásunk magvetését érleld meg az örök
aratásra. Krisztus, a mi Urunk által.

\subsection*{Kedd, 14. hét - Előtábor 4. nap}

\subsubsection*{Imádságra hívás}

\printsheet{RegemMagnum}

\subsubsection*{Olvasmányos imaóra}

\parttitle{Himnusz}

\printsheet{NocteSurgentes}

Jó Királyunknak, akik együtt zengünk,\\
hadd legyünk méltók az egekbe jutva,\\
ott, hol a szentek vele laknak együtt,\\
boldogan élni.

Add meg ezt, kérünk, örök Isten: áldott\\
jó Atyánk, Krisztus, s velük egybeforrott\\
éltető Lélek, kit a földön minden\\
zengve magasztal. Ámen.

\parttitle{Zsoltározás}

\redtext{1. antifóna}

\printsheet{Revela}

\begin{psalmsheet}[tone=6, continuous, number={36. zsoltár}, title={A gonoszok és a jók sorsa}, motto={Boldogok a szelídek, mert övék lesz a föld (Mt 5, 5).}, numeral={I.}]
    \psalmtext{
        1. Ne bosszankodj a bűnösökre, * és ne háborogj a gonosztevők miatt,\\
        2. mert mint a fű, hamar elhervadnak, * és elszáradnak, mint a zöldellő növény. ---

        3. Bizakodj az Úrban, és jót cselekedj, * akkor megmaradsz földeden, és biztonságban élhetsz.\\
        4. Örömödet az Úrban keresd, * ő majd teljesíti szíved vágyait. ---

        5. Bízd sorsod az Úrra, benne remélj, * és ő gondoskodik rólad.\\
        6. Napfényre hozza igazságodat, * és igaz ügyed ragyog, mint a déli napsugár. ---

        7. Légy csendben, és várj az Úrra, † ne irigykedj arra, akinek jól megy a sora, * noha álnokságon jártatja eszét.\\
        8. Sose bosszankodj, és ne mérgelődj, * ne is irigykedj, mert az rosszra visz. ---

        9. A gonosztevők sorsa pusztulás, * ám azoké lesz a föld, akik az Úrban bíznak.\\
        10. Egy kis idő, és már nem is lesz a bűnös, * keresve sem találod helyét.\\
        11. Ám a szelídek öröklik a földet, * és élvezik majd a béke áldásait.
}
\end{psalmsheet}

\redtext{2. antifóna}

\printsheet{IuniorFui}

\begin{psalmsheet}[tone=2, continuous, numeral={II.}]
    \psalmtext{
        12. Az igaz kárára cselt sző a gonosz, * és ellene vicsorgatja fogát.\\
        13. Az Úr azonban kineveti őt, * mert látja, hogy eljön rá a bűnhődés napja. ---

        14. Kardot rántanak a gonoszok, * íjat is feszítenek,\\
        14. hogy elejtsék a szegényt és a nyomorultat, * hogy megöljék azokat, akik szentül élnek.\\
        15. De kardjuk saját szívüket járja át, * és összetörnek íjaik. ---

        16. Jobb az igaz ember szűkös birtoka, * mint a gonosz bővelkedése,\\
        17. mert szertefoszlik a bűnös hatalma, * de az igazat pártolja az Úr. ---

        18. Ismeri az Úr az ártatlanok sorsát: * örökségük megmarad örökre.\\
        19. Nem szégyenülnek meg balsors idején, * s van eledelük ínséges időben. ---

        20. Mert elvesznek a bűnösök, † az Úr ellenségei elhullanak, mint a mezők dísze, * füstként enyésznek el.\\
        21. Kölcsönt kér a gonosz, de nem adja vissza; * ám az igaz jó szívvel osztogat. ---

        22. Mert azoké a föld, akiket Isten megáld, * s akiket megátkoz, elpusztulnak.\\
        23. A jámbor ember lépteit az Úr vezérli, * életútjában kedvét találja.\\
        24. Még ha elesik is, akkor sem éri baj, * hiszen az Úr fogja kezét. ---

        25. Ifjú voltam, s az idő eljárt fölöttem, † elhagyatott igazat mégsem láttam soha, * még ivadéka sem jutott koldusbotra.\\
        26. Mindenkor jószívű, és kölcsönt adni kész, * áldott lesz ezért sarjadéka is. ---

        27. Hagyd el a rosszat, s tedd a jót, * és megmaradsz örökkön-örökké,\\
        28. mert kedveli az Úr az igazságot, * és nem hagyja cserben híveit.\\
        28. A gonoszok elvesznek örökre, * a bűnösök sarjai semmivé válnak,\\
        29. ám az igazak öröklik a földet, * és mindörökké lakni fogják.
}
\end{psalmsheet}

\redtext{3. antifóna}

\printsheet{SperaInDomino}

\begin{psalmsheet}[tone=6, continuous, numeral={III.}]
    \psalmtext{
        30. Az igaz szája bölcsességet szól, * és nyelve igazságot beszél.\\
        31. Isten törvénye él a szívében, * nem is tétováznak soha léptei.\\
        32. A hívőre leselkedik a gonosz, * és arra tör, hogy elvegye életét.\\
        33. Az Úr azonban nem adja kezére, * s nem ítéli el, ha majd ítélőszéke elé kerül.

        34. Bízzál az Úrban, s maradj meg az ő útján, † melléd áll akkor, és a földet elnyered, * s meglátod a bűnösök vesztét.\\
        35. Láttam én már gőgös, gonosz embert, * mint viruló cédrus, fenn hordta fejét.\\
        36. Térültem-fordultam, és már sehol sem volt, * kerestem, de helyét sem találtam.

        37. Nézd az ártatlant, figyeld az igazat: * a békés embernek marad ivadéka.\\
        38. Ám a gonoszok hamar kipusztulnak, * megsemmisül a bűnös nemzedék.\\
        39. Az igazak üdve az Úrtól való, * az Úr oltalmazza őket szenvedés idején.\\
        40. Segítségük ő és szabadítójuk, † a bűnösök közül kiragadja, s megmenti őket, * mivel reményüket őbelé vetették.
}
\end{psalmsheet}

\versicle{Igazi jóságra, okosságra és bölcsességre oktass eng\underline{e}m.}{Mert én hiszek parancsolataidb\underline{a}n.}

\parttitle[Zsolt 54, 13. 14. 15; vö. 40, 10; 2 Sám 18, 33]{Válaszos ének}

\versicle{%
    Hogyha ellenségem átk\underline{o}zna \underline{e}ngem, † azt még
    elv\underline{i}s\underline{e}lném; * De te, olyan ember, aki egyenlő velem, akivel oly édes
    volt a barátságunk, sarkadat emelt\underline{e}d ell\underline{e}nem.
}{%
    Megrendült a k\underline{i}rály, † fölment a kapu fölötti szobába, \underline{é}s sírt, * s
    járva-kelve ezt mond\underline{o}g\underline{a}tta:
}

\parttitle[Ef 4, 1. 3. 4]{Válaszos ének}

\versicle{%
    Kérlek bennetek\underline{e}t az \underline{Ú}rban, † hogy éljetek méltón ahhoz a hivatáshoz,
    amelyet k\underline{a}pt\underline{a}tok. * Törekedjetek rá, hogy a béke kötelékével
    fenntartsátok a lelk\underline{i} egys\underline{é}get.
}{%
    Egy a Test és egy a L\underline{é}lek, * mint ahogy hivatástok is egy
    rem\underline{é}nyr\underline{e} szól.
}

\rubric{Könyörgés, mint vasárnap, a Reggeli dicséretben, \pageref{DXIV/L/Coll}.}

\subsubsection*{Reggeli dicséret}

\parttitle{Himnusz}

\printsheet{AeternaeLucisConditor}

\begin{hymnstanzas}
    Dereng az égnek alja már,\\
    eloszlik minden éji árny,\\
    kihunynak fent a csillagok,\\
    amint a hajnal felragyog.
\stanzabreak
    Ágyunkból frissen felkelünk,\\
    s már zeng is hálaénekünk,\\
    mert győzött álmos éjszakán\\
    ismét a felkelő sugár.
\stanzabreak
    Kérünk, Urunk, hogy testi vágy,\\
    s e csalfa csábító világ\\
    ne gyújtson szívünkben tüzet,\\
    hamis varázs ne ejtse meg.
\stanzabreak
    Haragtól ment legyen szavunk,\\
    és torkos vágytól asztalunk,\\
    ne vonzzon átkos kapzsiság,\\
    sem tékozló rút léhaság.
\stanzabreak
    Szilárd elmével, józanul,\\
    testben tisztán, nem álnokul,\\
    szolgálva híven, lelkesen\\
    napunk Krisztussal töltsük el.
\stanzabreak
    Ezt add meg, kérünk, jó Atya,\\
    s Atyának mása, egy Fia,\\
    s te, Szentlélek, Vigasztaló:\\
    örökre egy uralkodó. Ámen.
\end{hymnstanzas}

\parttitle{Zsoltározás}

\redtext{1. antifóna}

\printsheet{Salutare}

\begin{psalmsheet}[tone=6, continuous, number={42. zsoltár}, title={Sóvárgás a templom után}, motto={Világosságul jöttem a világba (Jn 12, 46).}]
    \psalmtext{
        1. Szolgáltass nekem igazságot, Istenem, † az istentelen néppel szemben védd ügyemet, * ments meg a gonosz és álnok emberektől!\\
        2. Mert te vagy az én megmentő Istenem, † miért taszítasz hát el magadtól? * Miért kell bánatban élnem, ellenségtől gyötörten? ---

        3. Világosságodat és hűségedet küldd el, † azok vezéreljenek engem * szent hegyedre, a te hajlékodba.\\
        4. Odalépek akkor Isten oltárához, † Istenhez, aki nekem ujjongó örömöm. * Citeraszóval magasztallak téged, Isten, én Istenem. ---

        5. Miért csüggedsz, lelkem, * miért vagy nyugtalan?\\
        6. Bízzál az Úrban: fogom még áldani * arcom szabadítóját, Istenemet.
}
\end{psalmsheet}

\redtext{2. antifóna}

\parttitle[Iz 38, 10-14. 17-20]{Kantikum}

\printsheet{CunctisDiebus}

\begin{psalmsheet}[tone=3, continuous, title={A halálos beteg szorongása és a meggyógyult öröme}, motto={Én vagyok az élő: meghaltam, … de most már nálam van a halál kulcsa (Jel 1, 17-18).}]
    \psalmtext{
        10. Így szóltam: „Már életem napjainak közepén * indulnom kell az alvilág kapui felé.”\\
        11. Kutattam, életemből mennyi van még hátra; * azt mondtam: „Nem látom meg az Urat az élők földjén,\\
        11. és nem látok többé embert * a föld lakói közül.” ---

        12. Sátramat lebontják, elveszik tőlem, * mint ahogy lebontják sátrukat a pásztorok.\\
        12. Életem fonalát mintha takács vágta volna el, † alig kezdődött szövése, máris megszakadt; * reggeltől estig végzel velem.\\
        13. Virradatig reménykedtem; † mintha oroszlán marcangolná minden csontomat; * reggeltől estig végzel velem.\\
        14. Mint a fecskefióka, úgy csipogok, * sírva nyögök, mint a gerle.\\
        15. Szemem elbágyad egészen, * s közben a magasba tekintek. ---

        17. Te azonban kimentettél, hogy el ne vesszek, * és hátad mögé dobtad minden bűnömet.\\
        18. Mert nem az alvilág dicsőít téged, † és a halál sem magasztal; * akik a sírba hulltak, hűségedben nem remélnek többé.\\
        19. Csak az élő, csak az élő dicsér, mint ma én is; * hűségedet az apák hirdessék fiaiknak.\\
        20. Uram, szabadíts meg, † hogy énekelhessük dalaidat * életünk minden napján az Úr házában.
}
\end{psalmsheet}

\redtext{3. antifóna}

\printsheet{TeDecetHymnus}

\begin{psalmsheet}[tone=2, continuous, number={64. zsoltár}, title={Ünnepélyes hálaadás}, motto={Sion alatt a mennyei várost kell értenünk (Origenész).}]
    \psalmtext{
        2. Téged illet, Isten, Sion hegyén a dicsérő ének, * neked teljesítik Jeruzsálemben a fogadalmakat.\\
        3. Mivel te meghallgatod az imát, * hozzád viszi bűne terhét minden ember.\\
        4. Terheljenek bár gonosz tetteink, * eltörlöd azokat irgalmasan.\\
        5. Boldog, akit kiválasztasz, és magadhoz emelsz, * hogy csarnokodban élhesse életét.\\
        5. Ó, hadd telhessünk el házad javaival, * templomod szentségével!\\
        6. Mivel igaz vagy, csodákat művelsz, † s úgy hallgatsz meg minket, szabadító Istenünk, * minden földhatárnak és messzi tengernek reménye! ---

        7. Erőddel szilárd hegyeket emelsz, * hatalommal övezed derekad.\\
        8. Megzabolázod a tenger harsogását, mennydörgő hullámait * és a népek háborgását.\\
        9. Csodajeleidet félik, akik a föld határain laknak: * Keletnek és Nyugatnak te vagy öröme. ---

        10. Meglátogatod a földet, és megöntözöd, * elhalmozod áldásoddal.\\
        10. Az egek csatornáit megtöltöd vízzel; † gabonát nevelsz az embernek, * és termékennyé teszed a földet.\\
        11. Barázdáit megöntözöd, rögeit elboronálod, * záporokkal porhanyítod, s termését megáldod. ---

        12. Az esztendőt áldásoddal koronázod, * s lábad nyomán elárad a bőség.\\
        13. Vízzel töltekeznek a szomjas legelők, * és ujjonganak körös-körül a halmok.\\
        14. A réteket nyájak lepik el, † a völgyekben gabona hullámzik; * minden ujjong, és dalol színed előtt.
}
\end{psalmsheet}

\parttitle[1 Tessz 5, 4-5]{Rövid olvasmány}
Ti, testvérek, nem jártok sötétségben, hogy az Úr napja tolvaj módjára lepjen meg
bennetek\underline{e}t. † Hiszen mindnyájan a világosság és a nappal fi\underline{a}i vagytok. * Nem
vagyunk az éjszakáé, sem a söt\underline{é}ts\underline{é}gé.

\parttitle{Rövid válaszos ének}

\printsheet{VocemMeamAudiDomine}

\parttitle[Lk 1,68-79]{Evangéliumi kantikum}

\printsheet{DeManuOmnium}

\parttitle{Fohászok}

\begin{preces}{Áldjuk Üdvözítőnket, aki feltámadásával fényt árasztott a világra, és hívjuk
    segítségül őt alázatos szóval:}{Urunk, őrizz meg bennünket ösvényeden!}
  Urunk, Jézusunk, reggeli imánkkal feltámadásodat köszöntjük,

  \reddash dicsőségedbe vetett reményünk ragyogja be napunkat!

  Fogadd el, Urunk, vágyakozásainkat és jófeltételeinket,

  \reddash mint mai napunk első ajándékait!

  Add, hogy ma gyarapodjunk az irántad való szeretetben,

  \reddash és bármi történik, mindnyájunk javára váljék!

  Világosodjék általunk a te világosságod az emberek előtt,

  \reddash hogy jócselekedeteink láttán dicsőítsék az Atyát!
\end{preces}
Mi Atyánk...

\parttitle{Könyörgés}

Urunk, Jézus Krisztus, igaz világosság, te megvilágítasz minden embert az üdvösségre. Kérünk, adj
erőt, hogy elkészítsük színed előtt a békesség és az igazság útjait. Aki élsz.

\subsubsection*{Befejező imaóra}

\parttitle{Zsoltározás}

\redtext{Antifóna}

\printsheet{NeIntres}

\begin{psalmsheet}[tone=8, continuous, number={142. zsoltár, 1-11}, title={Imádság megpróbáltatások idején}, motto={Az embert nem a törvény szerinti tettek teszik igazzá, hanem a Jézus Krisztusba vetett hit (Gal 2, 16).}]
    \psalmtext{
        1. Hallgasd meg, Uram, imádságomat, † te hűséges vagy, figyelj könyörgő szavamra, * igazságos vagy, hallgass meg engem!\\
        2. Ne szállj perbe szolgáddal, * színed előtt egy élő sem mondható igaznak. ---

        3. Az ellenség üldözőbe vett, † életemet a földre tiporja, * sötétségbe taszít, mint a régen megholtakat.\\
        4. Egészen elcsügged a lelkem, * és a szívem is megdermed bennem.\\
        5. A régmúlt napokra emlékezem, † végiggondolom minden tettedet, * elmélkedem azon, amit kezed művelt.\\
        6. Kezemet feléd tárom, * mint a kiszikkadt föld, lelkem utánad eped. ---

        7. Siess, hallgass meg, Uram, * már a lelkem is eleped.\\
        7. Ne rejtsd el előlem arcodat, * nehogy olyan legyek, mint akik a sírba szállnak.\\
        8. Hadd halljam minden reggel, hogy jóságos vagy, * mert csak benned bízom.\\
        9. Mutasd meg, melyik úton járjak, * mert hozzád vágyódik a lelkem. ---

        9. Ellenségeimtől ments meg, Uram, * nálad keresek oltalmat.\\
        10. Taníts meg akaratod teljesítésére, † mert te vagy az én Istenem. * Jó lelked vezessen az egyenes úton. ---

        11. Neved miatt, Uram, életemet tartsd meg, * gyötrelmeimből a te igazságod szerint vezess ki engem!
}
\end{psalmsheet}

\parttitle[1 Pét 5, 8-9]{Rövid olvasmány}
Józanok legyetek és vigyázzat\underline{o}k! † Ellenségetek, a sátán, ordító oroszlán módjára ott
kószál mindenütt, és keresi, k\underline{i}t nyeljen el. * Erősen álljatok neki ellen \underline{a}
h\underline{i}tben!

\parttitle{Könyörgés}
Jóságos Istenünk, világosítsd meg ezt az éjszakát, és add, hogy szolgáid, akik most békében
nyugovóra térnek, holnap örömmel ébredjenek az új nap világosságára. Krisztus, a mi Urunk által.

\subsection*{Szerda, 14. hét - Előtábor 5. nap}

\subsubsection*{Imádságra hívás}

\printsheet{DominumQuiFecitNos}

\subsubsection*{Olvasmányos imaóra}

\parttitle{Himnusz}

\printsheet{OSatorRerum}

Éjszaka zengjük hálaénekünket,\\
te segíts minket, hogy méltók lehessünk\\
dicsérni téged egy szívvel örökké,\\
fénynek adója.

Szítsd fel szívünknek szerelmét irántad,\\
add, hogy úgy éljünk, hogy halált ne lássunk,\\
s hogy tetteinkben örök dicsőséged\\
meg-megújuljon.

Add, hogy szívünket tüzed általjárja,\\
s mint Jegyesünket mindig ébren várjunk,\\
és a kezünkben ki nem alvó lánggal\\
égjen a mécses.

Egyként dicsérjük mennyei Atyánkat,\\
s téged, Királyunk, kegyes Üdvözítőnk,\\
Szentlélek Isten, dicséreted zengjen\\
földön, egekben. Ámen.

\parttitle{Zsoltározás}

\redtext{1. antifóna}

\printsheet{UtNonDelinquam}

\begin{psalmsheet}[tone=6, continuous, number={38. zsoltár}, title={Beteg ember imája}, motto={A természet mulandóságnak van alávetve … Isten akaratából (Róm 8, 20).}, numeral={I.}]
    \psalmtext{
        2. Így szóltam: „Vigyázok útjaimra, * nyelvemmel nem vétkezem;\\
        2. retesszel zárom le szájam, * amíg gonosz ember áll előttem.” ---

        3. Hallgattam, elnémultam, de mi haszna, * és fájdalmam kiújult.\\
        4. Szívem fölhevült bennem, * elmélkedtem, s egyre nagyobb lánggal égett. ---

        5. Nyelvem akkor megoldódott: * „Uram, tudasd velem életem végét:\\
        6. mennyi napjaimnak száma, * hadd tudjam meg, mily rövid az életem!” ---

        6. Íme, napjaimat néhány araszra mérted, * életem ideje előtted semmiség.\\
        7. Minden csak hiábavalóság, minden emberélet is, * 7mint az árnyék, az ember úgy múlik el.\\
        8. Hiábavaló az is, hogy vesződik, * kuporgat, de nem tudja, kinek.
}
\end{psalmsheet}

\redtext{2. antifóna}

\printsheet{NonSisMihi}

\begin{psalmsheet}[tone=7, continuous, numeral={II.}]
    \psalmtext{
        8. És most, Uram, mire várhatok? * Egyedül benned reménykedem.\\
        9. Minden gonoszságomból ments ki engem, * ne szolgáltass ki az esztelen gyalázatának!\\
        10. Elnémulok, már nem nyitom ki ajkamat, * mert te vagy, aki ezt végbevitted. ---

        11. Csapásaidat vedd le rólam, * elpusztulok sújtó kezed alatt.\\
        12. Fenyíted az embert gonoszsága miatt, † moly módjára erejét elemészted. * Mint egy lehelet, annyit ér minden ember. ---

        13. Hallgasd meg, Uram, imádságomat, * figyelj kiáltásomra!\\
        13. Ne légy süket siránkozásomra, † hiszen nálad jövevény vagyok, * vándor, mint minden ősöm.\\
        14. Vedd le rólam szemedet, hogy föllélegezzem, * mielőtt elmegyek, és nem leszek többé.
}
\end{psalmsheet}

\redtext{3. antifóna}

\printsheet{Expectabo}

\begin{psalmsheet}[tone=4, continuous, number={51. zsoltár}, title={A rágalmazók ellen}, motto={Aki dicsekszik, az Úrban dicsekedjék (1 Kor 1, 31).}]
    \psalmtext{
        3. Mit dicsekszel gonoszságoddal, * te, aki a gonoszságban vagy hatalmas?\\
        4. Egész nap álnokságon jár az eszed; * te cselszövő, a nyelved éles, mint a borotva!\\
        5. A rosszat jobban kedveled a jónál, † a hazugságot inkább, mint az igaz beszédet, * te álnok nyelv, minden bajt hozó beszéd kedves neked. ---

        7. Éppen ezért el is pusztít Isten mindörökre, † kigyomlál és kiköltöztet hajlékodból, * gyökerestől kitép az élők földjéből.\\
        8. Látják ezt az igazak, és megilletődnek, * azután nevetve mondják:\\
        9. „Íme, az ember, aki nem Istenben kereste oltalmát, † hanem gazdagságában volt reménye, * és csak gonoszsága minden ereje!” ---

        10. Én pedig mint viruló olajfa Isten házában, † Isten irgalmában bízom, * mindenkor és mindörökké.\\
        11. Örökké áldalak azért, amit tettél, † nevedben reménykedem, * mert jóságos vagy híveidhez.
}
\end{psalmsheet}

\versicle{Kitartok igéd mell\underline{e}tt.}{Az Úr minden remény\underline{e}m.}

\parttitle[Vö. Judit 9, 18; 1 Krón 21, 15; Sám 24, 17]{Válaszos ének}

\versicle{%
    Emlékezzél, Uram, szövetségedre, és mondd pusztító \underline{a}ngyal\underline{o}dnak: † Vond
    vissza \underline{a} k\underline{e}zed, * Hogy ne pusztuljon el a föld, és ne veszítsd el az
    él\underline{ő}lény\underline{e}ket.
}{%
    Én vétkeztem, én követtem el gonoszs\underline{á}got! † Ezek, a nyáj mit vét\underline{e}ttek? * Ezért
    hát kérlek, forduljon el haragod népedt\underline{ő}l, \underline{U}ram.
}

\parttitle[1 Kor 10, 16-17]{Válaszos ének}

\versicle{%
    Az áldás kelyhe, amely\underline{e}t meg\underline{á}ldunk, † nemde a Krisztus vérében való
    rész\underline{e}s\underline{e}dés? * S a kenyér, amelyet megtörünk, nemde a Krisztus testében
    való r\underline{é}szesedés?
}{%
    M\underline{i} ugyanis sokan egy kenyér és egy test v\underline{a}gyunk, * mivel mindnyájan egy
    kenyérből rész\underline{e}s\underline{ü}lünk.
}

\rubric{Könyörgés, mint vasárnap, a Reggeli dicséretben, \pageref{DXIV/L/Coll}.}

\subsubsection*{Reggeli dicséret}

\parttitle{Himnusz}

\printsheet{FulgentisAuctor}

\begin{hymnstanzas}
    Ím, szertefoszlik már az éj,\\
    a nagy világ kel, újra él,\\
    lelkünk is új erőre kap,\\
    s belőle sok jótett fakad.
\stanzabreak
    A felkelő nap int nekünk:\\
    téged dicsérjen énekünk,\\
    a fényes égnek arca már\\
    szívünk vidámságára vár.
\stanzabreak
    Kerüljük el, mi léhaság,\\
    lelkünk ne tűrjön bűnt, hibát,\\
    és életünk ne érje szenny,\\
    nyelvünk ne szóljon vétkesen.
\stanzabreak
    De míg a nap nyugodni száll,\\
    a forró, mély hit járjon át,\\
    reményünk várjon Életet,\\
    s fűzzön Krisztushoz szeretet.
\end{hymnstanzas}
\centerstanza{%
    Ezt add meg, kérünk, jó Atya,\\
    s Atyának mása, egy Fia,\\
    s te, Szentlélek, Vigasztaló:\\
    örökre egy uralkodó. Ámen.
}

\parttitle{Zsoltározás}

\redtext{1. antifóna}

\printsheet{DeusInSanctoViaTua}

\begin{psalmsheet}[tone=8, continuous, number={76. zsoltár}, title={Emlékezés az Úr nagy tetteire}, motto={Mindenfelől szorongatnak, de össze nem zúznak (2 Kor 4, 8).}]
    \psalmtext{
        2. Hangos szóval az Úrhoz kiáltok, * Istenhez kiáltok, és ő figyel reám.\\
        3. Keresem az Urat a szorongatás napján, † egész éjjel felé tárom kezemet, * és el nem lankadok.\\
        4. Lelkem nem tud megvigasztalódni; † Istenre gondolok, és sóhajtozom, * róla elmélkedem, és elcsügged a lelkem. ---

        5. Te tartod ébren szempilláimat, * zavart vagyok, szólni sem tudok.\\
        6. Emlékezem a régi napokról, * gondolkodom a rég letűnt esztendőkről.\\
        7. Egész éjjel elmélkedik szívem, * kutatok, és fürkészi a lelkem: ---

        8. Vajon mindörökre eltaszít az Úr, * sosem lesz már megbékéltebb?\\
        9. Vagy szeretete végleg elfogyott, * és ígérete nem valósul meg soha?\\
        10. Könyörülni is elfelejt talán az Isten, * és haragjában irgalmasságát szívébe zárja? ---

        11. És azt mondtam: „Onnét van sajgó sebem, * hogy az Úr keze most másként érint!”\\
        12. Emlékezem az Úr nagy tetteiről, * emlékezem kezdettől adott csodáiról.\\
        13. Minden művedről elmélkedem, * nagyszerű tetteidről gondolkodom. ---

        14. Isten, szent a te utad, * ki olyan nagy isten, mint a mi Istenünk?\\
        15. Te vagy az Isten, aki csodákat tehet, * megismertetted hatalmadat a népek között.\\
        16. Erős karral megváltottad népedet, * Jákob és József fiait. ---

        17. Megláttak téged a vizek, Istenem, † megláttak a vizek, és megremegtek, * a mély vizek megrendültek.\\
        18. A felhők ontották a vizet, † a fellegek mennydörögtek, * mert tüzes nyilaid cikáztak.\\
        19. Mennydörgésed harci szekérként dübörgött, † villámaid bevillogták a földkerekséget, * megingott és megreszketett a föld. ---

        20. Utad a tengeren át vezetett, ösvényed a nagy vizeken, * de lábad nyomát nem láthatta senki.\\
        21. Népedet, mint nyájat, úgy vezetted * Mózes és Áron kezével.
}
\end{psalmsheet}

\redtext{2. antifóna}

\parttitle[1 Sám 2, 1-10]{Kantikum}

\printsheet{Exsultavit}

\begin{psalmsheet}[tone=8, continuous, title={Az alázatosak örvendezése Istenben}, motto={Hatalmasokat ledöntött a trónról, és alázatosakat felmagasztalt, éhezőket minden jóval betöltött (Lk 1, 52-53).}]
    \psalmtext{
        1. Szívem ujjong az Úrban, * erővel tölt el Istenem.\\
        1. Ellenségeimmel szemben szám tágasra nyílhat, * mert örömmel tölt el segítséged. ---

        2. Nincs olyan szent senki, mint az Úr, † nincs senki más rajtad kívül, * nincs senki olyan erős, mint a mi Istenünk.\\
        3. Ne beszéljetek annyit gőgösen, * és ne büszkélkedjetek.\\
        3. Ne hagyja el szátokat kérkedő szó, † az Úr a mindentudás Istene, * ő méri le tetteinket. ---

        4. A hatalmasok íja eltörött, * de a gyengéket erő övezi.\\
        5. Aki jóllakott volt, elszegődik egy falat kenyérért, * de akik éheztek, már jóllakhatnak.\\
        5. Aki gyermektelen volt, az hétszeresen anya lesz, * a sokgyermekes anya meg elhervad.\\
        6. Az Úr adja a halált, és ő adja az életet, * alvilágra ő küld, és ő hoz fel onnan.\\
        7. Az Úr ad szegénységet, és ő ad gazdagságot, * ő az, aki megaláz és felmagasztal.\\
        8. A szűkölködőt ő emeli fel a porból, * és a sárból kihozza a nyomorultat.\\
        8. Az előkelők sorába ülteti, * hogy övé legyen a dicsőség széke. ---

        8. Mert az Úré a föld minden oszlopa, * azokra helyezte a földkerekséget.\\
        9. Szentjeinek lábát biztos úton vezeti, † a gonoszok elvesznek a sötétségben, * mert az ember nem győzhet a maga erejéből.\\
        10. Az Úr megfélemlíti ellenségeit, * az égből rájuk mennydörög.\\
        10. Az Úr megítéli a föld határait, † királyának ő ad hatalmat, * és Fölkentjének homlokát fölemeli!
}
\end{psalmsheet}

\redtext{3. antifóna}

\printsheet{DominusRegnavit}

\begin{psalmsheet}[tone=3, continuous, number={96. zsoltár}, title={Isten dicsőségesen ítélkezik}, motto={Ez a zsoltár a világ megváltását jelzi, és minden nép hitét Istenben (Szt. Atanáz).}]
    \psalmtext{
        1. Király az Úr, ujjongjon a föld, * † és örvendjen minden sziget!\\
        2. Sötét felhő lebegi őt körül, * trónját igazság és jog tartja szilárdan.\\
        3. Tűz jár színe előtt, * és elemészti minden ellenségét. ---

        4. Villámai bevilágítják a földkerekséget, * látja a föld, és beléremeg.\\
        5. Viaszként szétfolynak a hegyek az Úr előtt, * az egész föld Urának színe előtt.\\
        6. Igazságosságát hirdetik az egek, * és nagy fölségét meglátják a népek. ---

        7. Megszégyenülnek mind a bálványimádók, † akik faragott képeikkel dicsekszenek; * eléje borulnak mind a bálványistenek.\\
        8. Hallja ezt, és örvendez Sion, † Júda városai vigadoznak * a te ítéleted miatt, Uram.\\
        9. Mert magasságos Úr vagy te a világ fölött, * minden más istennél fölségesebb. ---

        10. Akik az Urat szeretitek, gyűlöljétek a rosszat! † Az Úr megoltalmazza hívei lelkét, * a gonoszok kezéből kimenti őket.\\
        11. Fényesség virrad az igazra, * és öröm a tiszta szívűekre.\\
        12. Örvendjetek, igazak, az Úrban, * és áldjátok szentséges nevét!
}
\end{psalmsheet}

\parttitle[Róm 8, 35. 37]{Rövid olvasmány}
Ki szakíthat el bennünket Krisztus szeretetét\underline{ő}l? † Nyomor vagy szükség? Üldözte\-tés vagy éhínség, ruhátlanság, életv\underline{e}szély vagy kard? * De mindezeken diadalmaskodunk őáltala, aki szer\underline{e}t m\underline{i}nket.

\parttitle{Rövid válaszos ének}

\printsheet{BenedicamDominum}

\parttitle[Lk 1,68-79]{Evangéliumi kantikum}

\printsheet{Liberati}

\parttitle{Fohászok}

\begin{preces}{Áldott legyen üdvözítő Istenünk, aki megígérte, hogy a világ végéig mindennap Egyházával marad. Adjunk hálát neki, és kérjük:}{Maradj velünk, Urunk!}
  Maradj velünk, Urunk, egész nap,

  \reddash sohase nyugodjék le fölöttünk kegyelmednek napja!

  Mai napunkat neked szenteljük,

  \reddash ígérjük: semmi rosszat nem teszünk és nem helyeslünk.

  Add, Urunk, hogy egész napunk folyamán téged sugározzunk,

  \reddash a föld sója és a világ világossága legyünk!

  Szentlelked irányítsa szívünket és ajkunkat,

  \reddash hogy mindig megmaradjunk igazságodban és dicséretedben.
\end{preces}
Mi Atyánk...

\parttitle{Könyörgés}

Istenünk, küldd ragyogó világosságodat szívünkbe, hogy mindig parancsaid útján járjunk, és mentve maradjunk minden eltévelyedéstől. A mi Urunk.

\subsubsection*{Befejező imaóra}

\parttitle{Zsoltározás}

\redtext{1. antifóna}

\printsheet{EstoMihi}

\begin{psalmsheet}[tone=6, continuous, number={30. zsoltár, 2-6}, title={Üldözött ember bizakodó imádsága}, motto={Atyám, kezedbe ajánlom lelkemet (Lk 23, 46).}]
    \psalmtext{
        2. Te vagy, Uram, én reményem, † ne hagyj soha szégyent érnem, * igazságodban szabadíts meg engem!\\
        3. Fordítsd felém füledet, * siess, Uram, ments meg engem!\\
        3. Te légy oltalmazó kősziklám, és megerősített házam, * hogy megszabadíts engem! ---

        4. Erősségem és menedékem valóban te vagy, * neved miatt vezess és irányíts engem!\\
        5. Kiszabadítasz a hálóból, amit titkon vetettek nekem, * mivel te vagy az én erősségem! ---

        6. Uram, kezedbe ajánlom lelkemet; * megváltottál engem, hűséges Istenem!
}
\end{psalmsheet}

\redtext{2. antifóna}

\printsheet{DeProfundis}

\begin{psalmsheet}[tone=8, continuous, number={129. zsoltár}, title={A mélységből kiáltok}, motto={Ő szabadítja meg népét bűneitől (Mt 1, 21).}]
    \psalmtext{
        1. A mélységből kiáltok, Uram, hozzád, * † Istenem, figyelj a szómra!\\
        2. Fordulj felém, és hallgass figyelmesen * könyörgésem hangos szavára! ---

        3. Ha bűneinket, Uram, felrovod, * Uram, ki az, aki előtted megállhat?\\
        4. De nálad bocsánatot nyer a vétek, * ezért félve tisztelünk téged.\\
        5. Remélek az Úrban; † benne remél a lelkem, * és bízom ígéretében.\\
        6. Lelkem várja az Urat, * jobban, mint éji őr a hajnalt. ---

        6. Bizony, jobban, mint éji őr a hajnalt, * várja Izrael az Urat,\\
        7. mert az Úrnál van az irgalom, * és bőséges a megváltás nála.\\
        8. Választott népét ő váltja meg * minden bűnétől.
}
\end{psalmsheet}

\parttitle[Ef 4, 26-27]{Rövid olvasmány}
Ne vétkezzet\underline{e}k! † A nap ne nyugodjék le harag\underline{o}tok fölött! * Ne adjatok teret a s\underline{á}t\underline{á}nnak!

\parttitle{Könyörgés}
Szelíd és alázatos szívű Urunk, Jézus Krisztus, aki édes igát és könnyű terhet raksz követőidre, fogadd jóságosan mai napunk jószándékait és tetteit. Adj nyugalmas éjszakát, hogy holnap készségesebbek legyünk szolgálatodra. Aki élsz és uralkodol mindörökkön-örökké.

\subsection*{Csütörtök, 14. hét - 1. nap}

\subsubsection*{Imádságra hívás}

\printsheet{AdoremusDominum}

\subsubsection*{Olvasmányos imaóra}

\parttitle{Himnusz}

\printsheet{AlesDieiNuntius}

\begin{hymnstanzas}
    „Hagyjátok már az ágyakat,\\
    ti ernyedt, kába tétlenek;\\
    szilárdan, tisztán, éberen\\
    virrasszatok; közel vagyok!”
\stanzabreak
    Hogy fényes szellők szárnyain\\
    ha kél a hajnal, s gyúl az ég,\\
    munkára készen kezdjük el\\
    az új napot s az új reményt,
\stanzabreak
    Jézust kiáltsuk hangosan,\\
    sírván, könyörgőn, éberen;\\
    imádkozzunk figyelmesen,\\
    s a tiszta szív nem szunnyad el.
\stanzabreak
    Álmunkat, Krisztus, messze űzd,\\
    törd szét az éj bilincseit,\\
    oldozd fel mind a régi bűnt,\\
    s új fényességet önts belénk!
\end{hymnstanzas}
\centerstanza{%
    Dicsérjük az örök Atyát,\\
    dicsérjük egyszülött Fiát\\
    s a Lelket, a Vigasztalót,\\
    időtlen századok során! Ámen.
}

\parttitle{Zsoltározás}

\redtext{1. antifóna}

\printsheet{Extende}

\begin{psalmsheet}[tone=8, continuous, number={43. zsoltár}, title={A nép sorscsapásai}, motto={Mindezeken diadalmaskodunk őáltala, aki szeret minket (Róm 8, 37).}, numeral={I.}]
    \psalmtext{
        2. Istenünk, hallottuk saját fülünkkel, † atyáink elbeszélték nekünk a csodákat, * amelyeket napjaikban tettél, egykor régen.\\
        3. Kezed népeket űzött el, őket meg letelepítetted, * nemzeteket zúztál össze, őket pedig megnövelted. ---

        4. Mert nem kardjukkal vették birtokba ezt a földet, * és nem is az ő karjuk segítette őket,\\
        4. hanem a te jobbod és a te karod és arcod fényessége, * mert úgy szeretted őket.\\
        5. Te vagy Istenem és Királyom, * parancsodra megszabadul Jákob.\\
        6. Veled szórtuk szét ellenségeinket; * a te nevedben tiportuk le azokat, akik ellenünk támadtak. ---

        7. Mert én nem az íjamban bízom, * kardom meg nem szabadíthat.\\
        8. De te megszabadítasz minket sanyargatóinktól, * és megszégyeníted ellenségeinket.\\
        9. Mindennap Isten a mi dicsekvésünk, * és nevedet mindörökké magasztaljuk!
}
\end{psalmsheet}

\redtext{2. antifóna}

\printsheet{Confundantur}

\begin{psalmsheet}[tone=2, continuous, numeral={II.}]
    \psalmtext{
        10. De most elvetettél és megaláztál minket, * Istenünk, már nem vonulsz ki együtt seregünkkel.\\
        11. Megfutamítottál ellenségeink elől, * és kifosztottak, akik gyűlölnek minket.\\
        12. Odaadtál, mint vágásra szánt juhokat, * a nemzetek között széjjelszórtál. ---

        13. Eladtad népedet haszon nélkül, * és cserébe érte nem adtak sokat.\\
        14. Szomszédaink előtt gyalázattá tettél minket, * gúnyol és csúfol környezetünk.\\
        15. Szállóigévé tettél a pogányok számára, * csak csóválják fejüket a népek. ---

        16. Gyalázatom szemem előtt van egész nap, * és szégyenpírban ég az arcom\\
        17. a gúnyolók és szemrehányók szava miatt, * ellenségem és támadóm színe előtt.
}
\end{psalmsheet}

\redtext{3. antifóna}

\printsheet{MementoDomine}

\begin{psalmsheet}[tone=2, continuous, numeral={III.}]
    \psalmtext{
        18. Mindezek most ránk szakadtak, † téged mégsem feledtünk el, * és nem szegtük meg szövetségedet.\\
        19. Szívünk nem fordult el tőled, * lépteink utadról le nem tértek,\\
        20. pedig megaláztál a sakálok földjén minket, * ránk borítottad a halál árnyát. ---

        21. Ha Istenünk nevét elfelednénk, * és ha kezünket más isten felé tárnánk,\\
        22. Isten vajon ezt nem kérné számon? * A szív titkait csak ő ismeri.\\
        23. Teérted gyilkolnak mindennap bennünket, * és leölésre szánt juhoknak tekintenek! ---

        24. Kelj fel, miért alszol, Uram, * kelj fel, és ne taszíts el mindörökre!\\
        25. Miért fordítod el arcodat? * Elfeleded ínségünket és gyötrelmeinket? ---

        26. A föld porába alázták lelkünket, * a földhöz tapadt egész testünk.\\
        26. Kelj fel, Uram, siess segítségünkre, * irgalmasságodban válts meg minket!
}
\end{psalmsheet}

\versicle{Uram, kihez menn\underline{é}nk?}{Tiéd az örök életet adó tanít\underline{á}s.}

\parttitle[1 Krón 22, 19; Zsolt 131, 7; Iz 56, 7]{Válaszos ének}

\versicle{%
    R\underline{a}jta, † egész szívvel és lélekkel építsetek szentélyt Istennek, \underline{a}z \underline{Ú}rnak; * Induljunk hát sátorába, boruljunk le lábának zsám\underline{o}lya \underline{e}lőtt!
}{%
    \underline{E}zt mondja \underline{a}z Úr: * házam minden nép számára az imádság h\underline{á}z\underline{a} lesz. 
}

\parttitle[Jel 3, 20; Mt 24, 46]{Válaszos ének}

\versicle{%
    Nézd, az ajtóban állok, \underline{é}s kop\underline{o}gok. † Aki meghallja szavam, és \underline{a}jt\underline{ó}t nyit, * Bemegyek hozzá, vele eszem, ő m\underline{e}g énv\underline{e}lem.
}{%
    Boldog a sz\underline{o}lga, * ha ura hazatérve ilyen munkában t\underline{a}l\underline{á}lja.
}

\rubric{Könyrgés, mint vasárnap, a Reggeli dicséretben, \pageref{DXIV/L/Coll}.}

\subsubsection*{Reggeli dicséret}

\parttitle{Himnusz}

\printsheet{IamLucisOrto}

\begin{hymnstanzas}
    Nyelvünket fogja fékre ma:\\
    ne szóljon rút perek szava;\\
    szemünket védőn óvja meg:\\
    a hívságot ne lássa meg!
\stanzabreak
    Lakjék szívünkben tisztaság,\\
    távozzék minden dőreség;\\
    a testnek dölyfét törje meg\\
    étel- s italban hősi fék,
\stanzabreak
    hogy majd a nap ha távozott,\\
    s az óra újra éjt hozott,\\
    lemondásunk szent éneke\\
    legyen az Úr dicsérete.
\stanzabreak
    Dicsérjük az örök Atyát,\\
    dicsérjük egyszülött Fiát\\
    s a Lelket, a Vigasztalót,\\
    most és örök időkön át! Ámen.\\
\end{hymnstanzas}

\parttitle{Zsoltározás}

\redtext{1. antifóna}

\printsheet{TuamDomine}

\begin{psalmsheet}[tone=2, continuous, number={79. zsoltár}, title={Látogasd meg szőlődet, Uram!}, motto={Jöjj el, Uram, Jézus! (Jel 22, 20)}]
    \psalmtext{
        2. Figyelj ránk, Izrael pásztora, * aki nyájadként vezeted József népét.\\
        3. Aki a kerubokon trónolsz, jelenj meg ragyogva * Efraim, Benjamin és Manassze törzse előtt.\\
        3. Mutasd meg hatalmadat, és jöjj, * hogy megszabadíts minket. ---

        4. Újíts meg minket, Istenünk, * ragyogtasd ránk arcodat, és szabadok leszünk! ---

        5. Urunk, mindenható Isten, * meddig haragszol még néped imádsága ellenére?\\
        6. Könnyek kenyerével táplálsz minket, * és italul könnyeink árját adtad nekünk.\\
        7. Szomszédaink előtt csúffá lettünk, * ellenségeink gúnyolódnak rajtunk.\\
        8. Újíts meg minket, mindenható Isten, * ragyogtasd ránk arcodat, és szabadok leszünk! ---

        9. Egyiptomból telepítetted át szőlődet, * népeket űztél el, hogy elültethesd.\\
        10. Helyét jól megtisztítottad, * gyökeret vert, és betöltötte már a földet.\\
        11. Beborítja árnyéka a hegyeket, * ágai magasabbak a hatalmas cédrusoknál.\\
        12. Indái a tengerig elértek, * hajtásai meg a Folyamig. ---

        13. Miért romboltad le kerítését, * hogy akik arra járnak, mind leszüreteljék? ---

        14. Elpusztította az erdei vadkan, * és a magányos vad lelegelte.\\
        15. Újíts meg minket, mindenható Isten, * tekints le ránk, lásd, és látogasd meg ezt a szőlőt! ---

        16. Védd meg, amit jobbod ültetett, * és az emberfiát, akit magadhoz emeltél.\\
        17. Tűzzel égetik és feldúlják, * de arcod feddésétől elpusztulnak.\\
        18. Jobbod embere fölé tárd ki kezedet, * az emberfia fölé, akit magadhoz emeltél.\\
        19. Nem hagyunk el többé, éltess még bennünket, * és mi segítségül hívjuk nevedet. ---

        20. Újíts meg minket, Urunk, mindenható Istenünk, * ragyogtasd ránk arcodat, és szabadok leszünk!
}
\end{psalmsheet}

\redtext{2. antifóna}

\parttitle[Iz 12, 1-6]{Kantikum}

\printsheet{ConversusEst}

\begin{psalmsheet}[tone=8, continuous, title={A megszabadított nép hálaéneke}, motto={Aki szomjazik, jöjjön hozzám, és igyék! (Jn 7, 37)}]
    \psalmtext{
        1. Hálát adok neked, Uram,mert megharagudtál ugyan reám, * de elfordult haragod, és megvigasztaltál.\\
        2. Íme, Isten az én szabadítóm, * bizalom tölt el, és nem félek,\\
        2. mert az Úr az én erőm és dicsőségem, * ő lett nékem szabadulásom. ---

        3. Víg örömmel merítetek vizet * az üdvösség forrásaiból.\\
        4. Így szóltok majd azon a napon: † „Adjatok hálát az Úrnak, * és magasztaljátok nevét!\\
        4. Mondjátok el a népek közt tetteit, * és hirdessétek, hogy az ő neve fölséges.\\
        5. Zengjetek az Úrnak, mert nagyszerű, amit cselekedett; * híreszteljétek az egész földön.\\
        6. Ujjongj, és mondj dicséretet, Sion népe, * mert nagy közötted Izrael Szentje.”
}
\end{psalmsheet}

\redtext{3. antifóna}

\printsheet{ExsultateDeo}

\begin{psalmsheet}[tone=1, continuous, number={80. zsoltár}, title={A szövetség ünnepélyes megújítása}, motto={Vigyázzatok, hogy ne legyen hitetlenségre hajló, gonosz a szívetek (Zsid 3, 12).}]
    \psalmtext{
        2. Örvendjetek Istennek, a mi segítőnknek, * ujjongjatok Jákob Istenének!\\
        3. Csendüljön a zsoltár, hozzatok dobot, * és lágy hangú hárfát citerával!\\
        4. Fújjátok meg a harsonát újholdkor, * és holdtöltekor, a mi nagy ünnepünkön! ---

        5. Izraelben ez a törvény, * ezt így rendelte Jákob Istene.\\
        6. Ő hagyta meg ezt Józsefnek, * amikor Egyiptom földjéről kivonult. ---

        7. Beszédet hallottam, amit nem ismertem: † „Válláról levettem a terhet, * keze megszabadult a kosártól.\\
        8. Gyötrelmedben segítségül hívtál, és megszabadítottalak, † a vihar leple mögül meghallgattalak, * próbára tettelek Meriba vizénél. ---

        9. Népem, hallgass rám, hadd intselek, * Izrael, bárcsak figyelnél szavamra:\\
        10. Ne legyen soha más istened, * ne imádj soha idegen istent!\\
        11. Mert én vagyok a te Urad, Istened, † aki kivezetett téged Egyiptomból, * csak nyisd ki szádat, és én betöltöm. ---

        12. De népem nem hallgatta meg szavamat, * és nem figyelt rám Izrael.\\
        13. Szívük keménységére hagytam őket, * járjanak csak kedvteléseik szerint. ---

        14. Bárcsak hallgatna rám népem, * és Izrael követné utaimat,\\
        15. ellenségeit hamarosan megaláznám, * és gyötrőire emelném kezemet. ---

        16. Akik gyűlölik az Urat, meghódolnának neki, * és sorsuk örökké tartana.\\
        17. Népemet pedig a búza legjavából táplálnám, * és jóllakatnám őket mézzel a sziklából.”
}
\end{psalmsheet}

\parttitle[Róm 14, 17-19]{Rövid olvasmány}
Az Isten országa nem eszem-iszom, hanem igazságosság, béke és öröm a Szentlélekb\underline{e}n. † Aki így szolgál Krisztusnak, az kedves Isten előtt és rokonszenves az \underline{e}mbereknek. * Ezért arra törekedjünk, ami a békességre és kölcsönös épülésünkr\underline{e} sz\underline{o}lgál.

\parttitle{Rövid válaszos ének}

\printsheet{InMatutinisDomine}

\parttitle[Lk 1,68-79]{Evangéliumi kantikum}

\printsheet{DaScientiam}

\parttitle{Fohászok}

\begin{preces}{Áldott legyen Istenünk, Atyánk; ő megoltalmazza gyermekeit és nem veti meg kéréseiket. Alázatos könyörgéssel kérjük őt valamennyien:}{Adj világosságot szemünknek, Urunk!}
  Köszönjük, Urunk, hogy Fiad által elküldted nekünk a világosságot,

  \reddash add, hogy egész nap bőven áradjon ránk az ő fényessége!

  Urunk, a te bölcsességed vezessen a mai napon minket,

  \reddash hogy megújult lélekkel járjunk-keljünk!

  Add, hogy érted a kellemetlenségeket is erős lélekkel viseljük,

  \reddash és mindig nagylelkűen szolgáljunk neked!

  Irányítsd ma gondolatainkat, érzéseinket és cselekedeteinket,

  \reddash hogy tanulékony lélekkel fogadjuk gondviselő intézkedéseidet!
\end{preces}
Mi Atyánk...

\parttitle{Könyörgés}

Istenünk, igazi Világosság és világosság Alkotója, esedezve kérünk add, hogy hívő lélekkel gondolkodjunk mindarról, ami szent, és mindenkor a te fényességedben éljünk. A mi Urunk.

\subsubsection*{Befejező imaóra}

\parttitle{Zsoltározás}

\redtext{Antifóna}

\printsheet{CaroMea}

\begin{psalmsheet}[tone=4, continuous, number={15. zsoltár}, title={Az Úr az én örökrészem}, motto={Isten feloldotta a halál bilincseit, és feltámasztotta Jézust (ApCsel 2, 24).}]
    \psalmtext{
        1. Oltalmazz engem, Istenem, * mert én benned bízom. ---

        2. Így szólok az Úrhoz: „Istenem vagy nekem, * nincs más javam rajtad kívül.”\\
        3. A dicső és szent férfiakhoz, akik földünkön élnek, * vonzódom egész szívemmel.\\
        4. Bajt bajra tetéznek maguknak, * akik idegen istenek után futnak.\\
        4. Véres áldozataikban nincsen részem, * még csak nevüket sem veszem ajkamra. ---

        5.Az Úr az én örökrészem és kelyhem, * te irányítod, Uram, sorsomat.\\
        6.A mérőzsinór értékes részt juttatott nekem, * valóban pompás az én örökrészem.\\
        7.Áldom az Urat, mert tanácsot ad nekem, * még éjszaka is a jóra int engem bensőm.\\
        8.Szüntelenül magam előtt látom az Urat, * ő áll jobbomon, hogy el ne essem.\\
        9.Ezért örül az én szívem, † és ujjong a lelkem, * sőt reménységben nyugszik testem is. ---

        10. Mert nem hagyod lelkemet a holtak honában, * és nem engeded, hogy pusztulást lásson szented.\\
        11. Az élet útját mutatod nekem, † örömmel töltesz el színed előtt, * és mindörökké tart az öröm jobbodon!
}
\end{psalmsheet}

\parttitle[1 Tessz 5, 23]{Rövid olvasmány}
A békesség Istene szenteljen meg bennetek\underline{e}t, † hogy tökélet\underline{e}sek legyetek. * Őrizze meg szellemeteket, lelketeket és testeteket feddhetetlenül Urunk, Jézus Krisztus eljövet\underline{e}l\underline{é}ig!

\parttitle{Könyörgés}
Urunk, Istenünk, ajándékozz nekünk pihentető álmot, mert elfáradtunk a napi munkában. Újíts meg minket állandó segítségeddel, hogy testben-lélekben a tieid maradjunk. Krisztus, a mi Urunk által.

\subsection*{Péntek, 14. hét - 2. nap}

\subsubsection*{Imádságra hívás}

\printsheet{DominumDeumNostrum}

\subsubsection*{Olvasmányos imaóra}

\parttitle{Himnusz}

\printsheet{GalliCantu}

\begin{hymnstanzas}
    Erős, éber Őrként jöttél,\\
    embereknek fényeként,\\
    mikor földön nagy csend honolt,\\
    s ráborult a néma éj,\\
    a halandó emberi nem\\
    halálának jeleként.
\stanzabreak
    Kérünk, Krisztus, ébressz minket,\\
    űzd el gonosz álmaink,\\
    és az éj e börtönéből\\
    kegyelmeddel szabadíts;\\
    add vissza a fényt szemünknek,\\
    életutunk társa légy.
\end{hymnstanzas}
\centerstanza{%
    Áldunk téged és Atyádat,\\
    s áldjuk együtt Lelkedet,\\
    a Hármas-egy Istenséget,\\
    békét, fényt és életet,\\
    s neve minden másnál drágább,\\
    minden méznél édesebb! Ámen.
}

\parttitle{Zsoltározás}

\redtext{1. antifóna}

\printsheet{NeInIraTua}

\begin{psalmsheet}[tone=8, continuous, number={37. zsoltár}, title={Végveszélybe került bűnös könyörgése}, motto={Jézus ismerősei távolabb álltak (Lk 23, 49).}, numeral={I.}]
    \psalmtext{
        2. Uram, ne feddj meg felindulásodban, * és ne büntess haragodban,\\
        3. mert nyilaid belém hatoltak, * és rám nehezedett kezed. ---

        4. Nincs testemen ép hely, mert haragszol énrám, * nincs nyugta csontjaimnak, mert sok a vétkem.\\
        5. Gonoszságaim fejemre gyűltek, * súlyos teherként nehezednek rám. ---
}
\end{psalmsheet}

\redtext{2. antifóna}

\printsheet{ConfundanturEtRevereantrur}

\begin{psalmsheet}[tone=2, continuous, numeral={II.}]
    \psalmtext{
        6. Üszkösödnek gennyedt sebeim * oktalanságom miatt.\\
        7. Nyomorúságomban meggörnyedtem, * naphosszat szomorúan járok-kelek. ---

        8. Ágyékomat gyulladás gyötri, * nincs testemben semmi épség.\\
        9. Elgyöngültem, és levert vagyok, * keserűen jajgatok szívem fájdalmában. ---

        10. Uram, előtted van minden kívánságom, * nincs elrejtve előtted sóhajom.\\
        11. Reszket a szívem, elhagy az erőm, * elhomályosul szemem világa. ---

        12. Barátaim és szeretteim dermedten állnak bajom előtt, * rokonaim messze maradnak. ---

        13. Csapdát állítanak, akik vesztemre törnek, † rosszakaróim álnokul szólnak, * és egész nap cselt szőnek ellenem.
}
\end{psalmsheet}

\redtext{3. antifóna}

\printsheet{NeDerelinquasMe}

\begin{psalmsheet}[tone=8, continuous, numeral={III.}]
    \psalmtext{
        14. Én meg, mint egy süket, semmit sem hallok, * mint egy néma, nem nyitom föl számat.\\
        15. Olyan vagyok, mint aki nem hall, * mint aki nem tud válaszolni. ---

        16. Mert, Uram, én benned remélek, * meg is hallgatsz, Uram, Istenem.\\
        17. Így szóltam ugyanis: „Rajtam senki ne nevessen, * és ha lábam tántorul, senki se dicsekedjék!” ---

        18. Mert már a bukás küszöbén állok, * a szenvedés nem marad el tőlem.\\
        19. Ezért gonoszságomat őszintén megvallom, * és bánkódom bűneim miatt. ---

        20. Ellenségeim azonban erősebbek nálam, * és sokan vannak, akik gonoszul gyűlölnek engem.\\
        21. Rágalmaznak, és a jóért rosszal fizetnek, * mivel mindig a jót akarom. ---

        22. Ne hagyj magamra, én Uram, * Istenem, ne távozz el tőlem!\\
        23. Siess, oltalmazz engem, * Uram, én menedékem!
}
\end{psalmsheet}

\versicle{Pilláim elernyednek, míg segítségedet vár\underline{o}m,}{És igazságodat, melyet megígért\underline{é}l.}

\parttitle[Én 3, 11; Zsolt 71, 1. 2]{Válaszos ének}

\versicle{%
    Sion leányai, gyert\underline{e}k, nézz\underline{é}tek: † Salamon király, fején a korona, amellyel szülőanyja megkor\underline{o}n\underline{á}zta * Szíve öröm\underline{é}nek n\underline{a}pján!
}{%
    Istenem, ítéletedet add a kir\underline{á}lynak, * hogy méltányosan ítélkezzék szegény\underline{ei}den. 
}

\parttitle[1 Jn 4, 21; Mt 22, 40]{Válaszos ének}

\versicle{%
    Ezt \underline{a} par\underline{a}ncsot † kaptuk \underline{I}st\underline{e}ntől: * Aki Istent szereti, szeresse t\underline{e}stvér\underline{é}t is.
}{%
    Ezen a két parancson al\underline{a}pszik * az egész törvény és a pr\underline{ó}f\underline{é}ták.
}

\rubric{Könyörgés, mint vasárnap, a Reggeli dicséretben, \pageref{DXIV/L/Coll}.}

\subsubsection*{Reggeli dicséret}

\parttitle{Himnusz}

\printsheet{DeusQuiCaeliLumen}

\begin{hymnstanzas}
    hajnal jön, tűnnek csillagok,\\
    nyomukban bíbor hab fakad,\\
    s az égi Szél gyöngyharmata\\
    új keresztvízként földre hull.\\
\stanzabreak
    Eltűnik már az éjsötét,\\
    az ég homálya messze száll,\\
    s Krisztust jelezve kelti már\\
    a hajnalcsillag a napot.\\
\stanzabreak
    Napoknak Napja, Istenünk,\\
    kinek fényében lát szemünk,\\
    egyedül vagy mindenható,\\
    Szentháromságban egy való.\\
\stanzabreak
    Megváltónk, buzgón kérlelünk,\\
    kik most elődbe térdelünk:\\
    Atyádat és Szentlelkedet\\
    hadd áldjuk szívből teveled. Ámen.\\
\end{hymnstanzas}

\parttitle{Zsoltározás}

\redtext{1. antifóna}

\printsheet{AmpliusLavaMe}

\begin{psalmsheet}[tone=1, continuous, number={50. zsoltár}, title={Könyörülj rajtam, Istenem!}, motto={Újuljatok meg lélekben és érzületben, s öltsétek magatokra az új embert! (Ef 4, 23-24)}]
    \psalmtext{
        3. Könyörülj rajtam, Istenem, * szerető jóságod szerint,\\
        3. és minthogy mérhetetlen a te irgalmad, * gonoszságomat töröld le rólam!\\
        4. Mosd le bűnömet teljesen, * és vétkemtől tisztíts meg engem! ---

        5. Gonoszságomat belátom, * bűnöm előttem lebeg szüntelen.\\
        6. Egyedül csak ellened vétkeztem, * ami színed előtt gonosz, olyat tettem.\\
        6. Te feddhetetlen vagy döntéseidben, * és igazságos, amikor ítélsz. ---

        7. Íme, én bűnben születtem, * anyám már vétekben fogant engem.\\
        8. Te pedig a szívben élő hűséget kedveled, * és kinyilatkoztattad nekem bölcsességed titkait. ---

        9. Hints meg izsóppal, és tiszta leszek, * moss meg, és fehérebb leszek, mint a hó!\\
        10. Add, hogy vidámságot és örömet halljak, * hadd örvendezzenek megtört csontjaim! ---

        11. Fordítsd el bűneimtől arcodat, * és töröld el minden vétkem.\\
        12. Tiszta szívet teremts bennem, Istenem, * új és erős lelket önts belém! ---

        13. Ne taszíts el színed elől, * és szent lelkedet ne vond meg tőlem!\\
        14. Add meg újra az üdvösség örömét, * a készséges lelkületet erősítsd bennem! ---

        15. A vétkeseknek megmutatom utadat, * és a bűnösök hozzád térnek.\\
        16. Ments meg a vérontástól, Isten, szabadító Istenem, * igazságosságodat ujjongva hirdeti nyelvem. ---

        17. Nyisd meg, Uram, ajkamat, * hogy dicséretedet hirdesse szavam!\\
        18. Hiszen te a véres áldozatot nem kedveled, * s ha égőáldozatot hoznék, nem vennéd szívesen.\\
        19. Isten előtt kedves áldozat a megtört lélek, * te nem veted meg a töredelmes és alázatos szívet. ---

        20. Jóságodban, Uram, tégy jót Sionnal, * építsd fel újra Jeruzsálem falait!\\
        21. Akkor majd neked tetsző áldozatot hoznak eléd, * és oltárodon fiatal tulkokat áldoznak.
}
\end{psalmsheet}

\redtext{2. antifóna}

\parttitle[Hab 3, 2-4. 13a. 15-19]{Kantikum}

\printsheet{DomineAudivi}

\begin{psalmsheet}[tone=tonus-irregularis, continuous, title={Isten megjelenik ítéletre}, motto={Emeljétek föl fejeteket, mert elérkezett megváltástok ideje (Lk 21, 28).}]
    \psalmtext{
        2. Hallottam, Uram, amit hirdettél, * és félelem fogott el.\\
        2. Uram, a te művedet * keltsd életre köztünk napjainkban!\\
        2. Napjainkban is hadd lássuk azokat, * haragodban is emlékezzél irgalmadra!\\
        3. Jön az Isten Támán felől, délről, * Fárán hegyéről jön népünk Szentje. ---

        3. Dicsősége elborítja az eget, * és fölsége betölti a földet.\\
        4. Úgy tündöklik, mint a napfény, † kezéből sugárnyaláb ragyog, * nagy ereje abban rejlik. ---

        13. Harcba szálltál, hogy megszabadítsd népedet, * és megvédd azt, akit fölkentél.\\
        15. Lovaiddal utat tapostál a tengeren, * a nagy vizek hullámain.\\
        16. Hallottam ezt, és beleremegett a szívem, * hangjától reszketett az ajkam.\\
        16. Csontjaimba kínos kór nyilallt, * lépteim bizonytalanná váltak,\\
        16. de én békén várom a szorongatás napját, * amely a minket szorongató népre virrad. ---

        17. Mert nem hajt ki többé a fügefa, * a szőlőtőkén nem lesz gyümölcs.\\
        17. Az olajfa nem hoz termést, * és a szántóföld nem terem eledelt.\\
        17. A juhok az aklokból kifogynak, * és a jászlak előtt nem áll jószág. ---

        18. Én azonban örvendezem az Úrban, * és ujjongok szabadító Istenemben.\\
        19. Az Isten, az Úr az én erőm, * aki lábamat hasonlóvá teszi a szarvaséhoz.\\
        19. A Győztes a magasságokra vezet engem, * ahol énekelhetem zsoltáraimat.
}
\end{psalmsheet}

\redtext{3. antifóna}

\printsheet{Lauda}

\begin{psalmsheet}[tone=tonus-irregularis, continuous, number={147. zsoltár}, title={Jeruzsálem megújul}, motto={Gyere, megmutatom neked a menyasszonyt, a Bárány jegyesét (Jel 21, 9).}]
    \psalmtext{
        12. Dicsérd, Jeruzsálem, az Urat, * Istenedet magasztald, Sion! ---

        13. Mert erőssé tette kapuid zárát, * és fiaidat megáldja benned.\\
        14. Határaidon ő ad békét, * és jóllakat kövér búzával.\\
        15. Elküldi szavát a földre, * gyorsan fut az ő igéje. ---

        16. Gyapjúként teríti szét a havat, * és mint a hamut, úgy szórja a zúzmarát.\\
        17. Jegét hullatja, mint a morzsát; * ki bírná elviselni annak hidegét?\\
        18. De kibocsátja igéjét, és szétolvasztja mindezt, * lehelete ráfúj, és a vizek árja zúg. ---

        19. Igéjét Jákobnak hirdeti, * törvényét és végzéseit Izrael népének.\\
        20. Egyetlen néppel sem bánt így, * végzéseit nem fedte fel senki másnak!
}
\end{psalmsheet}

\parttitle[Ef 2, 13-16]{Rövid olvasmány}
Most ti, akik „távol” voltatok, Krisztus Jézusban „közel” kerültetek, Krisztus vére ár\underline{á}n. † Ő, a mi békességünk a kettőt eggyé forrasztotta, és a közéjük emelt válaszfalat ledöntötte, az ellenségeskedést kiküszöbölte, a törvényt ugyanis parancsaival és rendelkezéseivel saját testében érvénytel\underline{e}nítette. * Mint békeszerző a két népet magában eggyé, új emberré teremtette, és egy testben mind a kettőt kiengesztelte az Istennel kereszthalála által, amellyel az ellenségeskedést megsz\underline{ü}nt\underline{e}tte.

\parttitle{Rövid válaszos ének}

\printsheet{ClamaboAdDominumAltissimum}

\parttitle[Lk 1,68-79]{Evangéliumi kantikum}

\printsheet{PerViscera}

\parttitle{Fohászok}

\begin{preces}{Krisztus Urunk a Szentlélek által, vére ontásával feláldozta magát az Atyának, hogy megtisztítsa lelkiismeretünket a holt cselekedetektől. Imádjuk őt, és valljuk meg őszinte szívvel:}{A te akaratodban van a békességünk, Urunk.}
  Jóságodnak köszönjük e mai nap kezdetét,

  \reddash segíts, hogy új életet kezdjünk!

  Te alkottál mindent, és gondot viselsz mindenre,

  \reddash add, hogy a teremtett dolgokban mindenkor meglássuk a te kezed alkotásait!

  Véreddel szentesítetted az új és örök szövetséget,

  \reddash add, hogy parancsaidat teljesítve hűek maradjunk hozzá!

  Amikor a kereszten függtél, vér és víz ömlött oldalad sebéből,

  \reddash vérednek üdvösséges patakjával mosd le bűneinket és örvendeztesd meg népedet!
\end{preces}
Mi Atyánk...

\parttitle{Könyörgés}

Add meg, kérünk, mindenható Istenünk, hogy a tiszteletedre most elvégzett dicséretet egykor mindörökre tovább énekelhessük szentjeiddel. A mi Urunk.

\subsubsection*{Befejező imaóra}

\parttitle{Zsoltározás}

\redtext{Antifóna}

\printsheet{DomineDeusSalutisMeae}

\begin{psalmsheet}[tone=4, continuous, number={87. zsoltár }, title={Súlyos beteg imája}, motto={Ez a ti órátok, a sötétség hatalmáé (Lk 22, 53).}]
  \psalmtext{
    2. Uram, szabadító Istenem, * éjjel-nappal tehozzád kiáltok. \\
    3. Imádságom jusson színed elé, * hajlítsd füledet kérésemre! ---

    4. Mert a lelkem tele gyötrelemmel, * életem közel van a holtak országához. \\
    5. A sírba szállók közé számítanak, * olyan ember vagyok, aki a végét járja. \\
    6. A halottak között van a fekhelyem, * mint aki megsebesítve a sírban alszik, \\
    7. akinek már emlékét sem őrzik többé, * és aki már kiszakadt a kezedből. ---

    8. A verem mélyére vetettél, * a sötétségbe és a halál árnyékába. \\
    9. Rám nehezedett súlyos haragod, * örvényeid fölöttem összecsapnak. ---

    10. Ismerőseimet távol tartod tőlem, * utálatossá tettél előttük. \\
    11. Fogoly vagyok, szabadulásom nincsen, * szemem elsorvad a kíntól. \\
    12. Uram, egész nap hozzád kiáltok, * és kezemet feléd tárom. ---

    13. Csodákat talán a holtaknak teszel, * vagy az árnyak fölkelnek, hogy áldjanak? \\
    14. Vajon irgalmadat hirdeti valaki a sírban, * vagy hűségedet a pusztulás helyén? \\
    15. Vajon csodáidat meglátja valaki a sötétben, * vagy igazságodat a feledés földjén? ---

    16. De én hozzád kiáltok, Uram, * már kora reggel eléd siet imádságom. \\
    17. Miért taszítasz el magadtól, Uram, * arcodat miért rejted el tőlem? \\
    18. Nyomorult vagyok, és halálra vált ifjúkorom óta, * megzavart és földre nyomott csapásaid súlya. \\
    19. Rám zúdult haragod, * és feldúlt a tőled való rettegés. \\
    20. Hullámzik körülöttem egész nap, mint az árvíz, * már-már összecsap fölöttem. \\
    21. Elszakítottad tőlem jó barátaimat, * már csak a sötétség ismer engem.
}
\end{psalmsheet}

\parttitle[Jer 14, 9]{Rövid olvasmány}
Itt élsz közöttünk, Ur\underline{a}m, † a te nev\underline{e}det viseljük. * Ne hagyj el hát minket, Urunk, \underline{I}st\underline{e}nünk!

\parttitle{Könyörgés}
Mindenható Istenünk, add, hogy sírjában nyugvó egyszülött Fiad mellett hűségesen kitartsunk, és méltók legyünk vele együtt új életre támadni. Aki él és uralkodik mindörökkön-örökké.

\subsection*{Szent Benedek apát, Európa fővédőszentje - 3. nap}

\begin{center}\redtext{Ünnep}\end{center}

\martyrology{Az umbriai Nursiában (ma Norcia) született 480 körül. Rómában nevelkedett, Subiaco vidékén remeteéletet kezdett, ahol összegyűjtötte tanítványait, majd Monte Cassinóra ment. Itt híres kolostort alapított, és Regulát állított össze, amely mindenfelé elterjedt. Ezért a nyugati szerzetesség atyjának is szokták őt nevezni. 547. március 21-én halt meg, de már a VIII. század végétől sokfelé július 11-én emlékeztek meg róla. VI. Pál pápa 1964. október 24-én Európa fővédőszentjévé nyilvánította.}

\subsubsection*{Imádságra hívás}

\printsheet{MirabilemInSanctisSuisDominum}

\subsubsection*{Olvasmányos imaóra}

\parttitle{Himnusz}

\printsheet{InterAeternas}


Testben ifjú még, de szívedben érett:\\
perzselő vágyak tüze nem hevített,\\
földi életnek gyönyörét kerülve\\
égre figyeltél.\\

Távol várostól, a szülői háztól,\\
sziklás erdőknek vadonába tértél,\\
s boldog életnek tükörét megírtad:\\
bölcs Reguládat.\\

Így a népek nagy nevelője lettél,\\
hintve Jézus szép tanait; te gyúrj át\\
minket is bátor követőddé, s gyújts fel\\
mennyei vágyra.\\

Zengje énekszó az Atyát Fiával,\\
és a Szentlelket, vigaszát ki adja:\\
Ők díszítettek kegyelemmel téged\\
gazdagon áldva. Ámen.

\parttitle{Zsoltározás}

\redtext{1. antifóna}

\printsheet{Vitam}

\begin{psalmsheet}[tone=8, continuous, number={20. zsoltár, 2-8. 14}, title={Hálaadás a király győzelméért}, motto={Amikor Krisztus feltámadt, életet kapott és a napok hosszú sorát mindörökre (Szt. Iréneusz).}]
    \psalmtext{
        2. Urunk, hatalmadnak örvend a király, * és ujjongva ujjong segítséged miatt.\\
        3. Szíve vágyát betöltötted, * nem tagadtad meg ajka óhaját.\\
        4. Kitüntetted boldogító áldásoddal, * színarany koronát tettél a fejére.\\
        5. Életet kért tőled, és te megadtad neki * a napok hosszú sorát mindörökre. ---

        6. Oltalmad alatt nagy az ő dicsősége, * méltóságot és díszt árasztasz reá.\\
        7. Áldássá teszed őt mindörökre, * örömmel töltöd el színed előtt.\\
        8. Mivel a király az Úrra hagyatkozik, * a Fölséges irgalma nem hagyja inogni. ---

        14. Kelj föl, Urunk, hatalmaddal, * mi pedig éneklünk, zsoltárt zengünk nagy tetteidről.
}
\end{psalmsheet}

\redtext{2. antifóna}

\printsheet{IustumDeduxit}

\begin{psalmsheet}[tone=7, continuous, number={91. zsoltár}, title={A teremtő Isten dicsérete}, motto={Dicsőítjük az egyszülött Fiút tettei miatt (Szt. Atanáz).}, numeral={I.}]
    \psalmtext{
        2. Jó dolog az Urat áldani, * énekkel magasztalni nevedet, Fölséges;\\
        3. már hajnalban hirdetni jóságodat, * és éjszaka zengeni hűségedet,\\
        4. tízhúrú lanton és hárfazenével, * citera mellett énekelve.\\
        5. Mert örömmel töltenek el alkotásaid, * kezed művei láttán örvendezem. ---

        6. Mily nagyszerűek a te műveid, Uram, * gondolataid mily mélységesek!\\
        7. Az oktalan ember föl nem fogja, * és nem érti meg őket az esztelen.\\
        8. Sarjadhatnak a bűnösök, akár a fű, * és virágozhatnak a gonosztevők, ---

        9. sorsuk mégis örökös pusztulás; * te azonban, Uram, fölséges vagy örökké.
}
\end{psalmsheet}

\redtext{3. antifóna}

\printsheet{IustusUtPalma}

\begin{psalmsheet}[tone=8, continuous, numeral={II.}]
    \psalmtext{
        10. Mert íme, ellenségeid, Uram, † mert íme, ellenségeid elvesznek, * és elpusztulnak mind a gonosztevők.\\
        11. Erőt adsz nekem, bikáéhoz hasonlót, * és fölkentél színtiszta olajjal.\\
        12. Szemem ellenségeimre lenéz, * s a rám törő gonoszok vesztéről hall a fülem. ---

        13. Virul az igaz, akár a pálma, * úgy növekszik, mint a Libanon cédrusa.\\
        14. Az Úr házában van elültetve, * Istenének csarnokában bont virágot.\\
        15. Gyümölcsöt érlel öreg korában is, * termékeny marad, szépen zöldellő.\\
        16. Hirdeti, hogy igaz az Isten, * ő a menedékem, nincs benne hamisság.
}
\end{psalmsheet}

\versicle{Egyenes utakon vezette az Úr az igaz\underline{a}t.}{És megmutatta neki Isten ország\underline{á}t.}

\parttitle[Lk 12, 35-36; Mt 24, 42]{Válaszos ének}

\versicle{%
    Csípőtök legyen f\underline{e}löv\underline{e}zve, † és kezetekben égjen a lámp\underline{á}s\underline{o}tok. * Hasonlítsatok azokhoz az emberekhez, akik Urukra várnak, hogy megérkezzék a m\underline{e}nyegz\underline{ő}ről.
}{%
    Legyetek hát éb\underline{e}rek, † mert nem tudj\underline{á}tok, * melyik órában jön el \underline{U}r\underline{a}tok.
}

\parttitle{Válaszos ének}

\versicle{%
    Szent Benedek elhagyta szülei ház\underline{á}t, vagy\underline{o}nát, † és egyedül Istennek akarván tetszeni, a szent szerzetesi életet vál\underline{a}szt\underline{o}tta. * A magányban élt, Isten sz\underline{í}ne \underline{e}lőtt.
}{%
    Visszavonult a vil\underline{á}gtól, † hogy ne tudja, amit már meg\underline{i}smert, * és megtapasztalja, amit még n\underline{e}m t\underline{u}dott.
}

\rubric{Te Deum, \pageref{TeDeum}.}

\parttitle{Könyörgés}
\rubric{Mint a Reggeli dicséretben.}

\subsubsection*{Reggeli dicséret}

\parttitle{Himnusz}

\printsheet{LegiferPrudens}

Népeket tettél egyazon családdá,\\
felvirágzott így a tökéletesség,\\
és szelíden lett követőid éke\\
Krisztus igája.\\

Szent szabályoddal szabadok s rabszolgák\\
Jézusunknak hű követői lettek;\\
s megvalósult bölcs Regulád világa:\\
munka, imádság.\\

Így vezéreld szent akaratra néped,\\
egy közösségben egyetértve éljünk,\\
és erősödjék az öröm közöttünk,\\
hozva a békét.\\

Zengje énekszó az Atyát Fiával,\\
és a Szentlelket, vigaszát ki adja:\\
ők díszítettek kegyelemmel téged\\
gazdagon áldva. Ámen.

\parttitle{Zsoltározás}

\redtext{1. antifóna}

\printsheet{Iucunditatem}

\begin{psalmsheet}[tone=1, continuous, number={62. zsoltár, 2-9}, title={Az Istenre szomjas lélek}, motto={Az keresi virrasztva Istent, aki a sötétség cselekedeteit elveti magától.}]
    \psalmtext{
        2. Isten, te vagy az én Istenem, * virrasztva kereslek.\\
        2. Terád szomjas a lelkem, testem utánad eped, * mint a puszta kiaszott földje.\\
        3. Szentélyedben hadd jelenjek meg előtted, * hogy lássam hatalmad és dicsőséged.\\
        4. Mert irgalmad többet ér, mint az élet, * hadd magasztaljon ezért ajkam! ---

        5. Neked mondok áldást, amíg csak élek, * a te nevedben tárom imára kezemet.\\
        6. Mint dús lakomával, teljék be lelkem, * ajkam ujjongjon, szám dicsőítsen téged.\\
        7. Még fekvőhelyemen is rád gondolok, * hajnalig rólad elmélkedem,\\
        8. mivel védelmezőm lettél, * és szárnyaid oltalmában örvendezem. ---

        9. Lelkem szorosan átölel téged, * és a jobbod szilárdan tart engem.
}
\end{psalmsheet}

\redtext{2. antifóna}

\parttitle[Dán 3, 57-88. 56]{Kantikum}

\printsheet{ServiDomini}

\begin{psalmsheet}[tone=tonus-peregrinus, continuous, title={Minden teremtmény dicsérje az Urat!}, motto={Dicsérjétek Istenünket, ti, összes szolgái! (Jel 19, 5)}]
    \psalmtext{
        57. Áldja az Urat az Úr minden műve, * dicsérje és magasztalja örökkön örökké!\\
        58. Áldjátok az Urat, az Úr angyalai mindnyájan; * áldjátok az Urat, ti, egek!\\
        60. Áldjátok az Urat, ti, vizek az égbolt fölött; * áldjátok az Urat, az Úr minden erői!\\
        62. Áldjátok az Urat, nap és holdvilág; * áldjátok az Urat, égi csillagok! ---

        64. Áldjátok az Urat, záporeső és harmat; * áldjátok az Urat, ti, viharok!\\
        66. Áldjátok az Urat, tűz és forróság; * áldjátok az Urat, hideg és hőség!\\
        68. Áldjátok az Urat, dér és harmat; * áldjátok az Urat, fagy és nagy hideg!\\
        70. Áldjátok az Urat, jég és hóesés; * áldjátok az Urat, éjek és nappalok!\\
        72. Áldjátok az Urat, világosság és sötétség; * áldjátok az Urat, villámok és fellegek! ---

        74. Áldja az Urat a szárazföld; * dicsérje és magasztalja örökké!\\
        75. Áldjátok az Urat, hegyek és halmok; * áldjátok az Urat mind, amik a földből sarjadtok!\\
        77. Áldjátok az Urat, ti, források; * áldjátok az Urat, folyamok és tengerek!\\
        79. Áldjátok az Urat, tengeri lények, s minden, ami mozog a vízben; * áldjátok az Urat, égi madarak!\\
        81. Áldjátok az Urat, vadak és háziállatok; * áldjátok az Urat, ti, emberek fiai! ---

        83. Áldja az Urat Izrael, * dicsérje és magasztalja örökké!\\
        84. Áldjátok az Urat, ti, az Úr papjai; * áldjátok az Urat, ti, az Úr szolgái!\\
        86. Áldjátok az Urat, igaz szellemek és lelkek; * áldjátok az Urat, szentek és alázatos szívűek!\\
        88. Áldjátok az Urat, Hananja, Azarja és Misael, * dicsérjétek és magasztaljátok örökké! ---

        88. Áldjuk az Atyát és a Fiút a Szentlélekkel együtt; * dicsérjük és magasztaljuk örökké!\\
        56. Áldott vagy, Urunk, az égbolt fölött, * dicséretre méltó, dicsőséges és magasztos örökké!
}
\end{psalmsheet}

\rubric{Ennek a kantikumnak a végén nem mondjuk: \textcolor{black}{Dicsőség az Atyának}.}

\redtext{3. antifóna}

\printsheet{Exsultabunt}

\begin{psalmsheet}[tone=4, continuous, number={149. zsoltár}, title={A szentek ujjongása
}, motto={Az Egyház fiai, Isten új népének fiai ujjongjanak királyukban: Krisztusban (Hesychius).}]
    \psalmtext{
        1. Új dalt énekeljetek az Úrnak, * dicsérete zengjen hívei gyülekezetében.\\
        2. Teremtőjének örvendjen Izrael, * és Sion fiai királyuk miatt ujjongjanak.\\
        3. Dicsérjék nevét körtánccal, * kézidobbal s citerával zsoltárt zengjenek neki.\\
        4. Mert népében kedvét találja az Úr, * és a megalázottakat diadallal koszorúzza.\\
        5. Örvendjenek dicsőségüknek a szentek, * ujjongjanak nyugvóhelyükön. ---

        6. Szájuk Istent dicsérje, * és kétélű kard legyen a kezükben,\\
        7. hogy a pogányokon bosszút álljanak, * és a népeket megfenyítsék.\\
        8. Királyaikat verjék béklyóba, * és uraikat vasbilincsbe,\\
        9. hogy végrehajtsák rajtuk az ítéletet, amint írva van: * Ez lesz dicsősége összes szentjeinek!
}
\end{psalmsheet}

\parttitle[Róm 12, 1-2]{Rövid olvasmány}
Testvérek, Isten irgalmára kérlek benneteket: Adjátok testeteket élő, szent, Istennek tetsző áldozat\underline{u}l. † Ez legyen szellemetek h\underline{ó}dolata. * Ne hasonuljatok a világhoz, hanem gondolkodástokban megújulva alakuljatok át, hogy felismerjétek, mi az Isten akarata, mi a helyes, mi a kedves előtte, és mi a tök\underline{é}l\underline{e}tes.

\parttitle{Rövid válaszos ének}

\printsheet{LexDeiEius}

\parttitle[Lk 1,68-79]{Evangéliumi kantikum}

\printsheet{FuitVir}

\parttitle{Fohászok}

\begin{preces}{Testvéreim, magasztaljuk Krisztust, szent Istenünket, és kérjük tőle, hogy szentségben járjunk színe előtt életünk minden napján. Mondjuk hangos szóval:}{Urunk, egyedül te vagy a Szent!}
  Te hasonló lettél hozzánk és vállaltad minden kísértésünket, de bűnbe nem estél,

  \reddash irgalmazz nekünk, Urunk, Jézus!

  Te mindnyájunkat meghívtál a tökéletes szeretetre,

  \reddash szentelj meg minket, Urunk, Jézus!

  Te úgy akartad, hogy a föld sója és a világ világossága legyünk,

  \reddash világosíts meg minket, Urunk, Jézus!

  Te szolgálni jöttél, és nem azért, hogy neked szolgáljanak,

  \reddash segíts, hogy alázatosan szolgáljunk testvéreinknek és neked, Urunk, Jézus!

  Te az Atya dicsőségének fénye és lényének képmása vagy,

  \reddash segíts, hogy meglássuk dicsőséges arcodat, Urunk, Jézus!
\end{preces}
Mi Atyánk...

\parttitle{Könyörgés}
Istenünk, te Szent Benedek apátot az istenszolgálat iskolájában kiváló mesterré tetted. Add, kérünk, hogy az irántad való szeretetet mindennél többre becsüljük, és szárnyaló szívvel járjunk parancsaid útján. A mi Urunk.

\subsection*{Évközi 15. vasárnap - 4. nap}

\rubric{I. Esti dicséret utáni Befejező imaóra, mint az előtábor 2. napjának I. Esti dicsérete után, \pageref{D/Compl1}.}

\rubric{Imádságra hívás, mint az előtábor 2. napján, \pageref{D/Inv}.}

\subsubsection*{Olvasmányos imaóra}

\parttitle{Himnusz}

\printsheet{PrimoDierumOmnium}

\begin{hymnstanzas}
    fürgébben serkenjünk ma fel,\\
    a kábult álmot űzzük el,\\
    Istent keressük éjszaka,\\
    ez prófétáknak szózata,
\stanzabreak
    hogy hallja meg kérő imánk,\\
    jobbját terjessze védve ránk,\\
    hogy szennyeinktől szabadon\\
    hívjon magához egykoron,
\stanzabreak
    kik a nyugvás óráiban,\\
    a napnak legszentebb szakán,\\
    zsolozsmaszóval áldjuk őt:\\
    boldog kegyelmét nyerjük el.
\stanzabreak
    Dicsérjük az örök Atyát,\\
    dicsérjük egyszülött Fiát\\
    s a Lelket, a Vigasztalót,\\
    örökre egy uralkodót. Ámen.
\end{hymnstanzas}

\parttitle{Zsoltározás}

\redtext{1. antifóna}

\printsheet{EcceDeusMeus}

\begin{psalmsheet}[tone=4, continuous, number={144. zsoltár}, title={Az isteni fölség dicsérete}, motto={Igazságos vagy, Uram, aki vagy, és aki voltál! (Jel 16, 5)}, numeral={I.}]
    \psalmtext{
        1. Magasztallak, Istenem, Királyom, † és áldom nevedet * mindörökkön-örökké.\\
        2. Naponta áldalak, † és nevedet magasztalom * mindörökkön-örökké. ---

        3. Nagy az Úr, és minden dicséretre méltó, * és nagyságát felfogni nem lehet.\\
        4. Nemzedékről nemzedékre dicsérjék tetteidet, * és hangosan hirdessék hatalmadat! ---

        5. Magasztos fölséged nagyságáról beszéljenek, * és beszéljék el csodatetteidet!\\
        6. Félelmetes tetteid erejéről szóljanak, * és hirdessék a te nagy voltodat!\\
        7. Terjesszék végtelen jóságod hírét, * és igazságosságodat ujjongva áldják!\\
        9. Jóságos az Úr mindenkihez, * irgalmas minden teremtményéhez.\\
        8. Irgalmas és kegyelmes az Úr, * hosszan tűrő és nagy irgalmú.
}
\end{psalmsheet}

\redtext{2. antifóna}

\printsheet{RegnumTuumDomine}

\begin{psalmsheet}[tone=1, continuous, numeral={II.}]
    \psalmtext{
        10. Magasztaljon téged, Uram, minden műved, * és szentjeid áldást mondjanak!\\
        11. Hirdessék, hogy országod milyen dicsőséges, * és hatalmadról beszéljenek,\\
        12. hogy megismertessék az emberekkel hatalmadat * és nagyszerű országod dicsőségét.\\
        13. Örökkévaló ország a te országod, * uralmad minden nemzedéken át megmarad.
}
\end{psalmsheet}

\redtext{3. antifóna}

\printsheet{InAeternum}

\begin{psalmsheet}[tone=8, continuous, numeral={III.}]
    \psalmtext{
        13. Hűséges az Úr minden szavában, * és szent minden cselekedetében.\\
        14. Az Úr fenntartja a botladozót, * és a megalázottat fölemeli.\\
        15. Mindenek szeme bízón rád tekint, * és te enni adsz nekik alkalmas időben.\\
        16. Megnyitod kezedet, * és eltöltesz minden élőt jótéteménnyel. ---

        17. Igazságos az Úr minden végzésében, * és szent minden cselekedetében.\\
        18. Közel az Úr azokhoz, akik őt szólítják, * mindazokhoz, akik őt igaz szívvel hívják.\\
        19. Az istenfélők vágyait teljesíti, † kérésüket meghallgatja, * és megmenti őket.\\
        20. Megóv az Úr mindenkit, aki őt szereti, * de elpusztít minden gonosztevőt. ---

        21. Hirdesse szám az Úr dicsőségét, † és szent nevét áldja minden élő * mindörökkön-örökké.
}
\end{psalmsheet}

\versicle{Figyelj, fiam, szavaimr\underline{a}.}{Nyisd meg füledet beszédemr\underline{e}!}

\parttitle[Jak 5, 17. 18; Sir 48, 1. 3]{Válaszos ének}

\versicle{%
    Illés esedezve imádkozott, hogy ne ess\underline{é}k az \underline{e}ső, † s nem \underline{i}s \underline{e}sett. * Majd ismét imádkozott, és esőt \underline{a}dott \underline{a}z ég.
}{%
    Mint a tűzvész, Illés prófét\underline{a} jött, † szava lángolt, mint az égő f\underline{á}klya, * az Úr szavával elzárta \underline{a}z \underline{e}get.
}

\parttitle[Tit 3, 3. 5; Ef 2, 3]{Válaszos ének}

\versicle{%
    Egykor mi magunk is balgák, engedetlenek és tév\underline{e}lygők v\underline{o}ltunk; † gonoszságunk, irigységünk miatt utálatra méltók voltunk, és gyűlölt\underline{ü}k \underline{e}gymást. * De Isten megmentett minket irgalmasságból, s a Szentlélekben való újjászületés és megújulás f\underline{ü}rdőj\underline{é}ben.
}{%
    Valamikor testi vágyainkban \underline{é}ltünk, * és születésünknél fogva a harag gyermeke\underline{i} v\underline{o}ltunk.
}

\rubric{Te Deum, \pageref{TeDeum}.}

\parttitle{Könyörgés}
\rubric{Mint a Reggeli dicséretben.}

\subsubsection*{Reggeli dicséret}

\parttitle{Himnusz}

\printsheet{AeterneRerumConditor}

\begin{hymnstanzas}
    szól már a napnak hírnöke,\\
    ki éber őr mély éjen át,\\
    s az utazónak éji fény,\\
    mert minden órát megkiált.
\stanzabreak
    A hajnalcsillag kél szaván,\\
    égről borút elűz tova;\\
    elhagyja ártó útjait\\
    az éji kóborlók hada.
\stanzabreak
    Hangján hajós erőre kap,\\
    a tengerárra csend borul;\\
    s az Egyház sziklaszál feje\\
    bűnét siratni elvonul.
\stanzabreak
    Úr Jézus! A bukóra nézz:\\
    térítsen meg tekinteted;\\
    ha ránk tekintesz, hull a bűn,\\
    s elmossa könny a vétkeket.
\stanzabreak
    Szent Fényesség, ragyogj belénk,\\
    a lélek álmát megszakaszd;\\
    első szavunk tiéd legyen,\\
    s amit fogadtunk, álljuk azt!
\stanzabreak
    Krisztus, kegyelmes nagy Király,\\
    neked s Atyádnak tisztelet,\\
    s Vigasztalónknak is veled\\
    most és örökre szüntelen. Ámen.
\end{hymnstanzas}

\parttitle{Zsoltározás}

\redtext{1. antifóna}

\printsheet{RegnavitDominusDecorem}

\begin{psalmsheet}[tone=4, continuous, number={92. zsoltár}, title={A teremtő Isten fölsége}, motto={Átvette uralmát Urunk, mindenható Istenünk: örvendezzünk, ujjongjunk, és dicsőítsük őt! (Jel 19, 6. 7)}]
  \psalmtext{
    1. Király az Úr, fönségbe öltözött; * felöltözött az Úr, felövezte magát erősséggel. \\
    2. Szilárddá tette a földet, hogy meg ne inogjon. † Trónod szilárdan áll ősidők óta, * öröktől fogva vagy te. ---

    3. Uram, a vizek árja zúg, † harsogva zúgnak a folyamok, * mennydörögve zúg a tornyosuló vizek árja. \\
    4. De hatalmasabb, mint a tenger morajlása, † hatalmasabb, mint a nagy vizek zúgása, * mindennél hatalmasabb az Isten a magasságban. ---

    5. Bizonyságaid valóban igazak, * házadat szentség illeti, Uram, mindörökké.
}
\end{psalmsheet}

\redtext{2. antifóna}

\parttitle[Dán 3, 57-88. 56]{Kantikum}

\printsheet{Benedicite}

\begin{psalmsheet}[tone=7, continuous, title={Minden teremtmény dicsérje az Urat!}, motto={Dicsérjétek Istenünket, ti, összes szolgái! (Jel 19, 5)}]
    \psalmtext{
        57. Áldja az Urat az Úr minden műve, * dicsérje és magasztalja örökkön örökké!\\
        58. Áldjátok az Urat, az Úr angyalai mindnyájan; * áldjátok az Urat, ti, egek!\\
        60. Áldjátok az Urat, ti, vizek az égbolt fölött; * áldjátok az Urat, az Úr minden erői!\\
        62. Áldjátok az Urat, nap és holdvilág; * áldjátok az Urat, égi csillagok! ---

        64. Áldjátok az Urat, záporeső és harmat; * áldjátok az Urat, ti, viharok!\\
        66. Áldjátok az Urat, tűz és forróság; * áldjátok az Urat, hideg és hőség!\\
        68. Áldjátok az Urat, dér és harmat; * áldjátok az Urat, fagy és nagy hideg!\\
        70. Áldjátok az Urat, jég és hóesés; * áldjátok az Urat, éjek és nappalok!\\
        72. Áldjátok az Urat, világosság és sötétség; * áldjátok az Urat, villámok és fellegek! ---

        74. Áldja az Urat a szárazföld; * dicsérje és magasztalja örökké!\\
        75. Áldjátok az Urat, hegyek és halmok; * áldjátok az Urat mind, amik a földből sarjadtok!\\
        77. Áldjátok az Urat, ti, források; * áldjátok az Urat, folyamok és tengerek!\\
        79. Áldjátok az Urat, tengeri lények, s minden, ami mozog a vízben; * áldjátok az Urat, égi madarak!\\
        81. Áldjátok az Urat, vadak és háziállatok; * áldjátok az Urat, ti, emberek fiai! ---

        83. Áldja az Urat Izrael, * dicsérje és magasztalja örökké!\\
        84. Áldjátok az Urat, ti, az Úr papjai; * áldjátok az Urat, ti, az Úr szolgái!\\
        86. Áldjátok az Urat, igaz szellemek és lelkek; * áldjátok az Urat, szentek és alázatos szívűek!\\
        88. Áldjátok az Urat, Hananja, Azarja és Misael, * dicsérjétek és magasztaljátok örökké! ---

        88. Áldjuk az Atyát és a Fiút a Szentlélekkel együtt; * dicsérjük és magasztaljuk örökké!\\
        56. Áldott vagy, Urunk, az égbolt fölött, * dicséretre méltó, dicsőséges és magasztos örökké!
}
\end{psalmsheet}

\rubric{Ennek a kantikumnak a végén nem mondjuk: \textcolor{black}{Dicsőség az Atyának}.}

\redtext{3. antifóna}

\printsheet{Laudate}

\begin{psalmsheet}[tone=1, continuous, number={148. zsoltár}, title={A teremtő Isten magasztalása}, motto={A trónon ülőnek és a Báránynak áldás, tisztelet, dicsőség és hatalom örökkön-örökké! (Jel 5, 13)}]
  \psalmtext{
    1. Dicsérjétek az Urat a mennyből, * dicsérjétek a magasságban! \\
    2. Dicsérjétek, ti, angyalai, mind, * dicsérjétek, összes mennyei karok! ---

    3. Dicsérje őt a nap és a hold is, * dicsérje őt minden fényes csillag! \\
    4. Dicsérjék őt az egek egei, * és minden vizek az egek fölött! ---

    5. Dicsérjék mind az Úr nevét, * mert ő parancsolt, és minden létrejött. \\
    6. Mindörökre helyükre állította őket, * törvényt szabott, amely el nem évül. ---

    7. Dicsérjétek az Urat a földről, * tengeri szörnyek és összes mély vizek, \\
    8. tűz és jégeső, hó és ködfelhő, * és igéjét teljesítő száguldó vihar, \\
    9. ti, hegyek és halmok, mindnyájan, * gyümölcsfák és összes cédrusok, \\
    10. erdei vadak és minden háziállat, * csúszómászók és szárnyas madarak! ---

    11. Föld királyai és minden népek, * fejedelmek és a föld minden bírája, \\
    12. ifjak és leányok, * fiatalok és öregek együtt, \\
    13. dicsérjék az Úr nevét, * mert egyedül csak az ő neve fölséges! ---

    14. Fölsége felülmúlja az eget és a földet, * és ő emeli fel népe homlokát. \\
    15. Ő teszi dicsővé minden szentjét, * Izraelt, a hozzá közel álló népet.
}
\end{psalmsheet}

\parttitle[Ez 37, 12b-14]{Rövid olvasmány}
Ezt mondja az Úr, az Isten: Nézzétek, kinyitom sírjaitokat, és kihozlak benneteket sírjaitokból, én népem, és elvezetlek benneteket Izrael földjér\underline{e}. † Akkor majd megtudjátok, hogy én vagyok az Úr, amikor kinyitom sírjaitokat, és kihozlak benneteket sírjaitokb\underline{ó}l, én népem. * Belétek oltom lelkemet, és életre keltek. Letelepítelek benneteket földeteken, és megtudjátok, hogy én, az Úr mondtam ezt, és végbevittem – mondj\underline{a} \underline{a}z Úr.

\parttitle{Rövid válaszos ének}

\printsheet{ChristiFiliDeiVivi}

\parttitle[Lk 1,68-79]{Evangéliumi kantikum}

\printsheet{CumTurbaPlurima}

\parttitle{Fohászok}

\begin{preces}{Isten elküldötte a Szentlelket, hogy minden szív boldogító világossága legyen. Kérjük állhatatosan:}{Világosítsd meg népedet, Urunk!}
  Istenünk, igaz fényességünk, áldunk téged,

  \reddash mert úgy rendelted, hogy dicsőségedre megérjük a mai napot.

  Fiad feltámadása által fényt derítettél a világra,

  \reddash áraszd ezt a világosságot Egyházad közreműködésével minden emberre!

  Fiad tanítványaiban kinyilvánítottad dicsőségedet Szentlelked által,

  \reddash küldd el őt, hogy Egyházad hűséges maradjon hozzád!

  Népek megvilágosítója, emlékezz meg azokról, akik még a sötétben élnek,

  \reddash nyisd meg lelki szemüket, hogy felismerjenek téged, egyedül igaz Istent!
\end{preces}
Mi Atyánk...

\label{DXV/L/Coll}

\parttitle{Könyörgés}
Istenünk, te megmutatod igazságod fényét a tévelygőknek, hogy visszataláljanak a helyes útra. Add meg híveidnek, hogy elutasítsunk mindent, ami a keresztény névvel ellenkezik, és vállaljuk, ami méltó hozzá. A mi Urunk.

\rubric{II. Esti dicséret utáni Befejező imaóra, mint az előtábor 2. napjának II. Esti dicsérete után, \pageref{D/Compl2}.}

\subsection*{Hétfő, 15. hét - 5. nap}

\rubric{Imádságra hívás, mint az előtábor 3. napján, \pageref{f2/Inv}.}

\subsubsection*{Olvasmányos imaóra}

\parttitle{Himnusz}

\printsheet{SomnoRefectisArtubus}

\begin{hymnstanzas}
    Szóljon hozzád első szavunk,\\
    hozzád a szív első heve,\\
    hogy minden tettnek ezután,\\
    Szent Úr, te légy a kezdete.
\stanzabreak
    A fény az éjt szüntesse már,\\
    sötét homályt a napsugár.\\
    Az éjben búvó bűnöket\\
    a napfény semmisítse meg.
\stanzabreak
    Imánk tehozzád esdekel:\\
    a rossz hatalmát űzzed el,\\
    és zengjen ajkunk teneked\\
    meg nem szűnő dicséretet.
\stanzabreak
    Add meg, kegyelmes jó Atya,\\
    Atyának egyszülött Fia\\
    és Szentlélek, Vigasztalónk:\\
    egyetlen Úr, örök Király. Ámen.
\end{hymnstanzas}

\parttitle{Zsoltározás}

\redtext{1. antifóna}

\printsheet{DeusDeorum}

\begin{psalmsheet}[tone=8, continuous, number={49. zsoltár}, title={A helyes vallásosság}, motto={Nem azért jöttem, hogy megszüntessem a törvényt, hanem hogy teljessé tegyem (vö. Mt 5, 17).}, numeral={I.}]
    \psalmtext{
        1. Megszólal az Úr, az istenek Istene, * napkelettől napnyugatig szólítja a földet.\\
        2. A szépséges Sionból ragyogva jelenik meg Isten, * jön a mi Istenünk, és nem hallgat.\\
        3. Emésztő tűz lobog színe előtt, * körülötte szélvihar tombol. ---

        4. Szólítja onnan felülről az eget és a földet, * hogy ítéletet tartson népe fölött.\\
        5. „Gyűjtsétek elém választottaimat, * akik velem áldozatban szövetségre léptek!”\\
        6. Igazságosságát hirdetik az egek, * Isten maga tart ítéletet.
}
\end{psalmsheet}

\redtext{2. antifóna}

\printsheet{ImmolaDeo}

\begin{psalmsheet}[tone=tonus-irregularis, continuous, numeral={II.}]
    \psalmtext{
        7. „Halljad, én népem, hadd szóljak, † Izrael, hadd tanúskodjam ellened: * Isten, a te Istened én vagyok.\\
        8. Nem áldozataid miatt dorgállak, * égőáldozataidat látom szüntelen.\\
        9. Nem kell nekem házadból a növendékbika, * sem nyájaidból a kecskebak. ---

        10. Úgyis enyém az erdő és minden vadja * és a hegyek ezernyi állata.\\
        11. Ismerem én az ég minden madarát, * és enyém, ami csak mozdul a réten.\\
        12. Ha éhezném, nem szólnék neked, * mert enyém a földkerekség és ami azt betölti. ---

        13. Megeszem talán a bikák húsát? * Vagy a bakok vérét megiszom?\\
        14. Ajánld föl Istennek a dicséret áldozatát, * és teljesítsd fogadalmadat a Fölséges előtt.\\
        15. Hívj segítségül, amikor bajba kerülsz, * megmentelek akkor, és te áldani fogsz.”
}
\end{psalmsheet}

\redtext{3. antifóna}

\printsheet{DixiIniquis}

\begin{psalmsheet}[tone=7, continuous, numeral={III.}]
    \psalmtext{
        16. A gonoszhoz pedig így beszél Isten: † „Törvényeimet mit emlegeted, * és miért veszed ajkadra szövetségemet?\\
        17. Hiszen a fegyelmet gyűlölöd, * és megveted szavaimat. ---

        18. Hogyha tolvajt látsz, társául szegődöl, * és szövetkezel a házasságtörővel.\\
        19. Szádból gonosz beszéd árad, * nyelved csalárdságot sző. ---

        20. Leülsz, és máris gyalázod testvéred, * és mocskolod anyád fiát.\\
        21. Te így tettél, és én hallgassak? * Azt véled, hogy hozzád hasonló vagyok?\\
        21. Pörbe szólítalak most, * és mindezt számon kérem tőled. ---

        22. Ti, akik Istenről megfeledkeztek, szálljatok magatokba; * különben elvesztek, s nincs, aki megmenthet titeket.\\
        23. Aki a dicséret áldozatát áldozza, az tisztel engem; † és akinek útja szeplőtelen, annak megadom, * hogy meglássa Isten üdvösségét.”
}
\end{psalmsheet}

\versicle{Halljad, én népem, hadd szólj\underline{a}k,}{Isten, a te Istened én vagy\underline{o}k.}

\parttitle[1 Kir 18, 21; Mt 6, 24]{Válaszos ének}

\versicle{%
    Illés odalépett az egész nép színe elé és m\underline{e}gkérd\underline{e}zte: † Meddig akartok még kétfelé sánt\underline{i}k\underline{á}lni? * Ha az Úr az Isten, akkor őt k\underline{ö}vess\underline{é}tek!
}{%
    Senki sem szolgálhat két Úrn\underline{a}k; † nem szolgálhattok az Istenn\underline{e}k is, * a Mamm\underline{o}nn\underline{a}k is.
}

\parttitle[Iz 44, 3. 4; Jn 4, 14]{Válaszos ének}

\versicle{%
    Elárasztom vízzel a t\underline{i}kkadt m\underline{e}zőt, † és bővizű patakokkal a kiasz\underline{o}tt f\underline{ö}ldet. * Kiárasztom Lelkemet, úgy sarjadzanak majd, mint a fűzfák a vízfoly\underline{á}sok m\underline{e}llett.
}{%
    Örök életre szök\underline{e}llő * vízf\underline{o}rr\underline{á}s lesz.
}

\rubric{Könyörgés, mint vasárnap, a Reggeli dicséretben, \pageref{DXV/L/Coll}.}

\end{document}
